% Leçon 26 — Expression écrite : caractères complexes liés à l'exode rural — Chinois 1ère A — Cours signé OnBuch+
% Programme officiel MINESEC. Compiler avec : tectonic ZHA26-ecrit-caracteres-complexes-lies-a-l-exode-rural.tex
\newcommand{\DOCMATIERE}{Chinois}
\newcommand{\DOCNIVEAU}{1ère A}
\newcommand{\DOCMODULE}{Module 3 : Citoyenneté et ouverture au monde}
\newcommand{\DOCLECON}{ZHA26}
\newcommand{\DOCTITRE}{Leçon 26 — Expression écrite : caractères complexes liés à l'exode rural}
% OnBuch+ — préambule « cours signés » ENRICHI (compilé avec tectonic / XeLaTeX)
% Système visuel « Pop » d'OnBuch+ : fond crème (jamais blanc), bordures encre
% épaisses, ombres « dures » décalées, titres Archivo Black en capitales.
% Version MATHÉMATIQUES : graphes (pgfplots/tikz), boîtes pédagogiques
% (définition, propriété, méthode, attention, le savais-tu, exercice résolu,
% exercice, TP, bilan), page de garde et signature OnBuch+ ancrée en bas.
%
% Macros attendues dans main.tex AVANT \input{preamble} :
%   \DOCMATIERE  \DOCNIVEAU  \DOCMODULE  \DOCLECON (numéro)  \DOCTITRE
\documentclass[11pt]{article}

\usepackage{fontspec}
\usepackage[french]{babel} % français = langue principale ; le chinois via \zh{..} / textebox (xeCJK)
\usepackage{xeCJK}
\usepackage{pifont} % \ding{51} \ding{55} utilisés par les modèles
\usepackage[a4paper,margin=2cm,top=3.4cm,bottom=2.8cm,headheight=1.4cm,footskip=1.3cm]{geometry}
\usepackage{amsmath,amssymb,amsfonts}
\usepackage{mathtools} % \xmapsto, \coloneqq... souvent utilisés par les modèles
\usepackage{cancel} % simplifications barrées (bilans de forces)
% \abs{z} (module, valeur absolue) et \norm{u} (norme), utilisés par les modèles
\providecommand{\abs}{}\let\abs\relax\DeclarePairedDelimiter{\abs}{\lvert}{\rvert}
\providecommand{\norm}{}\let\norm\relax\DeclarePairedDelimiter{\norm}{\lVert}{\rVert}
\usepackage[table]{xcolor} % [table] : fournit aussi \rowcolors (alternance de lignes)
\usepackage{pagecolor}
\usepackage{tikz}
\usetikzlibrary{arrows.meta,positioning,calc,shapes.geometric,shapes.misc,shapes.arrows,decorations.pathmorphing,decorations.pathreplacing,decorations.markings,patterns,shadows,mindmap,trees,fit,backgrounds,matrix,intersections,angles,quotes}
\usepackage{pgfplots}
\usepackage[version=4]{mhchem} % équations nucléaires \ce{^{235}_{92}U}
\usepackage[european,siunitx]{circuitikz} % schémas électriques
\pgfplotsset{compat=1.18}
\usepgfplotslibrary{fillbetween,groupplots} % aires entre courbes (calcul intégral)
\usepackage{enumitem}
\usepackage{fancyhdr}
% [most] charge toutes les bibliothèques tcolorbox (listings, minted, raster,
% documentation...) et gonfle inutilement la mémoire TeX sur un document long
% (~300 boîtes) : on ne charge que ce qui est réellement utilisé.
\usepackage[skins,breakable]{tcolorbox}
\usepackage{graphicx}
\usepackage{colortbl}
\usepackage{booktabs}
\usepackage{tabularx}
\usepackage{diagbox} % case d'en-tête diagonale des tableaux à double entrée
% Colonnes extensibles centrée / gauche / droite (C, L, R), souvent utilisées par les modèles
\newcolumntype{C}{>{\centering\arraybackslash}X}
\newcolumntype{L}{>{\raggedright\arraybackslash}X}
\newcolumntype{R}{>{\raggedleft\arraybackslash}X}
\usepackage{array}
\usepackage{multicol}
\usepackage{multirow}
\usepackage{siunitx}
% parse-numbers=false : accepte aussi \SI{2,5 \times 10^{-3}}{...}
\sisetup{locale=FR,output-decimal-marker={,},group-minimum-digits=4,parse-numbers=false}
\DeclareSIUnit\ohm{\ensuremath{\symup{\Omega}}} % U+2126 absent de la police
\DeclareSIUnit\division{div} % réglages d'oscilloscope (ms/div, V/div)
\DeclareSIUnit\tr{tr} % tours (vitesse de rotation, ex. \SI{1500}{\tr\per\minute})
\DeclareSIUnit\molar{M} % concentration molaire (ex. \SI{3}{\molar}, \SI{1}{\micro\molar})
\DeclareSIUnit\an{an} % année (le modèle confond parfois avec \year, un compteur TeX) — ex. \SI{5}{\kilo\meter\per\an}
\DeclareSIUnit\FCFA{FCFA} % franc CFA (ex. \SI{500}{\FCFA\per\kilo\gram})
\DeclareSIUnit\degreeBrix{\textdegree Brix} % teneur en sucre d'un sirop/jus (réfractométrie)
\DeclareSIUnit\month{mois} % le modèle utilise parfois \month, un compteur TeX, comme unité de durée
\usepackage{qrcode}
% \degree : commande fréquemment utilisée par les modèles pour un angle ou une
% température en dehors de \SI{}{} (non fournie par les paquets chargés).
\providecommand{\degree}{^{\circ}}
% \qed (fin de démonstration), utilisé par les modèles sans amsthm
\providecommand{\qed}{\hfill$\square$}
% \barre : double barre des valeurs interdites dans un tableau de variations (tkz-tab)
\providecommand{\barre}{\ensuremath{\|}}
% environnement proof (amsthm non chargé) utilisé par les modèles
\ifdefined\proof\else\newenvironment{proof}[1][Démonstration]{\par\noindent\textit{#1.}\ }{\qed\par}\fi
\usepackage{lastpage}
\usepackage{hyperref}
\hypersetup{hidelinks}

% ============================================================
% Palette « Pop » OnBuch+ — mêmes hex que la classe P (plus_ui.dart)
% ============================================================
\definecolor{poporange}{HTML}{F59321}
\definecolor{popdark}{HTML}{E07A0C}
\definecolor{popcream}{HTML}{FFF6E8}
\definecolor{popink}{HTML}{1C1714}
\definecolor{poppurple}{HTML}{7A5AE0}
\definecolor{popgreen}{HTML}{2E9E6A}
\definecolor{poppink}{HTML}{E0457B}
\definecolor{popblue}{HTML}{2F7DE1}
\definecolor{popgold}{HTML}{C98A12}
\definecolor{popmuted}{HTML}{8A7F76}
\definecolor{popline}{HTML}{EADFD2}
\definecolor{popgray}{HTML}{8A7F76}
\colorlet{popgrey}{popgray}
% Alias tolérants (le modèle nomme parfois les couleurs différemment)
\colorlet{popwhite}{white}
\colorlet{popblack}{popink}
\colorlet{popred}{poppink}
\colorlet{popyellow}{popgold}
% Teintes claires (fonds de tableaux/encadrés) — le modèle nomme parfois
% les couleurs avec un suffixe L (ex. popblueL) sans les avoir définies.
\colorlet{poporangeL}{poporange!20}
\colorlet{popdarkL}{popdark!20}
\colorlet{poppurpleL}{poppurple!20}
\colorlet{popgreenL}{popgreen!20}
\colorlet{poppinkL}{poppink!20}
\colorlet{popblueL}{popblue!20}
\colorlet{popgoldL}{popgold!20}
\colorlet{popredL}{popred!20}
\colorlet{popgrayL}{popgray!20}
% Teintes claires pour fonds de tableaux / zones
\colorlet{popblueL}{popblue!10!white}
\colorlet{poporangeL}{poporange!14!white}
\colorlet{popgreenL}{popgreen!12!white}
\colorlet{poppurpleL}{poppurple!12!white}
\colorlet{poppinkL}{poppink!12!white}
\colorlet{popgoldL}{popgold!14!white}

\pagecolor{popcream}

% ============================================================
% Typographie
% ============================================================
\defaultfontfeatures{Ligatures=TeX}
\setmainfont{PlusJakartaSans-Medium.ttf}[Path=../../../fonts/,BoldFont=PlusJakartaSans-Bold.ttf]
% Symboles, API et emojis absents de la police principale : repli sur DejaVu Sans (caractères actifs)
\newfontface\symfont{DejaVuSans.ttf}[Path=../../../fonts/]
\makeatletter
\newcommand\symactive[1]{\catcode#1=13 \begingroup\lccode`\~=#1 \lowercase{\endgroup\def~}{{\symfont\char#1}}}
\newcommand\symblank[1]{\catcode#1=13 \begingroup\lccode`\~=#1 \lowercase{\endgroup\def~}{}}
\newcommand\symhyph[1]{\catcode#1=13 \begingroup\lccode`\~=#1 \lowercase{\endgroup\def~}{\mbox{-}}}
\newcount\symc
\newcommand\symrange[2]{\symc=#1 \@whilenum\symc<#2 \do{\expandafter\symactive\expandafter{\the\symc}\advance\symc by 1 }}
\newcommand\symblankrange[2]{\symc=#1 \@whilenum\symc<#2 \do{\expandafter\symblank\expandafter{\the\symc}\advance\symc by 1 }}
\makeatother
\symrange{"0250}{"02FF}\symactive{"2713}\symactive{"2714}\symactive{"2717}\symactive{"2718}\symactive{"2610}\symactive{"2611}\symactive{"2612}
\symrange{"2600}{"27BF}
\catcode"2705=13 \begingroup\lccode`\~="2705 \lowercase{\endgroup\def~}{{\symfont\char"2713}}\symblank{"26BD}\symblank{"26A1}
\symrange{"2460}{"257F}\symactive{"223C}\symactive{"2248}\symactive{"2260}
\symactive{"01F8}\symactive{"01F9}\symactive{"1E3E}\symactive{"1E3F}
\symhyph{"2011}\symhyph{"2E1A}\symblankrange{"1F300}{"1FAFF}\symblank{"FE0F}
% Caractères chinois : WenQuanYi Zen Hei (sous-ensemble GB2312 + pinyin), gras/italique simulés
\setCJKmainfont{WenQuanYiZenHei-subset.ttf}[Path=../../../fonts/,AutoFakeBold=3,AutoFakeSlant=0.2]
\xeCJKsetup{CJKspace=false}
% Pinyin : voyelles à caron (ǎ ǐ ǒ ǔ ǖ ǘ ǚ ǜ) absentes de la police principale -> caractères actifs
\newfontface\pyfont{WenQuanYiZenHei-subset.ttf}[Path=../../../fonts/]
\newcommand\pyactive[1]{\catcode#1=13 \begingroup\lccode`\~=#1 \lowercase{\endgroup\def~}{{\pyfont\char#1}}}
\pyactive{"01CD}\pyactive{"01CE}\pyactive{"01CF}\pyactive{"01D0}\pyactive{"01D1}\pyactive{"01D2}\pyactive{"01D3}\pyactive{"01D4}\pyactive{"01D5}\pyactive{"01D6}\pyactive{"01D7}\pyactive{"01D8}\pyactive{"01D9}\pyactive{"01DA}\pyactive{"01DB}\pyactive{"01DC}
% Maths en sans-serif (Fira Math) avec chiffres et lettres droites de Plus
% Jakarta, pour que les formules se fondent dans le texte.
\usepackage[math-style=ISO,bold-style=ISO]{unicode-math}
\setmathfont{FiraMath-Regular.otf}[Path=../../../fonts/]
\setmathfont{PlusJakartaSans-Medium.ttf}[Path=../../../fonts/,range={up/num,up/{Latin,latin},\mathup}]
\usepackage{icomma} % virgule décimale sans espace en mode math
% Caractères mathématiques Unicode absents des polices de texte (Plus Jakarta,
% Archivo Black) : rendus par leur équivalent mathématique, en texte comme en
% formule — sinon ils s'affichent en carré vide (ex. « Arithmétique dans ℤ »).
\usepackage{newunicodechar}
\newunicodechar{ℝ}{\ensuremath{\mathbb{R}}}
\newunicodechar{ℤ}{\ensuremath{\mathbb{Z}}}
\newunicodechar{ℂ}{\ensuremath{\mathbb{C}}}
\newunicodechar{ℕ}{\ensuremath{\mathbb{N}}}
\newunicodechar{ℚ}{\ensuremath{\mathbb{Q}}}
\newunicodechar{𝔻}{\ensuremath{\mathbb{D}}}
\newunicodechar{α}{\ensuremath{\alpha}}
\newunicodechar{β}{\ensuremath{\beta}}
\newunicodechar{γ}{\ensuremath{\gamma}}
\newunicodechar{δ}{\ensuremath{\delta}}
\newunicodechar{ε}{\ensuremath{\varepsilon}}
\newunicodechar{θ}{\ensuremath{\theta}}
\newunicodechar{λ}{\ensuremath{\lambda}}
\newunicodechar{μ}{\ensuremath{\mu}}
\newunicodechar{σ}{\ensuremath{\sigma}}
\newunicodechar{τ}{\ensuremath{\tau}}
\newunicodechar{φ}{\ensuremath{\varphi}}
\newunicodechar{ψ}{\ensuremath{\psi}}
\newunicodechar{ω}{\ensuremath{\omega}}
\newunicodechar{Σ}{\ensuremath{\Sigma}}
% majuscules produites par \MakeUppercase dans les titres (α-aminés -> Α)
\newunicodechar{Α}{\ensuremath{\alpha}}
\newunicodechar{Β}{\ensuremath{\beta}}
\newunicodechar{Γ}{\ensuremath{\Gamma}}
\newunicodechar{Δ}{\ensuremath{\Delta}}
\newunicodechar{Ε}{\ensuremath{\varepsilon}}
\newunicodechar{Θ}{\ensuremath{\Theta}}
\newunicodechar{Λ}{\ensuremath{\Lambda}}
\newunicodechar{Μ}{\ensuremath{\mu}}
\newunicodechar{Τ}{\ensuremath{\tau}}
\newunicodechar{Φ}{\ensuremath{\Phi}}
\newunicodechar{Ψ}{\ensuremath{\Psi}}
\newunicodechar{Ω}{\ensuremath{\Omega}}
\newunicodechar{↦}{\ensuremath{\mapsto}}
\newunicodechar{∈}{\ensuremath{\in}}
\newunicodechar{∉}{\ensuremath{\notin}}
\newunicodechar{∧}{\ensuremath{\wedge}}
\newunicodechar{⇔}{\ensuremath{\Leftrightarrow}}
\newunicodechar{⟺}{\ensuremath{\Longleftrightarrow}}
\newunicodechar{⇒}{\ensuremath{\Rightarrow}}
\newunicodechar{⟹}{\ensuremath{\Longrightarrow}}
\newunicodechar{∘}{\ensuremath{\circ}}
\newunicodechar{∪}{\ensuremath{\cup}}
\newunicodechar{∩}{\ensuremath{\cap}}
\newunicodechar{≡}{\ensuremath{\equiv}}
\newunicodechar{⊂}{\ensuremath{\subset}}
\newunicodechar{⊥}{\ensuremath{\perp}}
\newunicodechar{∃}{\ensuremath{\exists}}
\newunicodechar{∀}{\ensuremath{\forall}}
\newunicodechar{∛}{\ensuremath{\sqrt[3]{\,}}}
\newunicodechar{∜}{\ensuremath{\sqrt[4]{\,}}}
\newunicodechar{⃗}{\ensuremath{{}^{\to}}}
\newunicodechar{ⁿ}{\ifmmode^{n}\else\textsuperscript{\ensuremath{n}}\fi}
\newunicodechar{ˣ}{\ifmmode^{x}\else\textsuperscript{\ensuremath{x}}\fi}
\newunicodechar{ᵇ}{\ifmmode^{b}\else\textsuperscript{\ensuremath{b}}\fi}
\newunicodechar{ʳ}{\ifmmode^{r}\else\textsuperscript{\ensuremath{r}}\fi}
\newunicodechar{ᵘ}{\ifmmode^{u}\else\textsuperscript{\ensuremath{u}}\fi}
\newunicodechar{ʸ}{\ifmmode^{y}\else\textsuperscript{\ensuremath{y}}\fi}
\newunicodechar{ᵃ}{\ifmmode^{a}\else\textsuperscript{\ensuremath{a}}\fi}
\newunicodechar{ᵉ}{\ifmmode^{e}\else\textsuperscript{\ensuremath{e}}\fi}
\newunicodechar{ᵏ}{\ifmmode^{k}\else\textsuperscript{\ensuremath{k}}\fi}
\newunicodechar{ᵖ}{\ifmmode^{p}\else\textsuperscript{\ensuremath{p}}\fi}
\newunicodechar{ᴺ}{\ifmmode^{N}\else\textsuperscript{\ensuremath{N}}\fi}
\newunicodechar{ᶜ}{\ifmmode^{c}\else\textsuperscript{\ensuremath{c}}\fi}
\newunicodechar{ʰ}{\ifmmode^{h}\else\textsuperscript{\ensuremath{h}}\fi}
\newunicodechar{ᵠ}{\ifmmode^{q}\else\textsuperscript{\ensuremath{q}}\fi}
\newunicodechar{ₙ}{\ifmmode_{n}\else\textsubscript{\ensuremath{n}}\fi}
\newunicodechar{ₐ}{\ifmmode_{a}\else\textsubscript{\ensuremath{a}}\fi}
\newunicodechar{ᵢ}{\ifmmode_{i}\else\textsubscript{\ensuremath{i}}\fi}
\newunicodechar{ᵣ}{\ifmmode_{r}\else\textsubscript{\ensuremath{r}}\fi}
\newunicodechar{ₖ}{\ifmmode_{k}\else\textsubscript{\ensuremath{k}}\fi}
\newunicodechar{ᵦ}{\ifmmode_{\beta}\else\textsubscript{\ensuremath{\beta}}\fi}
\newunicodechar{ₓ}{\ifmmode_{x}\else\textsubscript{\ensuremath{x}}\fi}
\newfontfamily\popdisplay{ArchivoBlack-Regular.ttf}[Path=../../../fonts/]
\newfontfamily\poplogo{SpaceGrotesk-Bold.ttf}[Path=../../../fonts/]
\newfontfamily\popsemi{PlusJakartaSans-SemiBold.ttf}[Path=../../../fonts/]
\newfontfamily\popxbold{PlusJakartaSans-ExtraBold.ttf}[Path=../../../fonts/]
\newfontfamily\popxboldls{PlusJakartaSans-ExtraBold.ttf}[Path=../../../fonts/,LetterSpace=3.0]

\setlist[enumerate]{leftmargin=1.7em,itemsep=5pt,topsep=4pt}
\setlist[itemize]{leftmargin=1.7em,itemsep=4pt,topsep=4pt,label={\color{poporange}\textbullet}}
\setlength{\parindent}{0pt}
\setlength{\parskip}{5pt}
\renewcommand{\arraystretch}{1.25}

% pgfplots : style Pop par défaut
\pgfplotsset{
  every axis/.append style={axis line style={popink,line width=1pt},tick style={popink},
    grid style={popline,line width=0.5pt},label style={font=\small},tick label style={font=\footnotesize}},
  popcurve/.style={poporange,line width=1.8pt,no markers,smooth},
}
% Même style disponible dans un \draw TikZ simple (hors pgfplots)
\tikzset{popcurve/.style={poporange,line width=1.8pt}}
\tikzset{popgrid/.style={popline,very thin}}
\colorlet{popcurve}{poporange} % le nom du style est parfois utilisé comme couleur

% ============================================================
% Pill / squiggle
% ============================================================
\newcommand{\poppill}[3]{%
  \tikz[baseline=-0.55ex]\node[draw=popink,line width=1.2pt,rounded corners=2.6pt,%
    fill=#2,inner xsep=6.5pt,inner ysep=3pt]{%
    \popxboldls\fontsize{7}{7}\selectfont\color{#3}\MakeUppercase{#1}};%
}
\newcommand{\popsquiggle}[1]{%
  \tikz[baseline=-0.4ex]{
    \draw[#1,line width=2.6pt,line cap=round]
      plot [smooth] coordinates {(0,0.09) (0.34,0.01) (0.68,0.21) (1.02,0.01) (1.36,0.10)};
  }%
}
% Mot-clé surligné (marqueur orange sous le texte)
\newcommand{\cle}[1]{{\popxbold\color{popdark}#1}}

% ============================================================
% En-tête / pied
% ============================================================
\pagestyle{fancy}
\fancyhf{}
\renewcommand{\headrulewidth}{0pt}
\renewcommand{\footrulewidth}{0pt}
\lhead{\raisebox{-0.15ex}{\poplogo\fontsize{13}{13}\selectfont\color{popink}OnBuch\color{poporange}+}}
\rhead{\poppill{\DOCMATIERE\ \textbullet\ \DOCNIVEAU\ \textbullet\ Leçon \DOCLECON}{white}{popink}}
\lfoot{\popxbold\fontsize{7.6}{7.6}\selectfont\color{popmuted}\MakeUppercase{Cours signé OnBuch+}\ \textbullet\ {\popsemi\DOCTITRE}}
\rfoot{\tikz[baseline=-0.6ex]\node[rounded corners=8.5pt,fill=popink,minimum height=17pt,minimum width=17pt,inner xsep=5pt,inner sep=0]{\popxbold\fontsize{8}{8}\selectfont\color{popcream}\thepage\,/\,\pageref*{LastPage}};}

% ============================================================
% Titres
% ============================================================
\newcommand{\chaptitre}[2]{%
  \begin{tcolorbox}[enhanced,colback=poporange,colframe=popink,boxrule=2.6pt,arc=11pt,
      drop shadow=popink,left=15pt,right=15pt,top=12pt,bottom=11pt,before skip=3pt,after skip=18pt]
    \poppill{#2}{popink}{popcream}\par\vspace{9pt}
    {\raggedright\hyphenpenalty=10000\exhyphenpenalty=10000\popdisplay\fontsize{22}{24}\selectfont\color{popcream}\MakeUppercase{#1}\par}
  \end{tcolorbox}
}
% Section numérotée : pastille orange avec le numéro + titre Archivo
\newcounter{popsec}
\newcommand{\coursec}[1]{%
  \stepcounter{popsec}\setcounter{popsub}{0}%
  \par\vspace{16pt}\noindent
  \begin{minipage}{\linewidth}
  \tikz[baseline=-0.7ex]\node[draw=popink,line width=1.6pt,fill=poporange,rounded corners=4pt,minimum size=22pt,inner sep=0]{\popdisplay\fontsize{12}{12}\selectfont\color{popcream}\thepopsec};%
  \hspace{8pt}{\popdisplay\fontsize{15}{17}\selectfont\color{popink}\MakeUppercase{#1}}\par
  \vspace{4pt}\noindent{\color{poporange}\rule{36pt}{3.6pt}}
  \end{minipage}\par\nopagebreak\vspace{8pt}
}
\newcounter{popsub}
\newcommand{\courssub}[1]{%
  \stepcounter{popsub}%
  \par\vspace{9pt}\noindent{\popxbold\fontsize{11.5}{13}\selectfont\color{popink}{\color{poporange}\thepopsec.\thepopsub}\hspace{6pt}#1}\par\nopagebreak\vspace{4pt}
}

% ============================================================
% Boîtes pédagogiques — même recette : fond blanc, bordure épaisse colorée,
% ombre dure encre, badge plein. Titre optionnel [#1].
% ============================================================
\tcbset{popbase/.style={breakable,enhanced,colback=white,boxrule=2.3pt,arc=9pt,
  shadow={3pt}{-3pt}{0pt}{popink},left=14pt,right=14pt,top=11pt,bottom=11pt,before skip=11pt,after skip=15pt,
  fontupper=\popsemi\fontsize{10.6}{14.5}\selectfont\color{popink},
  fontlower=\popsemi\fontsize{10.6}{14.5}\selectfont\color{popink}}}
% Coche dessinée (le glyphe ✓ n'existe pas dans les polices du document)
\newcommand{\popcheck}[1]{\tikz[baseline=-0.5ex]\draw[#1,line width=1.6pt,line cap=round,line join=round](-0.13,0.0)--(-0.04,-0.09)--(0.13,0.1);}
\providecommand{\coche}{\popcheck{popgreen}}
\providecommand{\resultat}[1]{\par\smallskip\noindent\fcolorbox{popgreen}{popgreenL}{\parbox{\dimexpr\linewidth-2\fboxsep-2\fboxrule\relax}{\textbf{Résultat.} #1}}\par\smallskip}
\newcommand{\popboxhead}[3]{\poppill{#1}{#2}{white}\ifx\relax#3\relax\else\hspace{6pt}{\popxbold\fontsize{10.6}{12}\selectfont\color{popink}#3}\fi\par\vspace{7pt}}

\newtcolorbox{aretenir}[1][]{popbase,colframe=popgreen,before upper={\popboxhead{À retenir}{popgreen}{#1}}}
% Texte en chinois (document de lecture, dialogue, extrait) : cadre
% d'étude. Un titre optionnel (source, auteur) s'affiche dans l'en-tête.
\newtcolorbox{textebox}[1][]{popbase,colframe=popblue,before upper={\popboxhead{文本 · Texte}{popblue}{#1}},after upper={}}
% Marques ✓ / ✗ (forme correcte / fautive) souvent utilisées par les modèles
\providecommand{\cross}{\textcolor{poppink}{\ensuremath{\times}}}
\providecommand{\xmark}{\cross}\providecommand{\crossmark}{\cross}\providecommand{\cmark}{\ensuremath{\checkmark}}
% Ligne de réponse / justification à compléter (inventée par les modèles)
\providecommand{\justification}[1]{\par\noindent\textit{Justification~:} \dotfill\par}
% \textipa (package tipa non chargé) : la police ne couvre pas l'API -> simple passage
\providecommand{\textipa}[1]{#1}
% Soulignement (ulem non chargé) : \uline, \ul
\providecommand{\uline}[1]{\underline{#1}}\providecommand{\ul}[1]{\underline{#1}}
% \glossaire{\item ...} inventée par les modèles : liste de glossaire
\providecommand{\glossaire}[1]{\begin{itemize}#1\end{itemize}}
% \alert (beamer) inventée par les modèles : texte mis en évidence
\providecommand{\alert}[1]{\textbf{\textcolor{poppink}{#1}}}
% variantes ASCII d'ordinaux écrites en macro
\providecommand{\iere}{\textsuperscript{re}}\providecommand{\ieme}{\textsuperscript{e}}\providecommand{\ere}{\textsuperscript{re}}\providecommand{\eme}{\textsuperscript{e}}\providecommand{\er}{\textsuperscript{er}}\providecommand{\ie}{\textsuperscript{e}}
% Ordinaux français écrits en macro par les modèles : 1\ière, 3\ième
\newcommand{\ière}{\textsuperscript{re}}\newcommand{\ième}{\textsuperscript{e}}
% \exemple (inventée par les modèles) : étiquette « Exemple : »
\providecommand{\exemple}{\textbf{Exemple~:}\ }
% \square utilisable aussi hors mode math (cases à cocher)
\AtBeginDocument{\let\squareMath\square\renewcommand{\square}{\ensuremath{\squareMath}}}
% Mot ou phrase en chinois à l'intérieur d'un paragraphe français
\newcommand{\zh}[1]{\ifmmode\text{#1}\else{#1}\fi}
\let\eng\zh \let\esp\zh \let\all\zh % alias tolérés
% flèches d'intonation inventées par les modèles
\providecommand{\textsearrow}{}\renewcommand{\textsearrow}{\ensuremath{\searrow}}\providecommand{\textnearrow}{}\renewcommand{\textnearrow}{\ensuremath{\nearrow}}
% \fr{..} : mot français (inventé par les modèles)
\providecommand{\fr}[1]{\textit{#1}}
% zhbox : encadré de texte chinois inventé par les modèles
\newenvironment{zhbox}{\begin{quote}}{\end{quote}}
% \courssubsub : titre de niveau 3 inventé par les modèles
\providecommand{\courssubsub}[1]{\par\medskip\noindent\textbf{#1}\par\smallskip}
% \zhbf : texte chinois en gras inventé par les modèles
\providecommand{\zhbf}[1]{\textbf{#1}}
% \vs : « versus » inventé par les modèles
\providecommand{\vs}{\textit{vs}\ }
% \blank : case à compléter inventée par les modèles
\providecommand{\blank}[1][]{\underline{\hspace{1.6em}}}
\newtcolorbox{exemplebox}[1][]{popbase,colframe=poporange,before upper={\popboxhead{Exemple}{poporange}{#1}}}
\newtcolorbox{definition}[1][]{popbase,colframe=popblue,before upper={\popboxhead{Définition}{popblue}{#1}}}
\newtcolorbox{propriete}[1][]{popbase,colframe=poppurple,before upper={\popboxhead{Propriété}{poppurple}{#1}}}
\newtcolorbox{methode}[1][]{popbase,colframe=popgold,before upper={\popboxhead{Méthode}{popgold}{#1}}}
\newtcolorbox{attention}[1][]{popbase,colframe=poppink,before upper={\popboxhead{Attention}{poppink}{#1}}}
\newtcolorbox{experience}[1][]{popbase,colframe=popblue,colback=popblueL,before upper={\popboxhead{Expérience}{popblue}{#1}}}
% « Le savais-tu ? » : carte violette pleine (contexte, culture, Cameroun)
\newtcolorbox{savaistu}[1][]{popbase,colback=poppurple,colframe=popink,
  fontupper=\popsemi\fontsize{10.4}{14.2}\selectfont\color{white},
  before upper={\poppill{Le savais-tu\,?}{white}{poppurple}\ifx\relax#1\relax\else\hspace{6pt}{\popxbold\color{white}#1}\fi\par\vspace{7pt}}}
% Exercice résolu : énoncé en haut, \tcblower puis la solution sur fond crème
\newcounter{popexo}\newcounter{popexr}
\newtcolorbox{exoresolu}[1][]{popbase,colframe=poporange,colbacklower=popcream,
  segmentation style={popink,line width=1.2pt,dashed},
  before upper={\stepcounter{popexr}\popboxhead{Exercice résolu \thepopexr}{poporange}{#1}},
  before lower={\poppill{Solution}{popink}{popcream}\par\vspace{6pt}}}
% Exercice (non résolu en place) — la correction est dans la partie Corrigés
\newtcolorbox{exercice}[1][]{popbase,colframe=popink,
  before upper={\stepcounter{popexo}\popboxhead{Exercice \thepopexo}{popink}{#1}}}
% Corrigé
\newtcolorbox{corrige}[1][]{popbase,colframe=popgreen,colback=popgreenL,
  before upper={\popboxhead{Corrigé}{popgreen}{#1}}}
% Objectifs / prérequis / bilan
\newtcolorbox{objectifs}{popbase,colframe=popink,colback=poporangeL,
  before upper={\popboxhead{Objectifs de la leçon}{popink}{}}}
\newtcolorbox{prerequis}{popbase,colframe=popink,colback=popblueL,
  before upper={\popboxhead{Prérequis}{popblue}{}}}
\newtcolorbox{bilan}{popbase,colframe=popink,colback=white,boxrule=2.8pt,
  before upper={\popboxhead{Fiche bilan}{poporange}{L'essentiel à connaître pour l'examen}}}
% Figure encadrée avec légende
\newtcolorbox{popfigure}[1][]{enhanced,breakable=false,colback=white,colframe=popline,boxrule=1.4pt,arc=8pt,
  left=10pt,right=10pt,top=10pt,bottom=8pt,before skip=10pt,after skip=12pt,halign=center,
  lower separated=false,halign lower=center,
  fontlower=\popsemi\fontsize{9}{11.5}\selectfont\color{popmuted},
  #1}
\newcommand{\legende}[1]{\par\vspace{6pt}{\popsemi\fontsize{9}{11.5}\selectfont\color{popmuted}#1}}

% ============================================================
% Signature OnBuch+
% ============================================================
\newcommand{\poplogoinline}[1]{{\poplogo\fontsize{#1}{#1}\selectfont\color{popink}OnBuch\color{poporange}+}}
\newcommand{\signatureonbuch}{%
  \begin{tcolorbox}[enhanced,colback=popink,colframe=popink,boxrule=2.2pt,arc=10pt,
      drop shadow=poporange,left=14pt,right=14pt,top=10pt,bottom=10pt,before skip=0pt,after skip=0pt]
    \tikz[baseline=-0.7ex]\node[circle,fill=popgreen,draw=popcream,line width=1.3pt,minimum size=22pt,inner sep=0]{\popcheck{white}};%
    \hspace{9pt}\begin{minipage}[c]{0.62\linewidth}
      {\popxbold\fontsize{12}{13}\selectfont\color{popcream}Signé\ \textbullet\ {\poplogo OnBuch\color{poporange}+}}\par\vspace{2pt}
      {\raggedright\popsemi\fontsize{8}{10.5}\selectfont\color{popline}\MakeUppercase{Équipe pédagogique OnBuch+}\\\MakeUppercase{Conforme au programme officiel MINESEC}\par}
    \end{minipage}\hfill
    {\poplogo\fontsize{22}{22}\selectfont\color{popcream}OnBuch\color{poporange}+}
  \end{tcolorbox}%
}

% ============================================================
% Page de garde
% #1 titre · #2 matière · #3 niveau · #4 module · #5 n° leçon · #6 durée ·
% #7 description / objectifs (texte ou itemize)
% ============================================================
\newcommand{\pagegarde}[7]{%
  \thispagestyle{empty}
  \newgeometry{a4paper,margin=1.7cm,top=1.6cm,bottom=1.4cm}%
  \begin{tikzpicture}[remember picture,overlay]
    \fill[poporange!18] ([xshift=-3.2cm,yshift=-3.6cm]current page.north east) circle (4.6cm);
    \fill[poppurple!14] ([xshift=2.4cm,yshift=7.6cm]current page.south west) circle (3.8cm);
  \end{tikzpicture}%
  {\poplogo\fontsize{30}{30}\selectfont\color{popink}OnBuch\color{poporange}+}\hfill
  \raisebox{0.6ex}{\poppill{Cours signé \textbullet\ Premium}{popink}{popcream}}\par
  \vspace{1pt}\popsquiggle{poppurple}\par
  \vspace{8pt}
  \begin{tcolorbox}[enhanced,colback=poporange,colframe=popink,boxrule=2.8pt,arc=13pt,
      drop shadow=popink,left=18pt,right=18pt,top=12pt,bottom=13pt,after skip=10pt]
    \poppill{#2 \textbullet\ #3}{popink}{popcream}\hspace{5pt}\poppill{#4}{white}{popink}\par\vspace{8pt}
    {\popdisplay\fontsize{14}{14}\selectfont\color{popink}LEÇON #5}\par\vspace{4pt}
    {\raggedright\hyphenpenalty=10000\exhyphenpenalty=10000\popdisplay\fontsize{25}{27}\selectfont\color{popcream}\MakeUppercase{#1}\par}
    \vspace{8pt}
    {\popxbold\fontsize{9.5}{9.5}\selectfont\color{popink}\MakeUppercase{Durée conseillée : #6}}
  \end{tcolorbox}
  \begin{tcolorbox}[enhanced,colback=white,colframe=popink,boxrule=2.2pt,arc=10pt,
      drop shadow=popink,left=16pt,right=16pt,top=10pt,bottom=10pt,after skip=8pt]
    \poppill{Dans cette leçon}{poppurple}{white}\par\vspace{5pt}
    \popsemi\fontsize{9.3}{12.6}\selectfont\color{popink}#7
  \end{tcolorbox}
  \noindent\begin{minipage}[t]{0.487\linewidth}
    \begin{tcolorbox}[enhanced,colback=popgreen,colframe=popink,boxrule=2.2pt,arc=10pt,
        drop shadow=popink,left=11pt,right=11pt,top=10pt,bottom=10pt,height=3.1cm,valign=center]
      \tcbox[colback=white,colframe=popink,boxrule=1.3pt,arc=5pt,boxsep=0pt,nobeforeafter,
        left=3pt,right=3pt,top=3pt,bottom=3pt,valign=center]{\qrcode[height=1.9cm]{https://play.google.com/store/apps/details?id=com.luvvix.onbuch&hl=fr}}%
      \hspace{8pt}\begin{minipage}[c]{0.5\linewidth}
        {\popdisplay\fontsize{11}{12.5}\selectfont\color{white}\MakeUppercase{Télécharge OnBuch}}\par\vspace{4pt}
        {\raggedright\popsemi\fontsize{8.4}{11}\selectfont\color{white}Gratuit \textbullet\ Google Play\\Tuteur IA, annales, concours, résultats\par}
      \end{minipage}
    \end{tcolorbox}
  \end{minipage}\hfill
  \begin{minipage}[t]{0.487\linewidth}
    \begin{tcolorbox}[enhanced,colback=poppurple,colframe=popink,boxrule=2.2pt,arc=10pt,
        drop shadow=popink,left=11pt,right=11pt,top=10pt,bottom=10pt,height=3.1cm,valign=center]
      \tikz[baseline=-0.7ex]\node[circle,fill=white,draw=popink,line width=1.3pt,minimum size=1.6cm,inner sep=0]{\popdisplay\fontsize{24}{24}\selectfont\color{poppurple}+};%
      \hspace{8pt}\begin{minipage}[c]{0.6\linewidth}
        {\popdisplay\fontsize{11}{12.5}\selectfont\color{white}\MakeUppercase{Passe à OnBuch+}}\par\vspace{4pt}
        {\raggedright\popsemi\fontsize{8.4}{11}\selectfont\color{white}Tous les cours signés, Tuteur IA illimité, fiches PDF — onglet OnBuch+ de l'app\par}
      \end{minipage}
    \end{tcolorbox}
  \end{minipage}
  \vspace{8pt}
  \popcontenu
  \vfill
  \signatureonbuch
  \restoregeometry
  \newpage
}

% Rangée « Ce cours contient » (page de garde)
\newcommand{\poptile}[3]{% #1 couleur #2 gros texte #3 légende
  \begin{tcolorbox}[enhanced,colback=#1,colframe=popink,boxrule=2pt,arc=9pt,shadow={2.5pt}{-2.5pt}{0pt}{popink},
      left=6pt,right=6pt,top=9pt,bottom=9pt,halign=center,valign=center,height=1.75cm,width=\linewidth,nobeforeafter]
    {\popdisplay\fontsize{12.5}{13}\selectfont\color{white}\MakeUppercase{#2}}\par\vspace{3pt}
    {\centering\popsemi\fontsize{7.8}{9.5}\selectfont\color{white}#3\par}
  \end{tcolorbox}}
\newcommand{\popcontenu}{%
  \par\noindent{\popxboldls\fontsize{7.5}{7.5}\selectfont\color{popmuted}CE COURS SIGNÉ CONTIENT}\par\vspace{7pt}
  \noindent\begin{minipage}[t]{0.235\linewidth}\poptile{popblue}{Cours}{complet, illustré, pas à pas}\end{minipage}\hfill
  \begin{minipage}[t]{0.235\linewidth}\poptile{poporange}{Méthodes}{et exercices résolus}\end{minipage}\hfill
  \begin{minipage}[t]{0.235\linewidth}\poptile{poppink}{Exercices}{type examen + corrigés}\end{minipage}\hfill
  \begin{minipage}[t]{0.235\linewidth}\poptile{popgreen}{Fiche bilan}{l'essentiel à retenir}\end{minipage}\par}

% Sommaire
\newtcolorbox{sommaire}{enhanced,colback=white,colframe=popink,boxrule=2.2pt,arc=10pt,
  drop shadow=popink,left=18pt,right=18pt,top=13pt,bottom=13pt,before skip=6pt,after skip=20pt,
  before upper={\poppill{Sommaire}{poppurple}{white}\par\vspace{9pt}\popsemi\fontsize{10.6}{14.5}\selectfont\color{popink}}}

% Signature de fin de cours — poussée tout en bas de la dernière page
\newcommand{\findecours}{%
  \par\vfill\nopagebreak
  \begin{center}\popsquiggle{poporange}\end{center}\vspace{4pt}
  \signatureonbuch
}
\AtBeginDocument{\def\llbracket{\mathopen{[\mkern-3.5mu[}}\def\rrbracket{\mathclose{]\mkern-3.5mu]}}} % intervalles d'entiers (glyphe ⟦ absent de la police)
% Glyphes absents de Fira Math : redéfinis à partir de glyphes disponibles
\AtBeginDocument{%
  \def\setminus{\mathbin{\backslash}}\let\smallsetminus\setminus
  \def\vdots{\mathord{\vbox{\baselineskip3.2pt\lineskiplimit0pt\kern4pt\hbox{.}\hbox{.}\hbox{.}}}}%
  \def\ddots{\mathinner{\mkern1mu\raise6.5pt\hbox{.}\mkern2mu\raise3.6pt\hbox{.}\mkern2mu\raise0.7pt\hbox{.}\mkern1mu}}%
  \def\triangle{\mathord{\scalebox{0.9}{$\symup{\Delta}$}}}%
  \def\nabla{\mathord{\raisebox{\depth}{\scalebox{1}[-1]{$\symup{\Delta}$}}}}%
  \def\checkmark{\popcheck{popdark}}%
  \def\star{\mathbin{\ast}}\def\bigstar{\mathord{\ast}}%
  \def\diamond{\mathbin{\scalebox{0.6}{$\Diamond$}}}%
  \def\bigcup{\mathop{\mathchoice{\raisebox{-0.3em}{\scalebox{1.7}{$\cup$}}}{\scalebox{1.25}{$\cup$}}{\cup}{\cup}}\displaylimits}%
  \def\bigcap{\mathop{\mathchoice{\raisebox{-0.3em}{\scalebox{1.7}{$\cap$}}}{\scalebox{1.25}{$\cap$}}{\cap}{\cap}}\displaylimits}%
  \def\lesseqqgtr{\mathrel{\lessgtr}}\def\bigoplus{\mathop{\oplus}\displaylimits}%
}
\providecommand{\ssi}{si et seulement si} % abréviation fréquente
\AtBeginDocument{\providecommand{\wideparen}{\overparen}} % arc de cercle

\begin{document}

\pagegarde{Leçon 26 — Expression écrite : caractères complexes liés à l'exode rural}{Chinois}{1ère A}{Module 3 : Citoyenneté et ouverture au monde}{ZHA26}{2 h}{Cette leçon d'expression écrite entraîne les élèves à produire et traduire un texte simple sur l'exode rural en chinois. Elle insiste sur l'écriture correcte des caractères complexes du thème (traits, radicaux, confusion entre caractères proches) et sur l'usage des connecteurs.
\vspace{4pt}
\begin{multicols}{2}\small
\begin{itemize}
\item À la fin de cette leçon, je sais écrire correctement les caractères complexes du thème de l'exode rural (traits, radicaux, ordre).
\item Je sais distinguer les caractères proches faciles à confondre (村/材, 搬/般, 挤/济).
\item Je sais analyser un modèle de texte simple sur l'exode rural et repérer sa structure.
\item Je sais utiliser les connecteurs 因为…所以…, 但是, 而且, 为了 dans une production écrite.
\item Je sais produire un texte simple de 6 à 10 phrases en chinois sur l'exode rural.
\item Je sais traduire des phrases simples du français vers le chinois (thème).
\end{itemize}\end{multicols}}

\begin{sommaire}
\begin{multicols}{2}
\begin{enumerate}
\item Modèle de texte commenté
\item Lexique écrit du thème
\item Écriture des caractères complexes
\item Structure et connecteurs du texte
\item Thème et version guidés
\item Production écrite et évaluation
\item Activité d'intégration
\item Méthodes et exercices résolus
\item Exercices
\item Corrigés des exercices
\item Fiche bilan
\end{enumerate}
\end{multicols}
\end{sommaire}

\begin{objectifs}
\begin{itemize}
\item À la fin de cette leçon, je sais écrire correctement les caractères complexes du thème de l'exode rural (traits, radicaux, ordre).
\item Je sais distinguer les caractères proches faciles à confondre (村/材, 搬/般, 挤/济).
\item Je sais analyser un modèle de texte simple sur l'exode rural et repérer sa structure.
\item Je sais utiliser les connecteurs 因为…所以…, 但是, 而且, 为了 dans une production écrite.
\item Je sais produire un texte simple de 6 à 10 phrases en chinois sur l'exode rural.
\item Je sais traduire des phrases simples du français vers le chinois (thème).
\item Je sais traduire un court texte du chinois vers le français (version).
\item Je sais m'auto-évaluer à l'aide d'une grille d'évaluation.
\end{itemize}
\end{objectifs}

\begin{prerequis}
\begin{itemize}
\item Vocabulaire de base de la ville et du village (城市, 农村, 工作, 生活) vu en début d'année
\item Radicaux courants (亻, 土, 木, 扌, 辶) et règles générales d'ordre des traits
\item Structure de base Sujet + Verbe + Complément et particule 了
\item Connecteurs simples 和, 因为…所以… vus en Seconde
\item Nombres et expressions de quantité (很多, 越来越…)
\end{itemize}
\end{prerequis}

\newpage

\coursec{Modèle de texte commenté}

\courssub{Situation d'accroche : partir du village}

À Maroua comme à Bamenda, de nombreux jeunes quittent leur village pour Douala ou Yaoundé à la recherche d'un emploi : c'est \cle{l'exode rural}. En Chine aussi, des millions de paysans ont rejoint les villes depuis les années 1980. Comment raconter ce phénomène en chinois, avec des caractères correctement tracés ? Cette leçon t'apprend à écrire et à traduire un texte simple sur l'exode rural, en maîtrisant les caractères complexes du thème comme \zh{搬} (déménager) ou \zh{拥挤} (bondé).

\textbf{Problématique :} comment un texte chinois simple raconte-t-il un départ de la campagne vers la ville, et quels mots-outils (connecteurs) organisent ce récit ?

\courssub{Lecture du texte modèle}

Avant toute analyse, lis le texte en entier, à voix haute, sans t'arrêter sur les caractères difficiles. Cherche d'abord le sens global : \emph{qui ? où ? pourquoi ?}

\begin{textebox}[我的表哥离开了农村 — Wǒ de biǎogē líkāi le nóngcūn]
我的表哥以前住在农村。因为农村工作太少，所以他去年搬到了雅温得。城市的生活很方便，但是房租很贵，人也很拥挤。他想先找工作，再学习中文。为了家人，他每天都很努力。\\
Wǒ de biǎogē yǐqián zhù zài nóngcūn. Yīnwèi nóngcūn gōngzuò tài shǎo, suǒyǐ tā qùnián bān dào le Yǎwēndé. Chéngshì de shēnghuó hěn fāngbiàn, dànshì fángzū hěn guì, rén yě hěn yōngjǐ. Tā xiǎng xiān zhǎo gōngzuò, zài xuéxí Zhōngwén. Wèile jiārén, tā měitiān dōu hěn nǔlì.\\
\emph{Traduction naturelle :} « Mon cousin habitait autrefois à la campagne. Comme il y avait trop peu de travail à la campagne, il a déménagé l'année dernière à Yaoundé. La vie en ville est très pratique, mais le loyer est cher et il y a beaucoup de monde. Il veut d'abord trouver un travail, puis apprendre le chinois. Pour sa famille, il travaille dur chaque jour. »
\end{textebox}

\textbf{Glossaire du texte modèle}

\begin{center}
\begin{tabularx}{\linewidth}{|X|X|X|X|}
\hline
\rowcolor{popblueL} Caractères & Pinyin & Nature & Traduction \\
\hline
表哥 & biǎogē & nom & cousin (aîné, côté maternel) \\
\hline
以前 & yǐqián & nom de temps & autrefois, avant \\
\hline
住 & zhù & verbe & habiter, résider \\
\hline
农村 & nóngcūn & nom & campagne, monde rural \\
\hline
因为 & yīnwèi & conjonction & parce que \\
\hline
所以 & suǒyǐ & conjonction & donc, c'est pourquoi \\
\hline
搬 & bān & verbe & déménager, déplacer \\
\hline
城市 & chéngshì & nom & ville \\
\hline
方便 & fāngbiàn & adjectif & pratique, commode \\
\hline
房租 & fángzū & nom & loyer \\
\hline
拥挤 & yōngjǐ & adjectif & bondé, surpeuplé \\
\hline
先 & xiān & adverbe & d'abord \\
\hline
再 & zài & adverbe & ensuite, puis \\
\hline
为了 & wèile & préposition & pour, afin de \\
\hline
努力 & nǔlì & adjectif/verbe & (faire) des efforts, travailler dur \\
\hline
\end{tabularx}
\end{center}

\courssub{Traduction mot à mot : comprendre la mécanique}

Comparer la traduction mot à mot et la traduction naturelle permet de voir comment le chinois « pense » la phrase.

\begin{exemplebox}[Mot à mot contre traduction naturelle]
\zh{我的表哥以前住在农村。}\\
\emph{Wǒ de biǎogē yǐqián zhù zài nóngcūn.}\\
Mot à mot : « moi + (possessif 的) + cousin + autrefois + habiter + à + campagne ».\\
Traduction naturelle : « Mon cousin habitait autrefois à la campagne. »\\[0.4em]
\zh{因为农村工作太少，所以他去年搬到了雅温得。}\\
\emph{Yīnwèi nóngcūn gōngzuò tài shǎo, suǒyǐ tā qùnián bān dào le Yǎwēndé.}\\
Mot à mot : « parce que + campagne + travail + trop + peu, donc + lui + année passée + déménager-vers + Yaoundé ».\\
Traduction naturelle : « Parce qu'il y avait trop peu de travail à la campagne, il a déménagé l'année dernière à Yaoundé. »\\[0.4em]
\zh{他想先找工作，再学习中文。}\\
\emph{Tā xiǎng xiān zhǎo gōngzuò, zài xuéxí Zhōngwén.}\\
Mot à mot : « lui + vouloir + d'abord + chercher + travail, ensuite + étudier + chinois ».\\
Traduction naturelle : « Il veut d'abord trouver un travail, puis apprendre le chinois. »
\end{exemplebox}

\begin{attention}[Le temps n'est pas dans le verbe]
En français, « habitait », « a déménagé », « veut » portent le temps. En chinois, le verbe ne change jamais : c'est le mot de temps (\zh{以前} autrefois, \zh{去年} l'année dernière) ou la particule \zh{了} qui situe l'action. Ne cherche donc jamais à « conjuguer » un verbe chinois : place le bon mot de temps au début de la phrase.
\end{attention}

\courssub{Repérage de la structure du texte}

Un bon texte simple en chinois suit une progression logique claire. Le texte modèle s'organise en six étapes :

\begin{enumerate}
\item \cle{Situation passée} : \zh{我的表哥以前住在农村。} (où vivait-il avant ?)
\item \cle{Cause} : \zh{因为农村工作太少…} (pourquoi est-il parti ?)
\item \cle{Départ} : \zh{…所以他去年搬到了雅温得。} (où est-il allé ?)
\item \cle{Contraste ville} : \zh{城市的生活很方便，但是房租很贵…} (avantages et inconvénients)
\item \cle{Projet} : \zh{他想先找工作，再学习中文。} (que veut-il faire ?)
\item \cle{Motivation} : \zh{为了家人，他每天都很努力。} (pour qui se bat-il ?)
\end{enumerate}

\begin{popfigure}
\begin{tikzpicture}[
  node distance=0.35cm,
  box/.style={rectangle, rounded corners=3pt, draw=popblue, fill=popblueL, align=center, font=\small, minimum width=3.1cm, minimum height=0.8cm},
  arr/.style={-{Stealth[length=2.5mm]}, thick, popdark}
]
\node[box, align=center] (n1){1. Situation passée\\以前住在农村};
\node[box, below=of n1, fill=poporangeL, draw=poporange, align=center] (n2){2. Cause\\因为…工作太少};
\node[box, below=of n2, fill=popgreenL, draw=popgreen, align=center] (n3){3. Départ\\所以…搬到了雅温得};
\node[box, below=of n3, fill=poppinkL, draw=poppink, align=center] (n4){4. Contraste ville\\方便，但是房租贵};
\node[box, below=of n4, fill=poppurpleL, draw=poppurple, align=center] (n5){5. Projet\\先找工作，再学习中文};
\node[box, below=of n5, fill=popgoldL, draw=popgold, align=center] (n6){6. Motivation\\为了家人…努力};
\draw[arr] (n1) -- (n2);
\draw[arr] (n2) -- (n3);
\draw[arr] (n3) -- (n4);
\draw[arr] (n4) -- (n5);
\draw[arr] (n5) -- (n6);
\end{tikzpicture}
\legende{Organigramme de la structure du texte modèle : six étapes logiques, de la situation passée à la motivation.}
\end{popfigure}

\courssub{Les connecteurs du texte : la charpente du récit}

Les \cle{connecteurs} sont les mots-outils qui relient les idées. Repérons-les dans le texte, puis énonçons les règles.

\begin{propriete}[因为…所以… : parce que… donc…]
Structure : \textbf{因为 + cause ，所以 + conséquence}.\\
Contrairement au français, où l'on dit « parce que… » \emph{ou} « donc… » mais pas les deux ensemble, le chinois emploie volontiers les deux dans la même phrase. La virgule chinoise \zh{，} sépare les deux parties.\\
Emploi : exprimer la cause d'un événement (ici, la cause de l'exode rural).
\end{propriete}

\begin{exemplebox}[因为…所以… en action]
\zh{因为农村工作太少，所以他搬到了城市。}\\
\emph{Yīnwèi nóngcūn gōngzuò tài shǎo, suǒyǐ tā bān dào le chéngshì.}\\
« Parce qu'il y avait trop peu de travail à la campagne, il a déménagé en ville. »\\[0.4em]
\zh{因为房租很贵，所以他住得很远。}\\
\emph{Yīnwèi fángzū hěn guì, suǒyǐ tā zhù de hěn yuǎn.}\\
« Parce que le loyer est cher, il habite très loin. »
\end{exemplebox}

\begin{propriete}[Quatre autres connecteurs du texte]
\begin{itemize}
\item \zh{但是} (dànshì) « mais » : introduit le contraste. \zh{城市的生活很方便，但是房租很贵。} « La vie en ville est pratique, mais le loyer est cher. »
\item \zh{也} (yě) « aussi » : se place \emph{avant} le verbe ou l'adjectif, jamais en fin de phrase comme le français « aussi ». \zh{人也很拥挤。} « Il y a aussi beaucoup de monde. »
\item \zh{先…再…} (xiān… zài…) « d'abord… puis… » : ordonne les actions d'un projet. \zh{先找工作，再学习中文。} « d'abord trouver un travail, puis étudier le chinois ».
\item \zh{为了} (wèile) « pour, afin de » : suivi d'un nom ou d'un groupe nominal, exprime le but. \zh{为了家人} « pour sa famille ».
\end{itemize}
\end{propriete}

\begin{center}
\begin{tabularx}{\linewidth}{|X|X|X|}
\hline
\rowcolor{popblueL} Connecteur & Exemple (caractères + pinyin) & Traduction \\
\hline
因为…所以… & 因为工作少，所以他走了。Yīnwèi gōngzuò shǎo, suǒyǐ tā zǒu le. & Parce qu'il y avait peu de travail, il est parti. \\
\hline
但是 & 城市很大，但是很贵。Chéngshì hěn dà, dànshì hěn guì. & La ville est grande, mais chère. \\
\hline
也 & 我也想去城市。Wǒ yě xiǎng qù chéngshì. & Moi aussi, je veux aller en ville. \\
\hline
先…再… & 先吃饭，再学习。Xiān chīfàn, zài xuéxí. & D'abord manger, puis étudier. \\
\hline
为了 & 为了孩子，她很努力。Wèile háizi, tā hěn nǔlì. & Pour ses enfants, elle travaille dur. \\
\hline
\end{tabularx}
\end{center}

\courssub{Carte mentale du thème : l'exode rural}

Pour préparer ta propre production écrite, organise le lexique autour du thème en quatre branches : les causes, la campagne, la ville, les projets.

\begin{popfigure}
\begin{tikzpicture}[
  mindmap,
  grow cyclic,
  every node/.style={concept, align=center, font=\small},
  root concept/.append style={fill=popblue, font=\bfseries},
  level 1/.append style={level distance=3.2cm, sibling angle=90},
  level 2/.append style={level distance=2.2cm, sibling angle=45, font=\footnotesize}
]
\node[root concept, align=center]{exode rural\\农村 → 城市}
  child[concept color=poporangeL]{ node[align=center]{原因\\causes}
    child { node[align=center]{工作少\\peu de travail} }
    child { node[align=center]{工资低\\salaires bas} }
  }
  child[concept color=popgreenL]{ node[align=center]{农村\\campagne}
    child { node[align=center]{家人\\famille} }
    child { node[align=center]{房子\\maison} }
  }
  child[concept color=poppinkL]{ node[align=center]{城市\\ville}
    child { node[align=center]{房租贵\\loyer cher} }
    child { node[align=center]{很拥挤\\bondé} }
    child { node[align=center]{很方便\\pratique} }
  }
  child[concept color=poppurpleL]{ node[align=center]{计划\\projets}
    child { node[align=center]{找工作\\trouver un emploi} }
    child { node[align=center]{学中文\\apprendre le chinois} }
  };
\end{tikzpicture}
\legende{Carte mentale lexicale : quatre branches pour construire un texte sur l'exode rural.}
\end{popfigure}

\courssub{Lexique écrit étendu du thème}

Le tableau ci-dessous regroupe le vocabulaire clé de la carte mentale, indispensable pour l'expression écrite et la traduction.

\begin{center}
\begin{tabularx}{\linewidth}{|X|X|X|X|}
\hline
\rowcolor{popblueL} Caractères & Pinyin & Nature & Traduction \\
\hline
工资 & gōngzī & nom & salaire, paye \\
\hline
低 & dī & adjectif & bas (prix, niveau) \\
\hline
家人 & jiārén & nom & famille, proches \\
\hline
房子 & fángzi & nom & maison, logement \\
\hline
找 & zhǎo & verbe & chercher, trouver \\
\hline
学习 & xuéxí & verbe & étudier, apprendre \\
\hline
中文 & Zhōngwén & nom & chinois (langue) \\
\hline
每天 & měitiān & nom de temps & chaque jour, tous les jours \\
\hline
去年 & qùnián & nom de temps & l'année dernière \\
\hline
生活 & shēnghuó & nom/verbe & vie, vivre \\
\hline
\end{tabularx}
\end{center}

\courssub{Écriture des caractères complexes}

Cette section cible les caractères à forte charge graphique du thème. Maîtriser leur décomposition (radical + phonétique) et l'ordre des traits évite les confusions fréquentes.

\begin{definition}[Caractère \zh{搬}  — bān, déménager]
\textbf{Radical :} \zh{扌} (main, main gauche) — action manuelle.\\
\textbf{Phonétique :} \zh{般} (bān) — donne le son.\\
\textbf{Structure :} Gauche-Droite.\\
\textbf{Ordre des traits (13 traits) :} 1. Horizontal (扌) 2. Crochet vertical (扌) 3. Trait montant (扌) → 4. Horizontal (般) 5. Vertical (般) 6. Horizontal (般) 7. Vertical (般) 8. Horizontal (般) 9. Point (般) 10. Horizontal (般) 11. Vertical (般) 12. Crochet horizontal (般) 13. Point (般).\\
\textbf{Astuce :} « La \zh{主} (main) \zh{般} (porte/gère) les affaires pour déménager. » Ne pas confondre avec \zh{般} seul (sorte, genre) ni \zh{盘} (pán, plateau, radical \zh{皿}).
\end{definition}

\begin{definition}[Caractère \zh{拥}  — yōng, embrasser / bondé (dans \zh{拥挤})]
\textbf{Radical :} \zh{扌} (main).\\
\textbf{Phonétique :} \zh{用} (yòng) — son proche (yōng vs yòng).\\
\textbf{Structure :} Gauche-Droite.\\
\textbf{Ordre des traits (19 traits) :} 1-3. Radical \zh{扌} (horizontal, crochet vertical, trait montant) → 4-19. \zh{用} : horizontal, vertical, horizontal, vertical, horizontal, horizontal, horizontal, point, horizontal, vertical, crochet horizontal, point, point, point, point, horizontal.\\
\textbf{Attention :} Ne pas écrire le bas de \zh{用} comme \zh{甘} (gān). Le bas de \zh{用} est deux traits « point + horizontal » répétés.
\end{definition}

\begin{definition}[Caractère \zh{挤}  — jǐ, serré, bondé / pousser]
\textbf{Radical :} \zh{扌} (main).\\
\textbf{Phonétique :} \zh{且} (qiě) — son différent mais composant graphique.\\
\textbf{Structure :} Gauche-Droite.\\
\textbf{Ordre des traits (15 traits) :} 1-3. Radical \zh{扌} → 4-15. \zh{且} : horizontal, vertical, horizontal, horizontal, horizontal, vertical, horizontal, horizontal, horizontal, vertical, crochet horizontal, point.\\
\textbf{Caractères proches :} \zh{齐} (qí, régulier — haut différent), \zh{剂} (jì, dose — radical \zh{刂} à droite).
\end{definition}

\begin{definition}[Caractère \zh{租}  — zū, louer, loyer]
\textbf{Radical :} \zh{禾} (hé, grain/riz) — lien avec la rente foncière historique (payer en grain).\\
\textbf{Phonétique :} \zh{且} (qiě).\\
\textbf{Structure :} Haut-Bas.\\
\textbf{Ordre des traits (13 traits) :} 1-5. \zh{禾} (horizontal, vertical, horizontal, horizontal, point) → 6-13. \zh{且} (horizontal, vertical, horizontal, horizontal, horizontal, vertical, horizontal, crochet horizontal).\\
\textbf{Ne pas confondre :} \zh{组} (zǔ, groupe — radical \zh{糸} soie), \zh{祖} (zǔ, ancêtre — radical \zh{示} autel).
\end{definition}

\begin{definition}[Caractères \zh{雅温得}  — Yǎwēndé, Yaoundé (transcription phonétique)]
\begin{itemize}
\item \zh{雅} (yǎ, élégant) : radical \zh{隹} (oiseau court), structure Haut-Bas. Ordre : haut \zh{牙} (dent) puis \zh{隹}.
\item \zh{温} (wēn, tiède) : radical \zh{氵} (eau), structure Gauche-Droite. Phonétique \zh{昷}. Ordre : \zh{氵} (point, point, trait montant) puis \zh{昷} (soleil \zh{日} dans \zh{皿} bol).
\item \zh{得} (dé, particule / děi, devoir / de, obtenir) : ici son \emph{dé} (transcription). Structure Gauche-Droite : \zh{彳} (pas) + \zh{寸} (pouce/main). Ordre : \zh{彳} (3 traits) puis \zh{寸} (3 traits).
\end{itemize}
\end{definition}

\begin{aretenir}[Règles d'ordre des traits rappelées pour ces caractères]
\begin{enumerate}
\item Horizontal avant vertical (ex: \zh{十} dans \zh{般}, \zh{用}, \zh{且}).
\item Gauche avant droite (ex: radical \zh{扌} avant phonétique dans \zh{搬}, \zh{拥}, \zh{挤}).
\item Haut avant bas (ex: \zh{禾} avant \zh{且} dans \zh{租} ; \zh{牙} avant \zh{隹} dans \zh{雅}).
\item Extérieur avant intérieur (ex: \zh{皿} dans \zh{昷} de \zh{温}).
\item Le trait central avant les ailes symétriques (ex: verticale centrale de \zh{用}, \zh{且}).
\end{enumerate}
\end{aretenir}

\courssub{Méthode : lire, repérer, traduire}

\begin{methode}[Comment aborder un texte chinois simple]
\begin{enumerate}
\item \textbf{Lecture globale} : lis le texte en entier, à voix haute, sans t'arrêter. Repère \emph{qui} fait \emph{quoi}, \emph{où} et \emph{quand}.
\item \textbf{Repérage des mots de temps} : entoure \zh{以前}, \zh{去年}, \zh{每天} : ils donnent la chronologie.
\item \textbf{Soulignage des connecteurs} : souligne \zh{因为…所以…}, \zh{但是}, \zh{也}, \zh{先…再…}, \zh{为了} avant de traduire : ils révèlent la logique du texte.
\item \textbf{Traduction par blocs} : traduis segment par segment (sujet / verbe / complément), puis reformule en français naturel.
\item \textbf{Vérification} : relis ta traduction en français seul ; si elle sonne faux, reviens au chinois.
\end{enumerate}
\end{methode}

\courssub{Thème et version guidés}

\begin{methode}[Thème (Français → Chinois) : 5 étapes]
\begin{enumerate}
\item \textbf{Analyser} : repère le sujet, le temps (passé/présent/futur), le lieu, l'action, le but.
\item \textbf{Ordonner} : applique l'ordre chinois \textbf{Sujet + Temps + Lieu + Verbe + (Objet/Complément)}.
\item \textbf{Choisir les connecteurs} : \zh{因为…所以…} (cause), \zh{但是} (contraste), \zh{先…再…} (séquence), \zh{为了} (but).
\item \textbf{Écrire les caractères} : vérifie l'ordre des traits, le radical, le classificateur si nombre.
\item \textbf{Relire} : vérifie la ponctuation chinoise (， 。) et les tons du pinyin.
\end{enumerate}
\end{methode}

\begin{methode}[Version (Chinois → Français) : 5 étapes]
\begin{enumerate}
\item \textbf{Segmenter} : coupe aux virgules chinoises \zh{，} et au point \zh{。}. Identifie les blocs [Sujet + Prédicat].
\item \textbf{Repérer les marqueurs} : entoure \zh{了} (accompli), \zh{过} (expérience), \zh{着} (duratif), mots de temps (\zh{去年}, \zh{以前}), connecteurs.
\item \textbf{Traduire littéralement} : mot à mot sous le chinois pour voir la structure.
\item \textbf{Conjuguer en français} : choisis le bon temps (imparfait, passé composé, présent, futur) selon le marqueur chinois.
\item \textbf{Polir} : remplace les connecteurs doubles (\zh{因为…所以…} → « Parce que… » \emph{ou} « …, donc… ») pour un français naturel.
\end{enumerate}
\end{methode}

\begin{exoresolu}[Thème guidé : phrase unique]
Énoncé : Traduis en chinois : « Parce que les salaires sont bas à la campagne, il veut d'abord trouver un logement en ville. »
\tcblower
Solution détaillée :
\begin{enumerate}
\item \textbf{Analyse} : Cause = salaires bas campagne. Conséquence = il veut d'abord trouver logement ville. Temps = présent/général.
\item \textbf{Mots-clés} : 因为 (yīnwèi), 农村, 工资, 低 (dī), 所以 (suǒyǐ), 他 (tā), 想 (xiǎng), 先 (xiān), 找 (zhǎo), 房子 (fángzi), 在 (zài), 城市.
\item \textbf{Structure} : [因为 + Cause], [所以 + Sujet + 想 + 先 + Verbe + Objet + 在 + Lieu].
\item \textbf{Chinois} : \zh{因为农村工资低，所以他想先在城市找房子。}
\item \textbf{Pinyin} : Yīnwèi nóngcūn gōngzī dī, suǒyǐ tā xiǎng xiān zài chéngshì zhǎo fángzi.
\item \textbf{Vérification} : \zh{在城市} (lieu) placé \emph{avant} le verbe \zh{找} (règle : Lieu avant Verbe d'action). \zh{先} avant \zh{找}. Correct.
\end{enumerate}
\end{exoresolu}

\begin{exoresolu}[Version guidée : phrase unique]
Énoncé : Traduis en français : \zh{为了工作，她去年搬到了广州，虽然房租贵，但是很方便。}
\tcblower
Solution détaillée :
\begin{enumerate}
\item \textbf{Segmentation} : \zh{为了工作} (But) / \zh{她去年搬到了广州} (Fait passé) / \zh{虽然房租贵} (Concession) / \zh{但是很方便} (Contraste).
\item \textbf{Marqueurs} : \zh{去年} (passé) + \zh{了} (accompli) → Passé Composé. \zh{虽然…但是…} (bien que… mais…).
\item \textbf{Littéral} : « Pour travail, elle année-passée déménager-vers Guangzhou, bien-que loyer cher, mais très pratique. »
\item \textbf{Français naturel} : « Pour son travail, elle a déménagé à Guangzhou l'année dernière ; bien que le loyer soit cher, c'est très pratique. »
\item \textbf{Note} : \zh{虽然…但是…} structure fixe. En français, « bien que » appelle le subjonctif (\emph{soit}).
\end{enumerate}
\end{exoresolu}

\begin{attention}[Erreurs fréquentes des francophones — Thème / Version]
\begin{itemize}
\item \textbf{Double marque en français} : « Parce qu'il n'y avait pas de travail, \emph{donc} il est parti » est lourd en français mais correct en chinois (\zh{因为…所以…}). Traduire par un seul connecteur français.
\item \textbf{Place de 也} : ne dis jamais \zh{我想去城市也} ; \zh{也} se met avant le verbe : \zh{我也想去城市。}
\item \textbf{再 contre 又} : \zh{再} sert pour « ensuite / une prochaine fois » (projet, futur) ; \zh{又} pour une répétition déjà accomplie. Dans un projet, c'est toujours \zh{先…再…}.
\item \textbf{Ordre Lieu / Verbe} : Complément de lieu \zh{在 + Lieu} se place \emph{avant} le verbe d'action (aller, habiter, mettre, chercher) : \zh{在城市找房子} (chercher une maison \emph{dans} la ville), jamais \zh{找房子在城市}.
\item \textbf{Particule 了 position} : Verbe + \zh{了} (accompli) vs Phrase + \zh{了} (changement d'état). \zh{搬到了} (Verbe + Complément de direction + \zh{了}) = action accomplie.
\item \textbf{Caractères proches} : \zh{租} (louer, radical \zh{禾}) vs \zh{组} (groupe, radical \zh{糸}) ; \zh{拥} (bondé, radical \zh{扌}) vs \zh{用} (utiliser, radical \zh{用}) ; \zh{挤} (pousser, radical \zh{扌}) vs \zh{齐} (aligné).
\end{itemize}
\end{attention}

\courssub{Méthode : rédiger un texte simple sur l'exode rural (Production écrite)}

\begin{methode}[Écrire ton propre texte en 6 étapes (modèle du texte commenté)]
\begin{enumerate}
\item \textbf{Personnage} : Qui part ? (\zh{我的表弟} — mon cousin cadet, \zh{我的同学} — mon camarade, \zh{我} — moi). Ajoute \zh{以前住在 + [Village/Ville]}.
\item \textbf{Cause} : \zh{因为 + [Lieu] + 工作少 / 工资低}，\zh{所以…} (utilise \zh{搬到} + [Ville : 雅温得, 杜阿拉, 北京, 上海] + \zh{了}).
\item \textbf{Contraste ville} : \zh{城市的生活很方便，但是房租很贵，人也很拥挤。} (Réutilise ce bloc, change adjectifs si besoin : \zh{累} fatigué, \zh{快} rapide).
\item \textbf{Projet} : \zh{他/她/我想先 + [Verbe 1 : 找工作, 赚钱], 再 + [Verbe 2 : 学中文, 买房子, 结婚]。}
\item \textbf{Motivation} : \zh{为了 + [家人, 父母, 孩子, 自己]，他/她/我每天都很努力。}
\item \textbf{Relecture finale} : Vérifie les 7 connecteurs, l'ordre Sujet-Temps-Lieu-Verbe, les caractères complexes (\zh{搬}, \zh{拥挤}, \zh{租}, \zh{雅温得}), la ponctuation chinoise.
\end{enumerate}
\end{methode}

\begin{experience}[Atelier d'écriture collaborative : « Mon histoire d'exode »]
\begin{enumerate}
\item \textbf{Individuel (10 min)} : Chaque élève rédige son texte de 6 phrases (suivant la méthode ci-dessus) sur brouillon. Utilise le lexique de la carte mentale.
\item \textbf{Binôme (10 min)} : Échange ton texte. Ton partenaire : (a) souligne les connecteurs, (b) entoure les 5 caractères complexes ciblés (\zh{搬}, \zh{拥}, \zh{挤}, \zh{租}, \zh{温} ou \zh{得}), (c) vérifie l'ordre des traits sur un dictionnaire ou l'appli, (d) propose une correction en pinyin.
\item \textbf{Groupe/Classe (10 min)} : 3 volontaires lisent leur texte au tableau. La classe valide : « Les connecteurs sont-ils tous là ? Les caractères sont-ils lisibles ? Le pinyin est-il juste ? ».
\item \textbf{Remise} : Version finale propre (caractères + pinyin + traduction française) à déposer sur OnBuch+.
\end{enumerate}
\end{experience}

\begin{savaistu}[Les 农民工, les travailleurs migrants chinois]
En Chine, les paysans partis travailler en ville s'appellent \zh{农民工} (\emph{nóngmíngōng}, littéralement « ouvriers-paysans »). Ils sont plus de 290 millions : c'est l'une des plus grandes migrations de l'histoire humaine. Beaucoup laissent leurs enfants au village avec les grands-parents et ne reviennent qu'au Nouvel An chinois. Ce phénomène est tout à fait comparable à l'exode rural vers Douala et Yaoundé : dans les deux pays, la ville attire par ses emplois, mais impose ses loyers chers et ses quartiers bondés.
\end{savaistu}

\begin{exercice}[Entraînement complet — Expression écrite et traduction]
\begin{enumerate}
\item \textbf{Écriture} : Recopie 3 fois chacun des 5 caractères complexes ci-dessous en respectant l'ordre des traits, puis écris leur pinyin et leur sens.
      \zh{搬} \quad \zh{拥} \quad \zh{挤} \quad \zh{租} \quad \zh{温}
\item \textbf{Version} : Traduis en français naturel.
      \begin{enumerate}
      \item \zh{因为城市很拥挤，所以他想先找房子。}
      \item \zh{为了家人，她每天都很努力学习中文。}
      \end{enumerate}
\item \textbf{Thème} : Traduis en chinois (caractères + pinyin).
      \begin{enumerate}
      \item « Mon cousin habitait à la campagne. Parce que le travail était peu, il a déménagé à Douala l'année dernière. »
      \item « La vie en ville est pratique, mais le loyer est cher. Il veut d'abord gagner de l'argent, puis acheter une maison. »
      \end{enumerate}
\item \textbf{Production} : Rédige un texte de 6 phrases (environ 80 caractères) sur le thème « Mon ami(e) a quitté le village pour la ville ». Suis la méthode en 6 étapes. Écris : Caractères — Pinyin — Traduction française.
\end{enumerate}
\end{exercice}

\begin{corrige}[Exercice : Entraînement complet]
\textbf{1. Écriture (modèle de tracé) :}
\begin{itemize}
\item \zh{搬} (bān) : radical \zh{扌}, 13 traits. Sens : déménager.
\item \zh{拥} (yōng) : radical \zh{扌}, 19 traits. Sens : embrasser / bondé (dans \zh{拥挤}).
\item \zh{挤} (jǐ) : radical \zh{扌}, 15 traits. Sens : serré, pousser.
\item \zh{租} (zū) : radical \zh{禾}, 13 traits. Sens : louer.
\item \zh{温} (wēn) : radical \zh{氵}, 12 traits. Sens : tiède / partie de Yaoundé.
\end{itemize}

\textbf{2. Version :}
\begin{enumerate}
\item « Parce que la ville est très bondée, il veut d'abord trouver une maison. » (ou « un logement »)
\item « Pour sa famille, elle étudie le chinois avec acharnement chaque jour. » (ou « fait des efforts pour étudier… »)
\end{enumerate}

\textbf{3. Thème :}
\begin{enumerate}
\item \zh{我的表哥以前住在农村。因为工作太少，所以他去年搬到了杜阿拉。}\\
      Wǒ de biǎogē yǐqián zhù zài nóngcūn. Yīnwèi gōngzuò tài shǎo, suǒyǐ tā qùnián bān dào le Dùālā.
\item \zh{城市的生活很方便，但是房租很贵。他想先赚钱，再买房子。}\\
      Chéngshì de shēnghuó hěn fāngbiàn, dànshì fángzū hěn guì. Tā xiǎng xiān zhuànqián, zài mǎi fángzi.
      \emph{Note :} \zh{赚钱} (zhuànqián) = gagner de l'argent. \zh{买} (mǎi) = acheter.
\end{enumerate}

\textbf{4. Production (exemple de réponse attendue) :}
\zh{我的同学小李以前住在村子。因为村子里工资太低，所以他前年搬到了雅温得。城市的生活很方便，但是房租很贵，人也很拥挤。他想先找一份好工作，再买一个房子。为了父母，他每天都很努力。}\\
Wǒ de tóngxué Xiǎo Lǐ yǐqián zhù zài cūnzi. Yīnwèi cūnzi lǐ gōngzī tài dī, suǒyǐ tā qiánnián bān dào le Yǎwēndé. Chéngshì de shēnghuó hěn fāngbiàn, dànshì fángzū hěn guì, rén yě hěn yōngjǐ. Tā xiǎng xiān zhǎo yī fèn hǎo gōngzuò, zài mǎi yī gè fángzi. Wèile fùmǔ, tā měitiān dōu hěn nǔlì.\\
\emph{Traduction :} « Mon camarade Xiao Li habitait autrefois au village. Parce que les salaires y étaient trop bas, il a déménagé à Yaoundé l'année d'avant. La vie en ville est pratique, mais le loyer est cher et c'est bondé. Il veut d'abord trouver un bon travail, puis acheter une maison. Pour ses parents, il fait des efforts chaque jour. »
\end{corrige}

\coursec{Lexique écrit du thème}

Dans la section précédente, nous avons lu et commenté un texte modèle sur l'exode rural : un jeune homme quitte son village pour s'installer en ville. Ce texte reposait sur un petit noyau de mots essentiels : \zh{农村}, \zh{城市}, \zh{离开}, \zh{搬}, \zh{工作}… Pour pouvoir, à votre tour, \emph{écrire} un texte simple sur ce thème, il ne suffit pas de reconnaître ces mots : il faut les \textbf{maîtriser en production}, c'est-à-dire connaître leur nature grammaticale, leurs combinaisons naturelles (les \cle{collocations}) et les constructions de phrases dans lesquelles ils s'insèrent. C'est l'objet de cette section : bâtir votre \cle{lexique écrit actif} du thème de l'exode rural.

\courssub{Observation : les mots du thème en contexte}

Avant tout tableau, observons ces mots dans des phrases réelles. Lisez chaque exemple, repérez le mot-clé et devinez son rôle dans la phrase : est-il un nom, un verbe, un adjectif ?

\begin{exemplebox}[Les douze mots du thème en situation]
\zh{越来越多的年轻人离开农村。}\\
\emph{Yuè lái yuè duō de niánqīngrén líkāi nóngcūn.}\\
« De plus en plus de jeunes quittent la campagne. »

\zh{他搬到杜阿拉找工作。}\\
\emph{Tā bān dào Dù'ālā zhǎo gōngzuò.}\\
« Il déménage à Douala pour chercher du travail. »

\zh{城市里的工作机会比较多。}\\
\emph{Chéngshì lǐ de gōngzuò jīhuì bǐjiào duō.}\\
« En ville, les opportunités de travail sont plus nombreuses. »

\zh{杜阿拉很拥挤，房租也很贵。}\\
\emph{Dù'ālā hěn yōngjǐ, fángzū yě hěn guì.}\\
« Douala est très bondée, et le loyer est aussi très cher. »

\zh{城市的生活很方便，但是不太安静。}\\
\emph{Chéngshì de shēnghuó hěn fāngbiàn, dànshì bú tài ānjìng.}\\
« La vie en ville est pratique, mais pas très calme. »

\zh{他希望努力学习，将来帮助家乡。}\\
\emph{Tā xīwàng nǔlì xuéxí, jiānglái bāngzhù jiāxiāng.}\\
« Il espère travailler dur pour aider son village natal plus tard. »
\end{exemplebox}

\textbf{Questions d'observation.} Dans la phrase \zh{他搬到杜阿拉找工作}, quel mot suit immédiatement \zh{搬} ? Et dans \zh{离开农村}, y a-t-il une préposition entre \zh{离开} et \zh{农村} ? Gardez ces deux remarques en mémoire : elles seront au cœur des règles de cette section.

\courssub{Glossaire du thème : les douze mots-clés}

Voici le glossaire complet du thème. Chaque mot est donné avec sa \textbf{nature grammaticale} : en chinois, connaître la nature d'un mot est indispensable, car elle détermine sa place dans la phrase.

\begin{center}
\begin{tabularx}{\linewidth}{|X|X|X|X|}
\hline
\rowcolor{popblueL} Caractères & Pinyin & Nature & Traduction \\
\hline
农村 & nóngcūn & nom & campagne, milieu rural \\
\hline
城市 & chéngshì & nom & ville \\
\hline
离开 & líkāi & verbe & quitter \\
\hline
搬 & bān & verbe & déménager, transporter \\
\hline
工作 & gōngzuò & nom / verbe & travail / travailler \\
\hline
机会 & jīhuì & nom & opportunité, occasion \\
\hline
生活 & shēnghuó & nom / verbe & vie / vivre \\
\hline
拥挤 & yōngjǐ & adjectif & bondé, surpeuplé \\
\hline
房租 & fángzū & nom & loyer \\
\hline
方便 & fāngbiàn & adjectif & pratique, commode \\
\hline
努力 & nǔlì & adjectif / verbe & appliqué / faire des efforts \\
\hline
希望 & xīwàng & verbe / nom & espérer / espoir \\
\hline
\end{tabularx}
\end{center}

\begin{popfigure}
\centering
\begin{tikzpicture}[mindmap, grow cyclic, every node/.style={concept, font=\small}, concept color=popblue!40, level 1/.append style={level distance=3.2cm, sibling angle=90}, level 2/.append style={level distance=2.4cm, sibling angle=45}]
\node[concept color=poporange!50, font=\small\bfseries] {Exode rural 离开农村}
  child[concept color=popgreen!40] { node {Lieux} child { node {农村 nóngcūn} } child { node {城市 chéngshì} } }
  child[concept color=poppurple!40] { node {Actions} child { node {离开 líkāi} } child { node {搬 bān} } child { node {找工作 zhǎo gōngzuò} } }
  child[concept color=popgold!50] { node {Vie en ville} child { node {拥挤 yōngjǐ} } child { node {房租 fángzū} } child { node {方便 fāngbiàn} } }
  child[concept color=poppink!40] { node {Avenir} child { node {机会 jīhuì} } child { node {努力 nǔlì} } child { node {希望 xīwàng} } };
\end{tikzpicture}
\legende{Carte mentale du lexique de l'exode rural : quatre familles de mots autour du thème central.}
\end{popfigure}

\begin{attention}[Mots à double nature : 工作, 生活, 努力, 希望]
Quatre mots de ce glossaire peuvent être à la fois \textbf{nom et verbe} (ou adjectif et verbe). Le français fait de même (« travail » / « travailler »), mais en chinois c'est \emph{le même mot, sans aucun changement de forme}. Seule la place dans la phrase indique la fonction :
\begin{itemize}
\item \zh{我在杜阿拉工作。} \emph{Wǒ zài Dù'ālā gōngzuò.} « Je travaille à Douala. » (\zh{工作} = verbe)
\item \zh{这个工作很好。} \emph{Zhège gōngzuò hěn hǎo.} « Ce travail est bien. » (\zh{工作} = nom)
\item \zh{他希望搬到城市。} \emph{Tā xīwàng bān dào chéngshì.} « Il espère déménager en ville. » (\zh{希望} = verbe, suivi d'une proposition)
\item \zh{年轻人有希望。} \emph{Niánqīngrén yǒu xīwàng.} « Les jeunes ont de l'espoir. » (\zh{希望} = nom)
\end{itemize}
\end{attention}

\courssub{Étude détaillée : 搬 et 离开, deux verbes à construire avec soin}

\begin{definition}[搬 bān]
\zh{搬} signifie « déménager, transporter ». En production écrite, il est \textbf{toujours suivi d'un complément de direction} : la structure est \cle{搬到 + lieu} (« déménager à/en + lieu »). Le caractère se compose de la clé de la main \zh{扌} à gauche (le sens : action de la main) et de \zh{般} à droite (le son).
\end{definition}

\begin{propriete}[Structure : Sujet + 搬到 + lieu (+ 找工作 / 生活)]
\zh{搬到} fonctionne comme un bloc inséparable : \zh{到} marque l'arrivée, la destination. Ne dites jamais \zh{搬杜阿拉} sans \zh{到}. Pour marquer une action accomplie, on ajoute \zh{了} en fin de phrase : \zh{他搬到杜阿拉了。} « Il a déménagé à Douala. »
\begin{itemize}
\item \zh{我们一家搬到雅温得。} \emph{Wǒmen yì jiā bān dào Yǎwēndé.} « Notre famille déménage à Yaoundé. »
\item \zh{很多农民搬到城市生活。} \emph{Hěn duō nóngmín bān dào chéngshì shēnghuó.} « Beaucoup de paysans déménagent en ville pour y vivre. »
\end{itemize}
\end{propriete}

\begin{attention}[离开 ne prend pas de 从]
En français on dit « partir \emph{de} la campagne », et le chinois possède bien la préposition \zh{从} (cóng, « de, depuis »). Mais le verbe \zh{离开} contient déjà l'idée de séparation : il se construit \textbf{directement avec son complément}, sans préposition.
\begin{itemize}
\item Correct : \zh{他离开了农村。} \emph{Tā líkāi le nóngcūn.} « Il a quitté la campagne. »
\item À éviter en production simple : \zh{他从农村离开。} (construction lourde et redondante à ce niveau)
\end{itemize}
Structure à retenir : \cle{Sujet + 离开 + lieu/personne}. On peut aussi quitter une personne : \zh{他离开了家人。} \emph{Tā líkāi le jiārén.} « Il a quitté sa famille. »
\end{attention}

\courssub{Les collocations utiles : apprendre les mots en blocs}

Une \cle{collocation} est une combinaison de mots qui vont naturellement ensemble. En chinois comme en français, on ne combine pas les mots au hasard : on dit « chercher du travail » (\zh{找工作}) tant qu'on est en recherche, et « trouver du travail » (\zh{找到工作}) seulement quand on l'a obtenu. Voici les quatre collocations essentielles du thème.

\begin{center}
\begin{tabularx}{\linewidth}{|X|X|X|}
\hline
\rowcolor{popblueL} Collocation & Pinyin & Traduction \\
\hline
离开家乡 & líkāi jiāxiāng & quitter son pays natal \\
\hline
搬到城市 & bān dào chéngshì & déménager en ville \\
\hline
找工作 & zhǎo gōngzuò & chercher du travail \\
\hline
生活很方便 & shēnghuó hěn fāngbiàn & la vie est pratique \\
\hline
\end{tabularx}
\end{center}

\begin{methode}[Apprendre par collocations, pas par mots isolés]
Pour écrire des phrases correctes sans hésiter, procédez ainsi :
\begin{enumerate}
\item Apprenez chaque mot \textbf{dans son bloc} : ne mémorisez pas « \zh{搬} = déménager », mémorisez « \zh{搬到城市} = déménager en ville ».
\item Remplacez un élément du bloc pour créer une phrase personnelle : \zh{搬到城市} devient \zh{搬到杜阿拉}, \zh{搬到雅温得}.
\item Ajoutez un sujet et un complément de but : \zh{他 + 搬到杜阿拉 + 找工作。}
\item Relisez en vérifiant les deux pièges : \zh{搬} a-t-il son \zh{到} ? \zh{离开} est-il sans \zh{从} ?
\end{enumerate}
\end{methode}

\textbf{Phrases types construites sur chaque collocation :}

\begin{exemplebox}[Une collocation, une phrase]
\zh{很多年轻人离开家乡，去城市工作。}\\
\emph{Hěn duō niánqīngrén líkāi jiāxiāng, qù chéngshì gōngzuò.}\\
« Beaucoup de jeunes quittent leur pays natal pour aller travailler en ville. »

\zh{去年，我哥哥搬到城市了。}\\
\emph{Qùnián, wǒ gēge bān dào chéngshì le.}\\
« L'année dernière, mon grand frère a déménagé en ville. »

\zh{他在杜阿拉找工作，找了三个月。}\\
\emph{Tā zài Dù'ālā zhǎo gōngzuò, zhǎo le sān ge yuè.}\\
« Il a cherché du travail à Douala pendant trois mois. »

\zh{城市生活很方便，但是房租很贵。}\\
\emph{Chéngshì shēnghuó hěn fāngbiàn, dànshì fángzū hěn guì.}\\
« La vie en ville est pratique, mais le loyer est très cher. »
\end{exemplebox}

\courssub{Expressions de quantité et de degré}

Pour décrire l'exode rural, il faut pouvoir dire que le phénomène \emph{augmente} et nuancer les qualités de la vie en ville. Trois outils suffisent à ce niveau.

\begin{propriete}[越来越多 + (的) + nom : « de plus en plus de … »]
Structure : \cle{越来越多 + 的 + nom + verbe}. Le \zh{的} est facultatif mais recommandé à l'écrit.
\begin{itemize}
\item \zh{越来越多的人离开农村。} \emph{Yuè lái yuè duō de rén líkāi nóngcūn.} « De plus en plus de gens quittent la campagne. »
\item \zh{越来越多的学生学汉语。} \emph{Yuè lái yuè duō de xuésheng xué Hànyǔ.} « De plus en plus d'élèves apprennent le chinois. »
\end{itemize}
Attention à la prononciation : \zh{越来越} se lit \emph{yuè lái yuè} — le \zh{越} répété exprime la progression, comme « more and more » en anglais.
\end{propriete}

\begin{propriete}[Marques de degré : 很 / 太 / 不太 + adjectif]
Devant un adjectif qualificatif (\zh{拥挤}, \zh{方便}, \zh{贵}), le chinois exige presque toujours un marqueur de degré :
\begin{itemize}
\item \zh{很} hěn : marqueur neutre (souvent simple liaison, pas forcément « très ») : \zh{城市很拥挤。} \emph{Chéngshì hěn yōngjǐ.} « La ville est bondée. »
\item \zh{太 … 了} tài … le : « trop … » : \zh{房租太贵了。} \emph{Fángzū tài guì le.} « Le loyer est trop cher. »
\item \zh{太少} tài shǎo : « trop peu » : \zh{工作机会太少了。} \emph{Gōngzuò jīhuì tài shǎo le.} « Les opportunités de travail sont trop rares. »
\item \zh{不太} bú tài : « pas très » : \zh{农村的生活不太方便。} \emph{Nóngcūn de shēnghuó bú tài fāngbiàn.} « La vie à la campagne n'est pas très pratique. »
\end{itemize}
\end{propriete}

\begin{attention}[L'adjectif seul ne suffit pas]
En français, « la ville est bondée » suffit. En chinois, une phrase affirmative avec un adjectif seul (\zh{城市拥挤}) sonne comme une comparaison inachevée (« la ville, elle, est bondée… contrairement à… »). En production simple, mettez toujours \zh{很}, \zh{太} ou \zh{不太} devant l'adjectif : \zh{城市很拥挤。}
\end{attention}

\courssub{Tableau récapitulatif : constructions à maîtriser}

\begin{center}
\begin{tabularx}{\linewidth}{|X|X|X|}
\hline
\rowcolor{popblueL} Structure & Exemple (caractères + pinyin) & Traduction \\
\hline
Sujet + 离开 + lieu & \zh{他离开农村。} Tā líkāi nóngcūn. & Il quitte la campagne. \\
\hline
Sujet + 搬到 + lieu & \zh{他搬到杜阿拉。} Tā bān dào Dù'ālā. & Il déménage à Douala. \\
\hline
(zài + lieu) + 找工作 & \zh{他在城市找工作。} Tā zài chéngshì zhǎo gōngzuò. & Il cherche du travail en ville. \\
\hline
越来越多 + 的 + nom + V & \zh{越来越多的人搬到城市。} Yuè lái yuè duō de rén bān dào chéngshì. & De plus en plus de gens déménagent en ville. \\
\hline
Nom + 很/太/不太 + adj. & \zh{房租太贵了。} Fángzū tài guì le. & Le loyer est trop cher. \\
\hline
Sujet + 希望 + proposition & \zh{他希望找到工作。} Tā xīwàng zhǎodào gōngzuò. & Il espère trouver du travail. \\
\hline
\end{tabularx}
\end{center}

\courssub{Exercice résolu et entraînement}

\begin{exoresolu}[Traduisez en chinois]
Énoncé : Traduisez les phrases suivantes en utilisant les collocations du thème. 1) De plus en plus de jeunes quittent leur village natal. 2) Elle a déménagé à Yaoundé pour chercher du travail. 3) La vie en ville est pratique, mais le loyer est trop cher.
\tcblower
Solution détaillée :
\begin{enumerate}
\item « De plus en plus de jeunes » = \zh{越来越多的年轻人} ; « quitter » = \zh{离开} (sans \zh{从} !) ; « leur village natal » = \zh{家乡}. Résultat : \zh{越来越多的年轻人离开家乡。} \emph{Yuè lái yuè duō de niánqīngrén líkāi jiāxiāng.}
\item « Déménager à Yaoundé » = \zh{搬到雅温得} (avec \zh{到}) ; « pour chercher du travail » = \zh{找工作} placé après ; l'action est accomplie, on ajoute \zh{了} en fin de phrase. Résultat : \zh{她搬到雅温得找工作了。} \emph{Tā bān dào Yǎwēndé zhǎo gōngzuò le.}
\item « La vie en ville » = \zh{城市的生活} ; « pratique » = \zh{方便} précédé de \zh{很} ; « mais » = \zh{但是} ; « trop cher » = \zh{太贵了}. Résultat : \zh{城市的生活很方便，但是房租太贵了。} \emph{Chéngshì de shēnghuó hěn fāngbiàn, dànshì fángzū tài guì le.}
\end{enumerate}
\end{exoresolu}

\begin{exercice}[À vous de jouer]
1) Traduisez en chinois : « Beaucoup de paysans déménagent en ville parce qu'il y a plus d'opportunités. » (indice : « parce que » = \zh{因为} yīnwèi ; « il y a » = \zh{有} yǒu ; « plus » = \zh{更多} gèng duō).\\
2) Complétez avec le bon mot (\zh{离开}, \zh{搬到}, \zh{机会}, \zh{拥挤}) : \zh{杜阿拉很……，房租也很贵。} / \zh{他……农村，去城市找工作。}\\
3) Corrigez la phrase fautive : \zh{他从家乡离开。}
\end{exercice}

\begin{corrige}[Exercice : lexique du thème]
1) \zh{很多农民搬到城市，因为城市有更多的机会。} \emph{Hěn duō nóngmín bān dào chéngshì, yīnwèi chéngshì yǒu gèng duō de jīhuì.} « Beaucoup de paysans déménagent en ville parce qu'en ville il y a plus d'opportunités. »\\
2) \zh{杜阿拉很拥挤，房租也很贵。} ; \zh{他离开农村，去城市找工作。}\\
3) \zh{他离开了家乡。} — \zh{离开} se construit directement avec son complément, sans \zh{从}.
\end{corrige}

\begin{aretenir}[L'essentiel du lexique de l'exode rural]
\begin{itemize}
\item Douze mots-clés, dont quatre à double nature : \zh{工作}, \zh{生活}, \zh{努力}, \zh{希望}.
\item Deux constructions à ne jamais fausser : \cle{搬到 + lieu} (avec \zh{到}) et \cle{离开 + lieu} (sans \zh{从}).
\item Quatre collocations à réciter par cœur : \zh{离开家乡}, \zh{搬到城市}, \zh{找工作}, \zh{生活很方便}.
\item Trois marques de degré : \zh{很}, \zh{太…了}, \zh{不太} ; et l'expression de progression \zh{越来越多}.
\end{itemize}
\end{aretenir}

Ce lexique est désormais votre boîte à outils. Dans la section suivante, nous passerons de la phrase au caractère : nous apprendrons à \emph{écrire} correctement les caractères complexes de ce thème — traits, radicaux et pièges des caractères proches.

\coursec{Écriture des caractères complexes}

Dans la section précédente, nous avons constitué le lexique écrit du thème de l'exode rural : \zh{农村} (campagne), \zh{城市} (ville), \zh{搬家} (déménager), \zh{离开} (quitter), \zh{拥挤} (bondé), \zh{租房} (louer un logement), \zh{愿意} (vouloir, être disposé à). Savoir reconnaître ces mots ne suffit pas : en expression écrite, l'examinateur attend que vous sachiez les \cle{tracer correctement}, trait par trait. C'est l'objet de cette section : maîtriser l'écriture des caractères complexes du thème, comprendre leur architecture (radical + composant phonétique) et éviter les confusions avec des caractères voisins.

\courssub{Observation : comment est construit un caractère complexe ?}

Observons d'abord une phrase extraite de notre modèle de texte :

\zh{很多年轻人离开农村，搬到拥挤的城市里租房子住。}

\emph{Hěn duō niánqīngrén líkāi nóngcūn, bān dào yōngjǐ de chéngshì lǐ zū fángzi zhù.}

« Beaucoup de jeunes quittent la campagne, déménagent dans des villes bondées et y louent un logement. »

Regardez les caractères \zh{搬}, \zh{离}, \zh{挤}, \zh{租} : aucun n'est un dessin unique et indivisible. Chacun est \cle{assemblé} à partir de composants plus simples, écrits dans un ordre précis. C'est la grande différence avec le français : en français, on écrit des lettres de gauche à droite sur une ligne ; en chinois, on \emph{construit} chaque caractère dans un carré imaginaire, en suivant des règles de séquencement.

\begin{propriete}[Règle d'écriture : la gauche avant la droite]
Dans un caractère composé de deux parties horizontales, le composant de \cle{gauche} — qui est très souvent le \cle{radical} (la clé de sens) — s'écrit \textbf{avant} la partie de droite. Ainsi, pour \zh{搬} : on écrit d'abord \zh{扌} (la main, à gauche), \textbf{puis} \zh{般} (à droite). De même : \zh{挤} = \zh{扌} puis \zh{齐} ; \zh{租} = \zh{禾} puis \zh{且} ; \zh{村} = \zh{木} puis \zh{寸}.
\end{propriete}

\begin{aretenir}[Les quatre règles générales d'ordre des traits]
\begin{enumerate}
\item \cle{De haut en bas} : \zh{三} (le trait du haut d'abord), \zh{愿} (la partie haute \zh{原} avant \zh{心} en bas).
\item \cle{De gauche à droite} : \zh{搬}, \zh{挤}, \zh{租}, \zh{村}.
\item \cle{L'extérieur avant l'intérieur} : \zh{同}, \zh{问}.
\item \cle{La fermeture en dernier} : dans \zh{国}, le trait horizontal du bas qui « ferme » le cadre se trace en tout dernier ; dans \zh{离}, l'élément interne \zh{凶} se termine avant le trait de clôture final.
\end{enumerate}
\end{aretenir}

\courssub{Étude détaillée des cinq caractères complexes du thème}

\textbf{1. \zh{搬} bān — déménager, transporter (13 traits)}\par

\zh{搬} se décompose en \zh{扌} (radical de la \cle{main}, 3 traits) + \zh{般} (10 traits, qui donne le son \emph{bān}). Ordre : la main d'abord (横 \emph{héng}, 竖钩 \emph{shùgōu}, 提 \emph{tí}), puis \zh{般} de gauche à droite : d'abord \zh{舟} (la barque, 6 traits), ensuite \zh{殳} (4 traits). Le sens est logique : déménager, c'est un travail de \emph{mains}.

\zh{他们下个月搬家。} \emph{Tāmen xià ge yuè bānjiā.} « Ils déménagent le mois prochain. »

\begin{popfigure}
\begin{tikzpicture}[font=\large]
\node[draw=popblue, thick, minimum width=1.3cm, minimum height=1.3cm, fill=popblueL] (shou) at (0,0) {\zh{扌}};
\node[draw=poporange, thick, minimum width=1.3cm, minimum height=1.3cm, fill=poporangeL] (ban) at (2.2,0) {\zh{般}};
\node[draw=popdark, very thick, minimum width=1.3cm, minimum height=1.3cm] (res) at (5.2,0) {\zh{搬}};
\node[above=0.15cm of shou, align=center, text=popblue] {\small 1\ier{} : la main\\ \small 3 traits};
\node[above=0.15cm of ban, align=center, text=poporange] {\small 2\ieme{} : \emph{bān}\\ \small 10 traits, gauche $\rightarrow$ droite};
\node[above=0.15cm of res, align=center] {\small résultat\\ \small 13 traits};
\draw[-{Stealth}, very thick, popblue] (shou.east) -- (ban.west);
\draw[-{Stealth}, very thick, popdark] (ban.east) -- (res.west);
\node[below=0.15cm of res, align=center, text=popmuted] {\small bān : déménager};
\end{tikzpicture}
\legende{Décomposition de \zh{搬} : le radical \zh{扌} (main) s'écrit d'abord, puis le composant phonétique \zh{般}.}
\end{popfigure}

\textbf{2. \zh{离} lí — quitter, être éloigné de (10 traits)}\par

Radical : \zh{亠} (la « coiffe », 2 traits). On écrit \cle{de haut en bas} : d'abord le point initial et le trait horizontal du haut (\zh{亠}), puis la partie interne \zh{凶} (4 traits), et enfin le trait de clôture horizontal du bas (1 trait) et le trait vertical descendant final (1 trait) — soit 2+4+2+? Attendez, standard simplifié : 亠(2) + 凶(4) + 两横一竖? Non. Ordre standard : 1.丶 2.一 (亠) 3.一 4.丨 5.一 6.乙 (凶) 7.一 8.丨 9.一 10.一 ? Vérification : \zh{离} 10 traits. Ordre : 亠 (2), puis 凶 (4), puis le bas (4 traits : horizontal, vertical, horizontal, horizontal). La règle « fermeture en dernier » s'applique au trait horizontal le plus bas.
\emph{Correction pédagogique :} On écrit \zh{亠}, puis \zh{凶}, puis les traits du bas (horizontal, vertical, horizontal, horizontal). Le dernier trait horizontal « ferme » le caractère.

\zh{他十八岁离开了村子。} \emph{Tā shíbā suì líkāi le cūnzi.} « Il a quitté le village à dix-huit ans. »

\textbf{3. \zh{挤} jǐ — presser, être bondé (9 traits)}\par

Radical \zh{扌} (main) à gauche, écrit en \cle{3 traits} : 横 (horizontal), 竖钩 (vertical-crochet), 提 (montant). Puis \zh{齐} à droite (6 traits : 米 + 丿 + 丶). On le retrouve dans \zh{拥挤} (bondé, congestionné), mot clé pour décrire les grandes villes comme Douala ou Pékin aux heures de pointe.

\zh{早高峰的时候，公交车非常拥挤。} \emph{Zǎo gāofēng de shíhou, gōngjiāochē fēicháng yōngjǐ.} « Aux heures de pointe du matin, les bus sont très bondés. »

\textbf{4. \zh{租} zū — louer (10 traits)}\par

Radical \zh{禾} (la \cle{céréale}, le riz sur pied, 5 traits) + \zh{且} (qui suggère le son, 5 traits). Point délicat : dans \zh{禾}, le dernier trait, qui serait normalement un 捺 (nà, trait tombant vers la droite), \cle{devient un point} 点 (diǎn) pour laisser la place au composant de droite. Historiquement, le loyer se payait en grain : d'où la clé de la céréale.

\zh{她在城里租了一个小房间。} \emph{Tā zài chéng lǐ zū le yí ge xiǎo fángjiān.} « Elle a loué une petite chambre en ville. »

\textbf{5. \zh{愿} yuàn — vouloir, être disposé à (14 traits)}\par

Radical : \zh{心} (le \cle{cœur}), placé \cle{en bas} (4 traits). On écrit d'abord \zh{原} (la partie haute, 10 traits : \zh{厂} 2, \zh{白} 5, \zh{小} 3), \textbf{puis} \zh{心} en dessous. Le cœur en bas exprime le « vouloir », le désir profond. On le trouve dans \zh{愿意} (consentir à, 4+4) et \zh{愿望} (souhait, 4+4).

\zh{很多年轻人愿意离开农村去城市工作。} \emph{Hěn duō niánqīngrén yuànyì líkāi nóngcūn qù chéngshì gōngzuò.} « Beaucoup de jeunes acceptent de quitter la campagne pour travailler en ville. »

\begin{methode}[Comment apprendre un caractère complexe]
\begin{enumerate}
\item \cle{Observer} : identifier le radical (sens) et le composant phonétique (son).
\item \cle{Compter} : écrire le caractère 5 fois en comptant les traits \textbf{à voix haute} (一、二、三…).
\item \cle{Décomposer} : réécrire le caractère en nommant chaque composant (« main, puis 般 »).
\item \cle{Mémoriser} : refermer le cahier et réécrire le caractère de mémoire, puis vérifier.
\item \cle{Réutiliser} : écrire une phrase complète avec le caractère, par exemple sur l'exode rural au Cameroun (Yaoundé, Douala, Buea).
\end{enumerate}
\end{methode}

\courssub{Caractères proches : ne pas se tromper}

Certaines paires de caractères se ressemblent beaucoup mais n'ont ni le même sens ni parfois le même son. La seule protection : vérifier le \cle{radical} et le \cle{sens} avant d'écrire.

\begin{popfigure}
\begin{tikzpicture}
\matrix[column sep=0.35cm, row sep=0.25cm, nodes={align=center, minimum height=0.85cm, anchor=center}] {
\node[fill=popblueL, minimum width=1.1cm, font=\large] {\zh{村}}; &
\node[fill=popblueL, minimum width=2.1cm] {\small cūn : village}; &
\node[fill=poporangeL, minimum width=1.1cm, font=\large] {\zh{材}}; &
\node[fill=poporangeL, minimum width=2.1cm] {\small cái : matériau}; \\
\node[fill=popgreenL, minimum width=1.1cm, font=\large] {\zh{搬}}; &
\node[fill=popgreenL, minimum width=2.1cm] {\small bān : déménager}; &
\node[fill=poppinkL, minimum width=1.1cm, font=\large] {\zh{般}}; &
\node[fill=poppinkL, minimum width=2.1cm] {\small bān : sorte (一般)}; \\
\node[fill=poppurpleL, minimum width=1.1cm, font=\large] {\zh{挤}}; &
\node[fill=poppurpleL, minimum width=2.1cm] {\small jǐ : presser}; &
\node[fill=popgoldL, minimum width=1.1cm, font=\large] {\zh{济}}; &
\node[fill=popgoldL, minimum width=2.1cm] {\small jì : économie (经济)}; \\
\node[fill=popblueL, minimum width=1.1cm, font=\large] {\zh{租}}; &
\node[fill=popblueL, minimum width=2.1cm] {\small zū : louer}; &
\node[fill=popgreenL, minimum width=1.1cm, font=\large] {\zh{组}}; &
\node[fill=popgreenL, minimum width=2.1cm] {\small zǔ : groupe}; \\
};
\end{tikzpicture}
\legende{Quatre paires de caractères proches à distinguer par le radical et le sens.}
\end{popfigure}

\begin{center}
\begin{tabularx}{\linewidth}{|X|X|X|X|}
\hline
\rowcolor{popblueL} Caractères & Pinyin & Nature & Traduction \\
\hline
村 & cūn & n. & village — \zh{农村} campagne \\
\hline
材 & cái & n. & matériau, bois — \zh{材料} matériel \\
\hline
搬 & bān & v. & déménager, transporter — \zh{搬家} déménager \\
\hline
般 & bān & n. & sorte, genre — \zh{一般} en général \\
\hline
挤 & jǐ & v. & presser, serrer — \zh{拥挤} bondé \\
\hline
济 & jì & v./n. & secourir, économie — \zh{经济} économie \\
\hline
租 & zū & v. & louer (payer) — \zh{租房} louer un logement \\
\hline
组 & zǔ & v./n. & former un groupe — \zh{小组} petit groupe \\
\hline
\end{tabularx}
\end{center}

\begin{attention}[Erreur fréquente des francophones]
Confondre \zh{村} (village) et \zh{材} (matériau) : les deux contiennent \zh{木} et \zh{寸}/\zh{才}, et les francophones, habitués à lire vite de gauche à droite, échangent les parties droites. Écrire \zh{农材} au lieu de \zh{农村} rend la phrase incompréhensible. \textbf{Réflexe} : avant d'écrire, demandez-vous « quel est le sens ? ». Village = \zh{村} (le bois \zh{木} + le pouce \zh{寸}) ; matériau = \zh{材} (le bois + le talent \zh{才}). De même, ne confondez pas \zh{租} (louer, clé de la céréale \zh{禾}) et \zh{组} (groupe, clé de la soie \zh{纟}) : \zh{组房子} n'existe pas, on dit \zh{租房子}. Attention aussi à \zh{挤} (jǐ, 3\ieme{} ton, main) vs \zh{济} (jì, 4\ieme{} ton, eau) : \zh{拥济} est faux.
\end{attention}

\begin{savaistu}[Le radical de la main \zh{扌}]
Le radical \zh{扌}, forme abrégée de \zh{手} (main), apparaît dans de très nombreux verbes d'action manuelle : \zh{搬} (déménager), \zh{打} (frapper), \zh{找} (chercher), \zh{挤} (presser), \zh{拉} (tirer), \zh{推} (pousser), \zh{拍} (claquer), \zh{扔} (jeter). Quand vous voyez \zh{扌} à gauche d'un caractère inconnu, vous pouvez deviner qu'il s'agit probablement d'une action faite avec les mains. C'est un formidable outil de lecture pour deviner le sens !
\end{savaistu}

\begin{exoresolu}[Identifier et corriger l'erreur]
Énoncé. Un élève de Première A a écrit : \zh{很多年轻人离开农材，到城市里组房子。} Relevez les deux erreurs de caractères, corrigez-les et traduisez la phrase corrigée.
\tcblower
Solution détaillée.
\begin{enumerate}
\item \zh{农材} est incorrect : \zh{材} (cái) signifie « matériau ». Le mot « campagne » s'écrit \zh{农村}, avec \zh{村} (cūn, village, radical \zh{木} + \zh{寸}).
\item \zh{组房子} est incorrect : \zh{组} (zǔ) signifie « grouper, organiser ». « Louer un logement » s'écrit \zh{租房子}, avec \zh{租} (zū, radical \zh{禾} céréale).
\end{enumerate}
Phrase corrigée : \zh{很多年轻人离开农村，到城市里租房子。}

\emph{Hěn duō niánqīngrén líkāi nóngcūn, dào chéngshì lǐ zū fángzi.}

Traduction : « Beaucoup de jeunes quittent la campagne et vont louer un logement en ville. »
\end{exoresolu}

\begin{exercice}[À vous de tracer]
\begin{enumerate}
\item Écrivez 5 fois chacun des caractères \zh{搬}, \zh{离}, \zh{挤}, \zh{租}, \zh{愿} en comptant les traits à voix haute, puis recopiez-les de mémoire sans modèle.
\item Complétez avec le bon caractère (\zh{村}/\zh{材}, \zh{租}/\zh{组}, \zh{挤}/\zh{济}) : \zh{农}＿＿ (campagne) ; \zh{＿＿房子} (louer un logement) ; \zh{拥}＿＿ (bondé) ; \zh{经}＿＿ (économie).
\item Traduisez en français : \zh{他不愿意离开农村。}
\item \textbf{Thème guidé} : Traduisez en chinois (caractères + pinyin) : « Ma sœur aînée ne veut pas déménager à Douala. » (Indications : sœur aînée = \zh{姐姐} jiějie ; Douala = \zh{杜阿拉} Dù'ālā ; déménager à = \zh{搬到} bān dào).
\end{enumerate}
\end{exercice}

\begin{corrige}[Exercice « À vous de tracer »]
1. Vérifiez le nombre de traits et l'ordre : \zh{搬} = 13 (扌+般), \zh{离} = 10 (亠+凶+bas), \zh{挤} = 9 (扌+齐), \zh{租} = 10 (禾+且, dernier trait de 禾 = point), \zh{愿} = 14 (原+心). Le radical de gauche (\zh{扌}, \zh{禾}) ou du haut (\zh{亠}) se trace toujours en premier ; pour \zh{愿}, \zh{原} avant \zh{心}.
2. \zh{农村} ; \zh{租房子} ; \zh{拥挤} ; \zh{经济}.
3. \zh{他不愿意离开农村。} \emph{Tā bù yuànyì líkāi nóngcūn.} « Il ne veut pas quitter la campagne. »
4. \zh{姐姐不愿意搬到杜阿拉。} \emph{Jiějie bù yuànyì bān dào Dù'ālā.}
\end{corrige}

\begin{aretenir}[L'essentiel de la section]
\begin{itemize}
\item Un caractère complexe se \cle{construit} : radical d'abord (souvent à gauche ou en haut), composant phonétique ensuite.
\item Cinq caractères clés du thème : \zh{搬} bān (13 traits), \zh{离} lí (10), \zh{挤} jǐ (9), \zh{租} zū (10), \zh{愿} yuàn (14).
\item Règles d'ordre : haut $\rightarrow$ bas, gauche $\rightarrow$ droite, extérieur $\rightarrow$ intérieur, fermeture en dernier.
\item Toujours vérifier le \cle{sens} pour distinguer les caractères proches : \zh{村}/\zh{材}, \zh{搬}/\zh{般}, \zh{挤}/\zh{济}, \zh{租}/\zh{组}.
\end{itemize}
\end{aretenir}

\coursec{Modèle de texte commenté : « Le chemin de Ah Yong »}

\begin{textebox}[Modèle : 阿勇的路 \zh{Ā Yǒng de Lù} — Niveau HSK 3]
\zh{阿勇来自西部的一个山村。因为家里土地少，而且收成不好，全家靠种玉米生活很困难。为了给弟弟交学费，阿勇先和朋友借了路费，再坐了两天火车到了深圳。} \\
\emph{Ā Yǒng lái zì xībù de yī gè shāncūn. Yīnwèi jiāli tǔdì shǎo, érqiě shōuchéng bù hǎo, quánjiā kào zhòng yùmǐ shēnghuó hěn kùnnán. Wèile gěi dìdi jiāo xuéfèi, Ā Yǒng xiān hé péngyou jiè le lùfèi, zài zuò le liǎng tiān huǒchē dào le Shēnzhèn.} \\
« Ah Yong vient d'un village de montagne de l'Ouest. \textbf{Parce que} la famille a peu de terres, \textbf{de plus} les récoltes sont mauvaises, toute la famille vit difficilement en cultivant du maïs. \textbf{Pour} payer les frais de scolarité de son petit frère, Ah Yong \textbf{d'abord} a emprunté l'argent du voyage à un ami, \textbf{ensuite} a pris le train pendant deux jours pour arriver à Shenzhen. »

\zh{刚开始，他在工厂上夜班，工资不低。但是城市生活节奏太快，可是他不习惯吃食堂的饭，而且很想念妈妈做的酸辣汤。} \\
\emph{Gāng kāishǐ, tā zài gōngchǎng shàng yèbān, gōngzī bù dī. Dànshì chéngshì shēnghuó jiézòu tài kuài, kěshì tā bù xíguàn chī shítáng de fàn, érqiě hěn xiǎngniàn māma zuò de suānlàtāng.} \\
« Au début, il travaille en équipe de nuit à l'usine, le salaire n'est pas bas. \textbf{Mais} le rythme de vie urbain est trop rapide, \textbf{pourtant} il n'a pas l'habitude de manger à la cantine, \textbf{de plus} il regrette beaucoup la soupe aigre-épicée de sa mère. »

\zh{现在，阿勇已经当了组长。为了将来在老家盖新房，他先把钱存进银行，再计划考个驾照。他相信：因为自己努力，所以家里的日子会越来越好。} \\
\emph{Xiànzài, Ā Yǒng yǐjīng dāng le zǔzhǎng. Wèile jiānglái zài lǎojiā gài xīnfáng, tā xiān bǎ qián cún jìn yínháng, zài jìhuà kǎo gè jiàzhào. Tā xiāngxìn: yīnwèi zìjǐ nǔlì, suǒyǐ jiāli de rìzi huì yuè lái yuè hǎo.} \\
« Maintenant, Ah Yong est déjà devenu chef d'équipe. \textbf{Pour} l'avenir, construire une nouvelle maison au village natal, il \textbf{d'abord} met l'argent à la banque, \textbf{ensuite} prévoit de passer le permis de conduire. Il croit : \textbf{parce qu'}il s'efforce lui-même, \textbf{donc} la vie de la famille ira de mieux en mieux. »

\begin{center}
\begin{tabularx}{\linewidth}{|X|X|X|X|}
\hline
\rowcolor{popblueL} \textbf{Caractères} & \textbf{Pinyin} & \textbf{Nature} & \textbf{Traduction} \\ \hline
\zh{山村} & shāncūn & n. & village de montagne \\ \hline
\zh{收成} & shōuchéng & n. & récolte \\ \hline
\zh{玉米} & yùmǐ & n. & maïs \\ \hline
\zh{路费} & lùfèi & n. & frais de voyage \\ \hline
\zh{夜班} & yèbān & n. & travail de nuit / équipe de nuit \\ \hline
\zh{节奏} & jiézòu & n. & rythme \\ \hline
\zh{食堂} & shítáng & n. & cantine, réfectoire \\ \hline
\zh{酸辣汤} & suānlàtāng & n. & soupe aigre-épicée (spécialité) \\ \hline
\zh{组长} & zǔzhǎng & n. & chef d'équipe, contremaître \\ \hline
\zh{盖房子} & gài fángzi & VO & construire une maison \\ \hline
\zh{存钱} & cún qián & VO & épargner, déposer de l'argent \\ \hline
\zh{驾照} & jiàzhào & n. & permis de conduire \\ \hline
\zh{越来越好} & yuè lái yuè hǎo & adj./structure & de mieux en mieux \\ \hline
\end{tabularx}
\end{center}

\textbf{Questions de compréhension / analyse structurelle :}
\begin{enumerate}
    \item Identifiez les 5 étapes du plan dans le texte (relevez la phrase correspondante).
    \item Soulignez tous les connecteurs étudiés (\zh{因为, 所以, 而且, 为了, 先, 再, 但是, 可是}) et notez leur fonction.
    \item Pourquoi l'auteur utilise-t-il \zh{可是} à la ligne 3 au lieu de \zh{但是} ? (Indice : subjectivité, ressenti physique).
\end{enumerate}
\end{textebox}

\coursec{Lexique écrit du thème : L'exode rural}

\begin{definition}[Vocabulaire thématique classé (HSK 2-3)]
Le lexique suivant regroupe les mots-outils et noms essentiels pour produire un texte sur l'exode rural. Maîtrisez l'écriture des caractères en gras.
\end{definition}

\begin{center}
\begin{tabularx}{\linewidth}{|X|X|X|X|}
\hline
\rowcolor{popblueL} \textbf{Caractères} & \textbf{Pinyin} & \textbf{Nature} & \textbf{Traduction} \\ \hline
\zh{农村} & nóngcūn & n. & campagne, zone rurale \\ \hline
\zh{城市} & chéngshì & n. & ville \\ \hline
\zh{扈口} & hùkǒu & n. & registre de résidence (hukou) \\ \hline
\zh{民工} & míngōng & n. & travailleur migrant (paysan-ouvrier) \\ \hline
\zh{打工} & dǎgōng & v. & travailler (souvent manuel, loin de chez soi) \\ \hline
\zh{离开} & líkāi & v. & quitter, partir \\ \hline
\zh{到达} & dàodá & v. & arriver à \\ \hline
\zh{适应} & shìyìng & v. & s'adapter à \\ \hline
\zh{想念} & xiǎngniàn & v. & regretter, avoir le mal de (qn/qch) \\ \hline
\zh{艰苦} & jiānkǔ & adj. & dur, pénible (vie, conditions) \\ \hline
\zh{穷} & qióng & adj. & pauvre \\ \hline
\zh{孤独} & gūdú & adj. & seul, solitaire \\ \hline
\zh{务工} & wùgōng & v./n. & travailler comme migrant / travail migrant \\ \hline
\zh{汇款} & huìkuǎn & v./n. & envoyer de l'argent / mandat, transfert \\ \hline
\zh{返乡} & fǎnxiāng & v. & retourner au village natal \\ \hline
\zh{盖房子} & gài fángzi & VO & construire une maison \\ \hline
\zh{交学费} & jiāo xuéfèi & VO & payer les frais de scolarité \\ \hline
\zh{看病} & kàn bìng & VO & se soigner, consulter un médecin \\ \hline
\end{tabularx}
\end{center}

\coursec{Écriture des caractères complexes : Analyse graphémique}

\begin{propriete}[Caractères clés du thème : traits, radicaux, pièges]
Voici l'analyse de 8 caractères complexes fréquents dans cette leçon. Respectez l'ordre des traits : \textbf{haut $\rightarrow$ bas, gauche $\rightarrow$ droite, horizontal avant vertical, trait central avant les ailes, cadre avant contenu, fermeture en dernier}.
\end{propriete}

\begin{center}
\begin{tabularx}{\linewidth}{|X|X|X|X|X|}
\hline
\rowcolor{popblueL} \textbf{Car.} & \textbf{Pinyin} & \textbf{Radical (Clé)} & \textbf{Nb traits} & \textbf{Analyse / Pièges / Caractères proches} \\ \hline
\zh{穷} & qióng & \zh{穴} (caverne, xué) & 10 & Haut : \zh{穴} (caverne) + Bas : \zh{力} (force). \emph{Sens :} force épuisée dans une caverne $\rightarrow$ pauvre. \textbf{Piège :} ne pas confondre avec \zh{究} (jiū, enquêter) qui a \zh{九} en bas. \\ \hline
\zh{迁} & qiān & \zh{辶} (marche, chuò) & 7 & Radical \zh{辶} (mouvement) + phonétique \zh{仟} (qiān, mille). \emph{Sens :} déplacer (mille fois). \textbf{Ordre :} écrire \zh{仟} (haut-gauche $\rightarrow$ bas-droite) PUIS \zh{辶} (point, horizontal, pli). \\ \hline
\zh{务} & wù & \zh{力} (force, lì) & 11 & Simplifié de \zh{務}. Haut : \zh{亠} + \zh{吴} (phonétique) ; Bas : \zh{力}. \emph{Sens :} affaire, devoir. \textbf{Piège :} ne pas écrire \zh{无} (wú, sans) à la place du centre. \\ \hline
\zh{偿} & cháng & \zh{亻} (homme, rén) & 11 & \zh{亻} + \zh{常} (cháng, constant). \emph{Sens :} rembourser (homme constant dans ses dettes). \textbf{Proche :} \zh{尝} (cháng, goûter) a \zh{小} en haut. \\ \hline
\zh{酸} & suān & \zh{酉} (alcool/vin, yǒu) & 14 & \zh{酉} (radical liquide fermenté) + \zh{夋} (phonétique). \emph{Sens :} goût du vin tourné. \textbf{Ordre :} \zh{酉} (horizontal, vertical, crochet, horizontal, horizontal, vertical, horizontal) puis le reste. \\ \hline
\zh{辣} & là & \zh{辛} (piquant, xīn) & 16 & \zh{辛} (piquant) + \zh{来} (lái, phonétique). \emph{Sens :} piquant qui vient. \textbf{Piège :} \zh{辛} $\neq$ \zh{幸} (xìng, bonheur). \zh{辛} a un trait horizontal de plus en haut. \\ \hline
\zh{组} & zǔ & \zh{丝} (fil, sī) & 10 & \zh{丝} (fil) + \zh{且} (qiě, phonétique). Simplifié de \zh{組}. \emph{Sens :} assembler des fils $\rightarrow$ groupe. \\ \hline
\zh{驾} & jià & \zh{马} (cheval, mǎ) & 8 & \zh{马} + \zh{加} (jiā, phonétique). Simplifié de \zh{駕}. \emph{Sens :} conduire un char (cheval + ajouter du joug). \\ \hline
\zh{盖} & gài & \zh{皿} (récipient, mǐn) & 11 & Haut : \zh{盍} (hé, couvercle + oiseau) + Bas : \zh{皿}. \emph{Sens :} couvrir un récipient $\rightarrow$ construire (toit). \textbf{Proche :} \zh{益} (yì, profit) a \zh{皿} en bas mais haut différent. \\ \hline
\zh{存} & cún & \zh{子} (enfant, zǐ) & 6 & Haut : \zh{在} (exister/être à) + Bas : \zh{子}. \emph{Sens :} l'enfant existe $\rightarrow$ conserver, déposer. \textbf{Proche :} \zh{在} (zài) n'a pas \zh{子}. \\ \hline
\end{tabularx}
\end{center}

\begin{methode}[Entraînement à l'écriture (Méthode des 3 temps)]
\begin{enumerate}
    \item \textbf{Observation :} Identifiez le radical (clé de dictionnaire) et la partie phonétique. Dites l'ordre des traits à voix haute.
    \item \textbf{Traçage guidé :} Reproduisez le caractère 3 fois en grand format sur papier quadrillé (格子纸), en respectant scrupuleusement l'ordre des traits.
    \item \textbf{Mémorisation contextuelle :} Écrivez 2 mots composés contenant ce caractère (ex: \zh{穷} $\rightarrow$ \zh{贫穷}, \zh{穷人}) et une phrase courte.
\end{enumerate}
\end{methode}

\coursec{Structure et connecteurs du texte}

Après avoir maîtrisé le vocabulaire et l'écriture des caractères complexes du thème (\emph{exode rural, ville, campagne, travail, famille, espoir}), nous abordons maintenant l'architecture du texte. En chinois, comme en français, un texte cohérent ne se résume pas à une suite de phrases ; il repose sur une \cle{structure logique} explicite, balisée par des \cle{connecteurs} (关联词 \zh{guānliáncí}) qui guident le lecteur. Cette section vous apprendra à organiser vos idées selon le plan type du récit migratoire et à utiliser les cinq connecteurs indispensables pour l'exprimer avec fluidité.

\courssub{Plan type d'un texte sur l'exode rural : les cinq étapes narratives}

Tout texte traitant de l'exode rural, qu'il soit un témoignage, une lettre ou une petite narration, suit généralement une progression en cinq temps. Respecter cet ordre permet de produire un texte \emph{cohérent} et \emph{convaincant}, attendu dans les épreuves d'expression écrite du baccalauréat.

\begin{aretenir}[Plan type en 5 étapes]
\begin{enumerate}
    \item \textbf{La situation au village (农村的情况)} : Description du cadre de vie initial (pauvreté, manque d'écoles, terres insuffisantes, vie dure).
    \item \textbf{La cause du départ (离开的原因)} : L'élément déclencheur (sécheresse, maladie, frais de scolarité, envie de changement). C'est le \cle{moteur} de l'action.
    \item \textbf{L'arrivée en ville (到达城市)} : Le choc, les premières impressions (bruit, bâtiments, circulation, recherche de logement).
    \item \textbf{Avantages et difficultés (优点与困难)} : Bilan contrasté : salaire vs loyer cher, liberté vs solitude, commodités vs pollution.
    \item \textbf{Projet ou espoir (计划与希望)} : Perspective d'avenir (économiser, faire venir la famille, retourner au village, études des enfants).
\end{enumerate}
\end{aretenir}

\begin{exemplebox}[Application du plan : Micro-texte structuré]
Voici un micro-texte de 5 phrases respectant exactement ce plan. Observez les connecteurs en \textbf{gras}.

\begin{center}
\begin{tabularx}{\linewidth}{|X|X|X|}
\hline
\rowcolor{popblueL} \textbf{Étape} & \textbf{Chinois (Caractères + Pinyin)} & \textbf{Traduction} \\ \hline
1. Situation & \zh{我家在北方的一个小村庄，生活很艰苦。} & Ma famille est dans un petit village du Nord, la vie est très dure. \\
& \emph{Wǒ jiā zài běifāng de yī gè xiǎo cūnzhuāng, shēnghuó hěn jiānkǔ.} & \\ \hline
2. Cause & \zh{\textbf{因为} 爸爸生病了，\textbf{所以} 我必须出来打工。} & \textbf{Parce que} papa est malade, \textbf{donc} je dois sortir travailler. \\
& \emph{Yīnwèi bàba shēngbìng le, suǒyǐ wǒ bìxū chūlái dǎgōng.} & \\ \hline
3. Arrivée & \zh{\textbf{先} 我到了广州，\textbf{再} 找到了一份工厂的工作。} & \textbf{D'abord} j'arrive à Guangzhou, \textbf{ensuite} je trouve un travail d'usine. \\
& \emph{Xiān wǒ dào le Guǎngzhōu, zài zhǎodào le yī fèn gōngchǎng de gōngzuò.} & \\ \hline
4. Bilan & \zh{工资不低，\textbf{但是} 房租很贵，\textbf{而且} 我很想家。} & Le salaire n'est pas bas, \textbf{mais} le loyer est cher, \textbf{de plus} j'ai le mal du pays. \\
& \emph{Gōngzī bù dī, dànshì fángzū hěn guì, érqiě wǒ hěn xiǎng jiā.} & \\ \hline
5. Espoir & \zh{\textbf{为了} 妹妹上大学，我会继续努力。} & \textbf{Pour} que ma sœur aille à l'université, je vais continuer à m'efforcer. \\
& \emph{Wèile mèimei shàng dàxué, wǒ huì jìxù nǔlì.} & \\ \hline
\end{tabularx}
\end{center}
\end{exemplebox}

\courssub{Les cinq connecteurs indispensables : forme, emploi et comparaison}

Nous allons analyser chaque connecteur à travers sa structure syntaxique, son emploi précis et les pièges de traduction depuis le français.

\courssub{La cause et la conséquence : \zh{因为 … 所以 …} (Yīnwèi … suǒyǐ …)}

C'est la structure reine pour expliquer \emph{pourquoi} quelque chose se produit.

\begin{propriete}[Règle de \zh{因为 … 所以 …}]
\textbf{Structure :} \cle{Sujet + 因为 + Cause , 所以 + Conséquence .}

\textbf{Points clés :}
\begin{itemize}
    \item \textbf{Les deux marqueurs coexistent souvent.} En chinois standard écrit, on garde fréquemment \zh{因为} au début de la proposition de cause ET \zh{所以} au début de la proposition de conséquence. C'est \textbf{correct et élégant}.
    \item \textbf{Attention francophones :} En français, « Parce que..., donc... » est une \emph{faute de syntaxe} (pléonasme syntaxique). On dit « Parce que..., ... » OU « ..., donc ... ». En chinois, \zh{因为…所以…} est la norme.
    \item \textbf{Séparation possible :} À l'oral ou dans des phrases courtes, on peut n'utiliser que \zh{因为} (en fin de phrase : \zh{...，因为...}) ou que \zh{所以} (en début de phrase : \zh{所以...}).
    \item \textbf{Position du sujet :} Si le sujet est le même pour les deux propositions, il se place généralement \textbf{avant} \zh{因为} (\zh{我因为...所以...}). Si les sujets diffèrent, il se place \textbf{après} \zh{因为} et \zh{所以} (\zh{因为下雨，所以路湿了}).
\end{itemize}
\end{propriete}

\begin{exemplebox}[Variantes de \zh{因为 … 所以 …}]
\begin{itemize}
    \item \zh{因为城市工作机会多，所以很多年轻人离开农村。} \\
    (Parce que la ville a beaucoup d'opportunités d'emploi, beaucoup de jeunes quittent la campagne.)
    \item \zh{他不想走，因为父母年纪大了。} (Il ne veut pas partir, \textbf{parce que} ses parents sont âgés.) \textit{[Seul 因为, cause en fin de phrase]}
    \item \zh{下雨了，所以比赛取消了。} (Il a plu, \textbf{donc} le match est annulé.) \textit{[Seul 所以, conséquence en début de phrase]}
\end{itemize}
\end{exemplebox}

\begin{popfigure}
\begin{tikzpicture}[
    node distance=1.2cm and 1.5cm,
    box/.style={draw=popblue, thick, rounded corners, fill=popblueL, minimum width=2.5cm, minimum height=1cm, align=center, font=\small},
    arrow/.style={-Stealth, thick, popdark},
    labelnode/.style={font=\tiny, text=popmuted, align=center}
]
% Nodes
\node[box, align=center] (cause){因为 \zh{(yīnwèi)}\\ \textbf{CAUSE}};
\node[box, right=of cause, align=center] (conseq){所以 \zh{(suǒyǐ)}\\ \textbf{CONSEQUENCE}};
\node[box, below=of cause, align=center] (opp){但是 / 可是\\ \zh{(dànshì / kěshì)}\\ \textbf{OPPOSITION}};
\node[box, right=of opp, align=center] (purp){为了 \zh{(wèile)}\\ \textbf{BUT}};
\node[box, below=of opp, align=center] (seq){先 … 再 …\\ \zh{(xiān … zài …)}\\ \textbf{SÉQUENCE}};
\node[box, right=of seq, align=center] (add){而且 \zh{(érqiě)}\\ \textbf{ADDITION}};

% Arrows
\draw[arrow] (cause) -- node[above, labelnode] {lien logique fort} (conseq);
\draw[arrow, poporange] (opp.west) -- ++(-1,0) |- (cause.south) node[pos=0.25, above left, labelnode] {contre-argument};
\draw[arrow, poppurple] (purp) -- node[above, labelnode] {objectif visé} (conseq);
\draw[arrow, popgreen] (seq) -| (cause) node[pos=0.5, left, labelnode] {ordre chrono};
\draw[arrow, poppink] (add) -- node[above, labelnode] {argument +} (opp);

% Title
\node[above=0.3cm of cause, font=\bfseries\small, text=popdark] {Chaîne logique des connecteurs (Exode rural)};
\end{tikzpicture}
\legende{Schéma des relations logiques : Cause $\rightarrow$ Conséquence (bleu), Opposition (orange), But (violet), Séquence (vert), Addition (rose).}
\end{popfigure}

\courssub{L'opposition : \zh{但是} (dànshì) / \zh{可是} (kěshì)}

Utilisés pour introduire un obstacle, une nuance ou une réalité contrastée (étape 4 du plan : avantages vs difficultés).

\begin{propriete}[Règle de \zh{但是} / \zh{可是}]
\textbf{Structure :} \cle{Proposition 1 , 但是/可是 + Proposition 2 .}

\textbf{Nuances :}
\begin{itemize}
    \item \zh{但是} : Neutre, standard, écrit et oral. \textbf{Le choix par défaut pour le bac.}
    \item \zh{可是} : Légèrement plus oral, plus subjectif, insiste sur le ressenti (« pourtant », « mais » avec émotion).
    \item \textbf{Position du sujet :} Si le sujet change, le répéter après \zh{但是}. S'il est identique, on peut l'omettre (ellipse du sujet).
\end{itemize}
\end{propriete}

\begin{exemplebox}[Opposition dans le contexte urbain]
\begin{itemize}
    \item \zh{城市很方便，\textbf{但是} 房租太贵了。} (La ville est commode, \textbf{mais} le loyer est trop cher.) \textit{[Sujets différents : ville / loyer]}
    \item \zh{我想回家，\textbf{可是} (我) 没有钱买票。} (Je veux rentrer, \textbf{mais} (je) n'ai pas d'argent pour le billet.) \textit{[Sujet identique '我', omis la 2e fois]}
\end{itemize}
\end{exemplebox}

\courssub{L'addition : \zh{而且} (érqiě)}

Indispensable pour accumuler les arguments (étape 4 : « le salaire est bas \textbf{et en plus} le patron est méchant »).

\begin{propriete}[Règle de \zh{而且}]
\textbf{Structure :} \cle{Proposition 1 , 而且 + Proposition 2 .}

\textbf{Emploi :}
\begin{itemize}
    \item Signifie « de plus », « en outre », « qui plus est ».
    \item Souvent utilisé pour \textbf{renforcer un trait négatif} ou \textbf{ajouter un argument dans le même sens}.
    \item Ne s'utilise \textbf{pas} pour une simple énumération d'actions successives (on utilise \zh{还} ou \zh{也} pour « aussi » simple). \zh{而且} implique une \textbf{progression dans l'argumentation}.
\end{itemize}
\end{propriete}

\begin{exemplebox}[Addition argumentative]
\begin{itemize}
    \item \zh{工作很累，\textbf{而且} 工资很低。} (Le travail est fatiguant, \textbf{de plus} le salaire est bas.) \textit{[Double argument négatif]}
    \item \zh{他很聪明，\textbf{而且} 很努力。} (Il est intelligent, \textbf{de plus} il est travailleur.) \textit{[Double argument positif]}
    \item \textit{Erreur fréquente :} \zh{我买了米，而且买了菜。} $\rightarrow$ Préférer \zh{我买了米，\textbf{还} 买了菜。} (J'ai acheté du riz, \textbf{aussi} des légumes.) \textit{[Simple énumération d'actions]}
\end{itemize}
\end{exemplebox}

\courssub{Le but : \zh{为了} (wèile)}

Le connecteur clé pour l'étape 5 (projet/espérance) et l'étape 2 (cause profonde).

\begin{propriete}[Règle de \zh{为了}]
\textbf{Structure :} \cle{为了 + Nom / Verbe (+ 的) , Sujet + Verbe d'action .}

\textbf{Règles de construction :}
\begin{itemize}
    \item \textbf{但 + Nom :} \zh{为了家人} (Pour la famille), \zh{为了钱} (Pour l'argent).
    \item \textbf{但 + Verbe :} Ajouter \zh{的} après le verbe si l'objet est court ou pour nominaliser : \zh{为了赚钱(的)} (Pour gagner de l'argent), \zh{为了让孩子上学(的)} (Pour que les enfants aillent à l'école).
    \item \textbf{Position :} La proposition de but (\zh{为了...}) se place \textbf{TOUJOURS AVANT} le sujet principal de la phrase (sujet + verbe principal). En français, « Pour..., je... » ou « Je... pour... ».
\end{itemize}
\end{propriete}

\begin{attention}[Piège majeur : \zh{为了} vs \zh{给} (gěi)]
Les francophones traduisent souvent « pour » par \zh{给} (donner à / pour [bénéficiaire direct]).
\begin{itemize}
    \item \textbf{\zh{为了} = But, finalité, motivation profonde} (Pourquoi on agit ?). \checkmark \zh{为了家人，我努力工作。}
    \item \textbf{\zh{给} = Bénéficiaire direct, destinataire d'une action de transfert} (À qui ?). \checkmark \zh{我给家人寄钱。} (J'envoie de l'argent \textbf{à} ma famille.)
    \item \textbf{Interdit :} \zh{*给家人，我努力工作。} (Incorrect pour exprimer le but).
\end{itemize}
\end{attention}

\begin{exemplebox}[Le but dans le récit migratoire]
\begin{itemize}
    \item \zh{\textbf{为了} 孩子的学习，他们搬到了城市。} (Pour les études des enfants, ils ont déménagé en ville.)
    \item \zh{他 \textbf{为了} \textbf{赚钱} ，\textbf{去} 了 深圳 。} (Il \textbf{est allé} à Shenzhen \textbf{pour gagner de l'argent}.)
    \item \zh{我们 \textbf{为了} \textbf{建设 家乡} ，\textbf{决定} 回去 。} (Nous avons décidé de retourner \textbf{pour construire notre village natal}.)
\end{itemize}
\end{exemplebox}

\courssub{La séquence temporelle : \zh{先 … 再 …} (xiān … zài …)}

Essentiel pour l'étape 3 (arrivée puis installation) et pour narrer la chronologie.

\begin{propriete}[Règle de \zh{先 … 再 …}]
\textbf{Structure :} \cle{Sujet + 先 + Action 1 , 再 + Action 2 .}

\textbf{Emploi :}
\begin{itemize}
    \item Indique une \textbf{succession stricte} : Action 1 terminée $\rightarrow$ Action 2 commence.
    \item \zh{先} (d'abord) et \zh{再} (ensuite/puis) sont des adverbes placés \textbf{avant le verbe}.
    \item Différence avec \zh{然后} (ránhòu) : \zh{然后} est une conjonction (début de phrase), \zh{先…再…} est une structure corrélative interne à la phrase.
\end{itemize}
\end{propriete}

\begin{exemplebox}[Chronologie de l'installation]
\begin{itemize}
    \item \zh{他想 \textbf{先} 找 房子 ，\textbf{再} 找 工作 。} (Il veut \textbf{d'abord} trouver un logement, \textbf{ensuite} un travail.)
    \item \zh{我们 \textbf{先} 到 车站 ，\textbf{再} 坐 地铁 。} (Nous allons \textbf{d'abord} à la gare, \textbf{puis} on prend le métro.)
\end{itemize}
\end{exemplebox}

\courssub{Tableau récapitulatif comparatif : Connecteurs Français / Chinois}

\begin{center}
\begin{tabularx}{\linewidth}{|>{\bfseries}X|X|X|X|}
\hline
\rowcolor{popblueL} \textbf{Fonction} & \textbf{Connecteur Chinois} & \textbf{Structure Type} & \textbf{Piège Français $\rightarrow$ Chinois} \\ \hline
Cause $\rightarrow$ Conséquence & \zh{因为 … 所以 …} & S + 因为 + Cause, 所以 + Cons. & Garder les DEUX mots (interdit en FR « parce que... donc »). \\ \hline
Opposition / Concession & \zh{但是} (standard) \newline \zh{可是} (oral/émotif) & Prop 1, 但是 + Prop 2 & Ne pas confondre avec \zh{不过} (néanmoins, plus faible). \\ \hline
Addition / Renforcement & \zh{而且} & Prop 1, 而且 + Prop 2 & Ne pas utiliser pour simple « aussi » (\zh{也/还}). \\ \hline
But / Finalité & \zh{为了} + N / V(+的) & 为了 + But, S + Action & \textbf{Jamais} \zh{给} pour le but. \zh{给} = destinataire. \\ \hline
Séquence temporelle & \zh{先 … 再 …} & S + 先 + V1, 再 + V2 & \zh{先/再} sont des adverbes (devant V), pas des conjonctions. \\ \hline
\end{tabularx}
\end{center}

\coursec{Thème et version guidés}

\courssub{Méthode : Rédiger un texte structuré sur l'exode rural (5 étapes)}

\begin{methode}[De l'idée au texte chinois en 4 étapes]
\begin{enumerate}
    \item \textbf{Planifier en français (5 lignes).} Rédigez une phrase par étape du plan type (Situation, Cause, Arrivée, Bilan, Espoir). Exemple : « 1. Mon village est isolé. 2. Mon père est malade. 3. J'arrive à Douala, je cherche du travail. 4. Le salaire est OK mais le loyer est cher. 5. Je veux soigner mon père. »
    \item \textbf{Associer un connecteur cible par étape.}
    \begin{itemize}
        \item Étape 2 (Cause) $\rightarrow$ \zh{因为…所以…}
        \item Étape 3 (Chronologie) $\rightarrow$ \zh{先…再…}
        \item Étape 4 (Contraste) $\rightarrow$ \zh{但是/而且}
        \item Étape 5 (But) $\rightarrow$ \zh{为了}
    \end{itemize}
    \item \textbf{Rédiger en chinois (phrase par phrase).} Traduisez chaque ligne du plan en appliquant la structure du connecteur choisi. Vérifiez l'ordre des mots (Sujet + Connecteur + Verbe).
    \item \textbf{Relier et peaufiner.} Ajoutez des mots de liaison légers (\zh{后来} « plus tard », \zh{现在} « maintenant »), vérifiez les particules (\zh{了, 的, 地}) et l'accord des temps (aspects).
\end{enumerate}
\end{methode}

\courssub{Exercice résolu : Transformation et assemblage}

\begin{exoresolu}[Assembler un paragraphe cohérent]
\textbf{Consigne :} Reliez les propositions suivantes en un seul paragraphe cohérent en utilisant \textbf{obligatoirement} : \zh{因为…所以…}, \zh{为了}, \zh{先…再…}, \zh{但是}, \zh{而且}. Respectez l'ordre logique (Cause $\rightarrow$ But $\rightarrow$ Séquence $\rightarrow$ Opposition/Addition).

\textbf{Propositions de base (désordonnées) :}
\begin{enumerate}
    \item Il veut acheter une maison au village. (\zh{想 在 老家 买 房子})
    \item Il est parti à Yaoundé. (\zh{去 了 雅温得})
    \item Le village est pauvre. (\zh{家乡 很 穷})
    \item Il a trouvé un travail de gardien. (\zh{找 到 看门 的 工作})
    \item Il a emprunté de l'argent pour le voyage. (\zh{借 了 路费})
    \item Le salaire est bas. (\zh{工资 低})
    \item Le patron est méchant. (\zh{老板 凶})
    \item Il veut aider ses parents. (\zh{帮助 父母})
\end{enumerate}

\tcblower

\textbf{Solution détaillée :}

\textbf{Étape 1 : Ordonner logiquement (Plan 5 étapes).}
1. Situation : \zh{家乡很穷。}
2. Cause/Départ : \zh{因为家乡很穷，所以他去了雅温得。} (Intégration \zh{因为…所以…})
3. But profond : \zh{为了帮助父母，他想在老家买房子。} (Intégration \zh{为了})
4. Séquence Arrivée : \zh{他先借了路费，再去了雅温得。} (Intégration \zh{先…再…}) \textit{Note : chronologiquement, emprunter précède le départ.}
5. Bilan Travail : \zh{他找到了看门的工作，但是工资低，而且老板凶。} (Intégration \zh{但是} + \zh{而且})

\textbf{Étape 2 : Fusion en paragraphe (niveau HSK 3).}

\zh{我的家乡很穷，\textbf{因为} 没有工厂，\textbf{所以} 年轻人都走光了。\textbf{为了} \textbf{帮助} 父母，\textbf{也} \textbf{为了} \textbf{在} 老家 \textbf{买} 房子，我 \textbf{先} 跟 哥哥 \textbf{借} 了 路费，\textbf{再} 坐 火车 \textbf{去} 了 雅温得。刚开始 我 \textbf{找} 到 了 看门 的 工作，\textbf{但是} 工资 很 低，\textbf{而且} 老板 很 凶，每天 都 吼 人。不过 我 \textbf{因为} 有 目标，\textbf{所以} 很 坚强。}

\textbf{Analyse des connecteurs utilisés :}
\begin{itemize}
    \item \zh{因为…所以…} (x2 : cause pauvreté/départ ; cause objectif/force).
    \item \zh{为了} (x2 : but aider parents, but acheter maison).
    \item \zh{先…再…} (chronologie prêt $\rightarrow$ voyage).
    \item \zh{但是…而且…} (contraste salaire/patron + addition argument négatif).
\end{itemize}
\end{exoresolu}

\courssub{Exercices d'entraînement}

\begin{exercice}[Mise en pratique : Connecteurs et Structure]
\textbf{Exercice 1 : Complétez avec le connecteur adapté (\zh{因为, 所以, 为了, 但是, 而且, 先, 再}).}
\begin{enumerate}
    \item \zh{\_\_\_\_\_\_} 下雨了，\_\_\_\_\_\_ 比赛取消了。
    \item \zh{我 \_\_\_\_\_\_ 买 票，\_\_\_\_\_\_ 上 车。}
    \item \zh{\_\_\_\_\_\_} 家人 的 健康，他 每天 跑步。
    \item \zh{这 个 手机 很 好看，\_\_\_\_\_\_ 太 贵了，\_\_\_\_\_\_ 电池 不 耐用。}
    \item \zh{他 不 来 了，\_\_\_\_\_\_ 他 生病 了。} (Seul connecteur de cause en fin de phrase).
\end{enumerate}

\textbf{Exercice 2 : Thème (Français $\rightarrow$ Chinois). Traduisez en utilisant les structures ciblées.}
\begin{enumerate}
    \item Parce que la campagne est pauvre, donc mon frère est parti travailler en ville. (Structure \zh{因为…所以…})
    \item Pour gagner de l'argent, il a d'abord appris la conduite, ensuite il est devenu chauffeur. (\zh{为了} + \zh{先…再…})
    \item La ville est belle, mais l'air n'est pas bon, de plus c'est bruyant. (\zh{但是} + \zh{而且})
    \item Je veux retourner au village pour construire une maison pour mes parents. (\zh{为了} + but nominalisé)
\end{enumerate}

\textbf{Exercice 3 : Version (Chinois $\rightarrow$ Français). Traduisez en français.}
\begin{enumerate}
    \item \zh{因为家里穷，所以他去城市打工。}
    \item \zh{为了给妈妈看病，他先借钱，再去了医院。}
    \item \zh{城市工作机会多，但是生活压力大，而且很想家。}
    \item \zh{他先存钱，再盖房子，为了将来生活更好。}
\end{enumerate}

\textbf{Exercice 4 : Production écrite guidée (Mini-rédaction \textasciitilde 80 caractères).}
Rédigez l'histoire de \zh{小李} (Xiǎo Lǐ) en suivant \textbf{strictement} le plan en 5 étapes et en utilisant \textbf{au moins 4 connecteurs différents} parmi ceux étudiés.
\begin{itemize}
    \item Situation : Village isolé, parents âgés.
    \item Cause : Frais médicaux élevés.
    \item Arrivée : Bus pour Douala, loge chez cousin.
    \item Bilan : Travail dur (chantier), mais gagne bien sa vie.
    \item Espoir : Dans 3 ans, rentrer au village, ouvrir petit commerce.
\end{itemize}
\end{exercice}

\begin{corrige}[Exercice 1]
\begin{enumerate}
    \item \zh{因为} 下雨了，\zh{所以} 比赛取消了。 (Standard)
    \item \zh{我 \textbf{先} 买 票，\textbf{再} 上 车。} (Séquence)
    \item \zh{\textbf{为了}} 家人 的 健康，他 每天 跑步。 (But)
    \item \zh{这 个 手机 很 好看，\textbf{但是} 太 贵了，\textbf{而且} 电池 不 耐用。} (Opposition + Addition)
    \item \zh{他 不 来 了，\textbf{因为} 他 生病 了。} (Cause en fin de phrase, seul \zh{因为}).
\end{enumerate}
\end{corrige}

\begin{corrige}[Exercice 2]
\begin{enumerate}
    \item \zh{因为 农村 很 穷，\textbf{所以} 我 哥哥 去 城市 打工 了。}
    \item \zh{\textbf{为了} 赚钱，他 \textbf{先} 学 了 开车，\textbf{再} 当 了 司机。}
    \item \zh{城市 很 漂亮，\textbf{但是} 空气 不 好，\textbf{而且} 很 吵。}
    \item \zh{\textbf{为了} \textbf{给} 父母 \textbf{盖} 房子，我 想 回 老家。}
\end{enumerate}
\end{corrige}

\begin{corrige}[Exercice 3]
\begin{enumerate}
    \item Parce que la famille est pauvre, il est parti travailler en ville.
    \item Pour soigner sa mère, il a d'abord emprunté de l'argent, ensuite il est allé à l'hôpital.
    \item La ville a beaucoup d'opportunités d'emploi, mais la pression de vie est forte, et en plus il a le mal du pays.
    \item Il économise d'abord, puis construit une maison, pour que la vie future soit meilleure.
\end{enumerate}
\end{corrige}

\begin{corrige}[Exercice 4 — Proposition de rédaction modèle]
\zh{小李 的 家 在 深山 里，父母 年纪 大 了。\textbf{因为} 爸爸 生病 住院，\textbf{所以} 需要 很多 钱。\textbf{为了} 赚 医药费，小李 \textbf{先} 坐 大巴 去 了 杜阿拉，\textbf{再} 住 在 表哥 家。现在 他在 工地 搬 砖，\textbf{虽然} 很 累，\textbf{但是} 工资 高。\textbf{为了} 三 年 后 回 老家 开 小店，他 很 努力。}

\textit{Note : J'ai ajouté \zh{虽然…但是…} (bien que... mais...) qui est une structure d'opposition renforcée (HSK 3), excellente pour le bilan.}
\end{corrige}

\coursec{Production orale et évaluation}

\begin{experience}[Jeu de rôle : « L'entretien au retour du village »]
\textbf{Contexte :} Vous êtes \zh{阿勇} (ou \zh{小李}) de retour au village pour le Nouvel An (\zh{春节}). Un voisin (camarade) vous interroge sur votre vie en ville.
\textbf{Rôles :}
\begin{itemize}
    \item \textbf{A (Migrant) :} Utilisez le plan en 5 étapes + connecteurs. Parlez de : situation passée, cause départ, arrivée, bilan (vrai/faux), projet.
    \item \textbf{B (Voisin) :} Posez des questions avec \zh{为什么} (pourquoi), \zh{怎么样} (comment), \zh{以后打算} (projets futurs).
\end{itemize}
\textbf{Contraintes :} Minimum 5 tours de parole chacun. Utiliser \zh{因为…所以…}, \zh{为了}, \zh{先…再…}, \zh{但是/而且} au moins une fois chacun.
\textbf{Évaluation par les pairs :} Grille : Prononciation (tons), Fluidité, Richesse vocabulaire, Exactitude structures cibles.
\end{experience}

\begin{experience}[Débat structuré : « Faut-il partir ou rester ? »]
\textbf{Sujet :} \zh{年轻人应该留在农村建设家乡，还是去城市发展？}
\textbf{Organisation :} Deux groupes (Partir / Rester). Préparation 10 min avec fiches arguments (vocabulaire fourni). Débat 15 min. Synthèse écrite commune : 3 arguments pour chaque camp avec connecteurs \zh{因为…所以…}, \zh{虽然…但是…}, \zh{为了}.
\end{experience}

\coursec{Point culturel : L'exode rural, miroir Chine-Cameroun}

\begin{savaistu}[Le « Mingong » (民工) et le « Bushfaller » : deux figures migrantes]
Le phénomène de l'exode rural est au cœur de l'histoire récente des deux pays.
\begin{itemize}
    \item \textbf{Chine :} Depuis les années 1980, plus de \textbf{290 millions} de \zh{民工} (míngōng — travailleurs migrants ruraux) ont quitté les provinces intérieures (Sichuan, Henan, Anhui) vers les zones côtières (Guangdong, Zhejiang, Shanghai). Ils ont bâti les villes modernes. Le \zh{春运} (chūnyùn, « transport du Nouvel An ») est la plus grande migration humaine annuelle : des centaines de millions de \zh{民工} rentrent au village pour le Nouvel An chinois, unique moment de retrouvailles familiales. Le \zh{户口} (hùkǒu, registre de résidence) crée encore des inégalités d'accès aux écoles et hôpitaux pour leurs enfants en ville.
    \item \textbf{Cameroun :} L'exode alimente Yaoundé et Douala depuis les années 60. Le terme « \emph{Bushfaller} » (celui qui « tombe de la brousse » / réussit à l'étranger) désignait à l'origine ceux qui partaient en Europe, mais s'applique aujourd'hui aux ruraux qui « réussissent » en ville. Le \zh{砖房} (zhuānfáng — maison en briques/ciment) au village reste le symbole ultime de la réussite du migrant (comme \zh{盖房子} dans notre texte).
    \item \textbf{Comparaison :} En Chine, le retour définitif au village (\zh{返乡} fǎnxiāng) est de plus en plus fréquent grâce au développement des counties. Au Cameroun, le retour est souvent lié à la retraite. Dans les deux cas, la \textbf{remise d'argent} (\zh{汇款} huìkuǎn / \emph{transferts d'argent}) structure l'économie familiale villageoise.
\end{itemize}
\end{savaistu}

\coursec{Thème et version guidés}

Dans la section précédente, nous avons appris à organiser un texte simple sur l'exode rural grâce aux connecteurs \zh{因为} (parce que), \zh{但是} (mais), \zh{为了} (pour, afin de), \zh{越来越} (de plus en plus) et \zh{虽然……但是……} (bien que… mais…). Nous allons maintenant mettre ces outils en action dans l'exercice roi du cours de chinois : la \cle{traduction}. Le \cle{thème} consiste à traduire du français vers le chinois ; la \cle{version} consiste à traduire du chinois vers le français. Ces deux exercices ne se ressemblent pas : dans le thème, la difficulté est de \emph{construire} une phrase chinoise correcte ; dans la version, elle est de \emph{comprendre} la phrase chinoise et de la restituer en français naturel.

\courssub{Observer avant de traduire : un texte témoin}

Avant toute règle, lisons un court texte original qui rassemble toutes les phrases que nous allons traduire dans cette section. Observez la place des connecteurs, la marque de l'accompli \zh{了} et l'ordre des mots.

\begin{textebox}[从农村到城市 — Du village à la ville]
我的哥哥离开了农村，因为农村没有工作。\\
\emph{Wǒ de gēge líkāi le nóngcūn, yīnwèi nóngcūn méiyǒu gōngzuò.}\\
« Mon frère aîné a quitté la campagne parce qu'il n'y avait pas de travail à la campagne. »\\[4pt]
越来越多的年轻人搬到城市找工作。\\
\emph{Yuè lái yuè duō de niánqīngrén bān dào chéngshì zhǎo gōngzuò.}\\
« De plus en plus de jeunes déménagent en ville pour chercher du travail. »\\[4pt]
城市的生活很方便，但是房租很贵。\\
\emph{Chéngshì de shēnghuó hěn fāngbiàn, dànshì fángzū hěn guì.}\\
« La vie en ville est pratique, mais le loyer est très cher. »\\[4pt]
虽然城市很拥挤，但是机会很多。\\
\emph{Suīrán chéngshì hěn yōngjǐ, dànshì jīhuì hěn duō.}\\
« Bien que la ville soit bondée, il y a beaucoup d'opportunités. »\\[4pt]
为了家人，他每天都很努力。他想先找工作，再租房子，然后把父母接到城市来。\\
\emph{Wèile jiārén, tā měitiān dōu hěn nǔlì. Tā xiǎng xiān zhǎo gōngzuò, zài zū fángzi, ránhòu bǎ fùmǔ jiē dào chéngshì lái.}\\
« Pour sa famille, il travaille dur tous les jours. Il veut d'abord trouver un travail, puis louer un logement, ensuite faire venir ses parents en ville. »
\end{textebox}

\textbf{Glossaire du texte}

\begin{center}
\begin{tabularx}{\linewidth}{|X|X|X|X|}
\hline
\rowcolor{popblueL} Caractères & Pinyin & Nature & Traduction \\
\hline
离开 & líkāi & verbe & quitter, partir de \\
\hline
农村 & nóngcūn & nom & campagne, monde rural \\
\hline
搬 & bān & verbe & déménager, déplacer \\
\hline
方便 & fāngbiàn & adjectif & pratique, commode \\
\hline
房租 & fángzū & nom & loyer \\
\hline
拥挤 & yōngjǐ & adjectif & bondé, surpeuplé \\
\hline
机会 & jīhuì & nom & occasion, opportunité \\
\hline
努力 & nǔlì & adjectif/verbe & faire des efforts, travailleur \\
\hline
租 & zū & verbe & louer \\
\hline
接 & jiē & verbe & accueillir, aller chercher \\
\hline
\end{tabularx}
\end{center}

\textbf{Questions de compréhension (oral ou écrit) :}
\begin{enumerate}
\item \zh{哥哥为什么离开了农村？} \emph{Gēge wèishénme líkāi le nóngcūn?} — Pourquoi le frère a-t-il quitté la campagne ?
\item \zh{城市的生活怎么样？} \emph{Chéngshì de shēnghuó zěnmeyàng?} — Comment est la vie en ville ?
\item \zh{他想先做什么，再做什么？} \emph{Tā xiǎng xiān zuò shénme, zài zuò shénme?} — Que veut-il faire d'abord, et ensuite ?
\end{enumerate}

\courssub{La procédure de traduction en quatre étapes}

Traduire n'est pas remplacer des mots français par des mots chinois. C'est un travail en quatre étapes, valable pour le thème comme pour la version.

\begin{popfigure}
\begin{tikzpicture}[node distance=0.9cm, every node/.style={align=center}]
\node[draw=popblue, fill=popblueL, rounded corners, minimum width=2.6cm, minimum height=1.1cm, font=\bfseries] (lire) {1. Lire};
\node[draw=poporange, fill=poporangeL, rounded corners, minimum width=2.6cm, minimum height=1.1cm, font=\bfseries, right=of lire] (seg) {2. Segmenter};
\node[draw=popgreen, fill=popgreenL, rounded corners, minimum width=2.6cm, minimum height=1.1cm, font=\bfseries, right=of seg] (trad) {3. Traduire};
\node[draw=poppurple, fill=poppurpleL, rounded corners, minimum width=2.6cm, minimum height=1.1cm, font=\bfseries, right=of trad] (relire) {4. Relire};
\draw[-{Stealth}, very thick, popdark] (lire) -- (seg);
\draw[-{Stealth}, very thick, popdark] (seg) -- (trad);
\draw[-{Stealth}, very thick, popdark] (trad) -- (relire);
\node[below=0.25cm of lire, font=\small, text width=2.8cm] {comprendre le sens global};
\node[below=0.25cm of seg, font=\small, text width=2.8cm] {sujet / verbe / complément, connecteurs};
\node[below=0.25cm of trad, font=\small, text width=2.8cm] {ordre chinois, temps, 了};
\node[below=0.25cm of relire, font=\small, text width=2.8cm] {comparer avec le texte source};
\end{tikzpicture}
\legende{Procédure de traduction en quatre étapes : Lire → Segmenter → Traduire → Relire.}
\end{popfigure}

\begin{methode}[Traduire un thème (français → chinois) sans se tromper]
\begin{enumerate}
\item \textbf{Lire} la phrase française entière et la reformuler mentalement en « chinois simple » : qui ? fait quoi ? où ? quand ? pourquoi ?
\item \textbf{Segmenter} : repérer le sujet, le verbe, le complément et le connecteur (\zh{因为}, \zh{但是}, \zh{为了}…).
\item \textbf{Traduire} en respectant l'ordre chinois : Sujet + (temps/lieu) + Verbe + Complément. Ajouter \zh{了} si l'action est accomplie.
\item \textbf{Relire} en comparant au texte source : ai-je oublié une information ? un ton ? un caractère ?
\end{enumerate}
Ne traduisez \cle{jamais} mot à mot depuis le français : le français dit « il n'y avait pas de travail », le chinois dit simplement \zh{没有工作} « ne pas avoir de travail ».
\end{methode}

\courssub{Thème guidé : trois phrases pas à pas}

\textbf{Thème 1.} « Mon frère aîné a quitté le village parce qu'il n'y avait pas de travail. »

Étape 2 — segmentation : sujet = \zh{我的哥哥} ; verbe = \zh{离开} (accompli → \zh{了}) ; complément = \zh{农村} ; connecteur de cause = \zh{因为}.

Étape 3 — traduction : \zh{我的哥哥离开了农村，因为农村没有工作。} \emph{Wǒ de gēge líkāi le nóngcūn, yīnwèi nóngcūn méiyǒu gōngzuò.}

Remarquez : le français « il n'y avait pas » (imparfait) devient simplement \zh{没有} : le chinois ne conjugue pas, le contexte suffit. Notez aussi que la négation de \zh{有} « avoir » est toujours \zh{没有}, jamais \zh{不有}.

\textbf{Thème 2.} « La vie en ville est pratique, mais le loyer est très cher. »

Segmentation : sujet 1 = \zh{城市的生活} (la vie \emph{de} la ville → \zh{的}) ; adjectif = \zh{方便} ; connecteur d'opposition = \zh{但是} ; sujet 2 = \zh{房租}.

Traduction : \zh{城市的生活很方便，但是房租很贵。} \emph{Chéngshì de shēnghuó hěn fāngbiàn, dànshì fángzū hěn guì.}

Piège : en chinois, l'adjectif attribut ne prend pas le verbe « être » : on dit \zh{房租很贵}, jamais \zh{房租是很贵}. Le mot \zh{很}, littéralement « très », sert surtout de lien naturel entre le sujet et l'adjectif.

\textbf{Thème 3.} « Pour sa famille, il travaille beaucoup tous les jours. »

Segmentation : complément de but = \zh{为了家人} (placé \emph{avant} le sujet, suivi d'une virgule) ; sujet = \zh{他} ; temps = \zh{每天} ; adverbe = \zh{都} ; verbe = \zh{努力}.

Traduction : \zh{为了家人，他每天都很努力。} \emph{Wèile jiārén, tā měitiān dōu hěn nǔlì.}

\begin{attention}[L'accompli : n'oubliez pas 了]
Dans une phrase qui raconte un événement accompli, le chinois marque l'accompli avec \zh{了} juste après le verbe : \zh{离开了} (a quitté), \zh{搬到了} (a déménagé), \zh{找到了} (a trouvé). Les francophones l'oublient souvent parce que le français conjugue le verbe, pas le chinois. Attention : \zh{了} ne s'emploie pas avec les états habituels ni après \zh{没有} (\zh{农村没有工作}, sans \zh{了}).
\end{attention}

\begin{exemplebox}[Exemples traduits supplémentaires]
\zh{他想先找工作，再租房子。} \emph{Tā xiǎng xiān zhǎo gōngzuò, zài zū fángzi.} « Il veut d'abord trouver un travail, puis louer un logement. » (structure \zh{先……再……} = d'abord… puis…)\\
\zh{弟弟搬到了杜阿拉，因为那里机会很多。} \emph{Dìdi bān dào le Dù'ālā, yīnwèi nàli jīhuì hěn duō.} « Le petit frère a déménagé à Douala parce qu'il y a beaucoup d'opportunités là-bas. »\\
\zh{虽然房租很贵，但是他不想回农村。} \emph{Suīrán fángzū hěn guì, dànshì tā bù xiǎng huí nóngcūn.} « Bien que le loyer soit très cher, il ne veut pas retourner à la campagne. »
\end{exemplebox}

\courssub{Version guidée : du chinois vers le français}

\textbf{Version 1.} \zh{越来越多的年轻人搬到城市找工作。}

Segmentation : \zh{越来越多的年轻人} = sujet (« de plus en plus de jeunes ») ; \zh{搬到城市} = verbe + direction (« déménagent en ville » : \zh{到} marque la destination) ; \zh{找工作} = but (« pour chercher du travail »).

Traduction : « De plus en plus de jeunes déménagent en ville pour chercher du travail. »

\textbf{Version 2.} \zh{虽然城市很拥挤，但是机会很多。}

Segmentation : \zh{虽然} = « bien que » ; \zh{城市很拥挤} = « la ville est bondée » ; \zh{但是} = « mais » ; \zh{机会很多} = « il y a beaucoup d'opportunités ».

Traduction : « Bien que la ville soit bondée, il y a beaucoup d'opportunités. » En version, un seul des deux connecteurs suffit en français (« bien que… » ou « …, mais… »), alors qu'en chinois \zh{虽然……但是……} fonctionnent en paire.

\begin{propriete}[Ponctuation chinoise : virgule et point pleine largeur]
En chinois, la virgule est \zh{，} et le point est \zh{。} (un petit cercle, pas un point). Ce sont des signes \emph{pleine largeur} : ils occupent la place d'un caractère et créent un espace visuel naturel. Il n'y a \textbf{pas de majuscules} en chinois, même en début de phrase : seuls les noms propres se distinguent éventuellement dans le pinyin (\emph{Kāmàilóng}). On n'écrit donc jamais de majuscule après \zh{。}.
\end{propriete}

\begin{savaistu}[L'exode rural, un phénomène camerounais et chinois]
Comme le Cameroun, la Chine connaît un important mouvement des campagnes vers les villes : des millions de travailleurs migrants, appelés \zh{农民工} \emph{nóngmíngōng} (« ouvriers paysans »), ont quitté leur village pour travailler dans les grandes villes comme Pékin ou Shanghai. Beaucoup envoient de l'argent à leur famille restée au village, exactement comme un jeune de Douala ou de Yaoundé peut aider ses parents restés au village. Ce parallèle vous aidera à rédiger vos propres textes sur ce thème.
\end{savaistu}

\courssub{Tableau récapitulatif et erreurs typiques}

\begin{center}
\begin{tabularx}{\linewidth}{|X|X|X|}
\hline
\rowcolor{popblueL} Structure & Exemple en caractères + pinyin & Traduction \\
\hline
因为 + cause & 因为农村没有工作 (yīnwèi nóngcūn méiyǒu gōngzuò) & parce qu'il n'y a pas de travail à la campagne \\
\hline
但是 + opposition & 但是房租很贵 (dànshì fángzū hěn guì) & mais le loyer est très cher \\
\hline
为了 + but & 为了家人 (wèile jiārén) & pour sa famille \\
\hline
越来越 + adj. & 越来越多 (yuè lái yuè duō) & de plus en plus (nombreux) \\
\hline
虽然……但是…… & 虽然城市很拥挤，但是机会很多 & bien que la ville soit bondée, il y a beaucoup d'opportunités \\
\hline
先……再…… & 先找工作，再租房子 & d'abord trouver un travail, puis louer un logement \\
\hline
\end{tabularx}
\end{center}

\begin{attention}[Erreurs typiques à corriger en collectif]
\begin{itemize}
\item \textbf{Ordre des mots} : le complément de temps se place avant le verbe. On dit \zh{他每天都很努力}, jamais \zh{他很努力每天}.
\item \textbf{Oubli de \zh{了}} : « Mon frère a quitté le village » exige \zh{离开了}, pas seulement \zh{离开}.
\item \textbf{Mauvais connecteur} : \zh{因为} introduit la cause, \zh{所以} la conséquence ; ne pas les inverser. « Il est parti \emph{parce qu}'il n'y avait pas de travail » = \zh{因为}, pas \zh{所以}.
\item \textbf{Verbe « être » en trop} : \zh{房租很贵} suffit ; pas de \zh{是} devant un adjectif.
\item \textbf{Négation de \zh{有}} : on dit \zh{没有工作}, jamais \zh{不有工作}.
\end{itemize}
\end{attention}

\begin{exoresolu}[Thème corrigé : phrase à traduire]
Énoncé. Traduisez en chinois : « Ma sœur aînée a déménagé à Yaoundé parce que la ville offre plus d'opportunités. »
\tcblower
Solution détaillée. Étape 1 : qui ? ma sœur aînée → \zh{我的姐姐}. Fait quoi ? a déménagé (accompli) → \zh{搬到了}. Où ? à Yaoundé → \zh{雅温得}. Pourquoi ? parce que la ville a plus d'opportunités → \zh{因为城市的机会更多}. Étape 3 : on assemble dans l'ordre chinois. Traduction : \zh{我的姐姐搬到了雅温得，因为城市的机会更多。} \emph{Wǒ de jiějie bān dào le Yǎwēndé, yīnwèi chéngshì de jīhuì gèng duō.} Vérification : \zh{了} présent (action accomplie), \zh{的} entre \zh{城市} et \zh{机会}, virgule pleine largeur \zh{，}, point \zh{。}.
\end{exoresolu}

\begin{exercice}[À vous de traduire]
\begin{enumerate}
\item Thème : « De plus en plus de jeunes quittent le village pour travailler en ville. »
\item Thème : « Bien que le loyer soit cher, il veut rester à Douala. »
\item Version : \zh{为了家人，哥哥每天都很努力。}
\item Version : \zh{虽然农村很安静，但是工作很少。}
\end{enumerate}
\end{exercice}

\begin{corrige}[Exercice : Thème et version]
\begin{enumerate}
\item \zh{越来越多的年轻人离开农村去城市工作。} \emph{Yuè lái yuè duō de niánqīngrén líkāi nóngcūn qù chéngshì gōngzuò.} (habitude présente : \zh{了} non obligatoire ici ; \zh{去} « aller » introduit le but.)
\item \zh{虽然房租很贵，但是他想留在杜阿拉。} \emph{Suīrán fángzū hěn guì, dànshì tā xiǎng liú zài Dù'ālā.}
\item « Pour sa famille, le grand frère travaille dur tous les jours. »
\item « Bien que la campagne soit calme, il y a peu de travail. »
\end{enumerate}
\end{corrige}

\begin{aretenir}[L'essentiel de la section]
Pour réussir thème et version : lire, segmenter (sujet-verbe-complément + connecteur), traduire dans l'ordre chinois, relire en comparant au source. Trois réflexes : marquer l'accompli avec \zh{了}, choisir le bon connecteur (\zh{因为} / \zh{但是} / \zh{为了} / \zh{虽然……但是……}), et utiliser la ponctuation pleine largeur \zh{，} et \zh{。} sans majuscules.
\end{aretenir}

\coursec{Production écrite et évaluation}

Dans la section précédente, vous avez appris à traduire des phrases simples sur l'exode rural, du français vers le chinois (thème) et du chinois vers le français (version). Vous disposez maintenant de tout le matériel nécessaire : le lexique du thème (\zh{农村} campagne, \zh{城市} ville, \zh{打工} travailler comme migrant, \zh{离开} quitter…), les connecteurs (\zh{因为…所以…}, \zh{但是}, \zh{为了}) et les règles d'écriture des caractères complexes. Il est temps de rassembler tout cela dans une \cle{production écrite complète}, exactement comme on vous le demandera à l'examen.

\courssub{La tâche finale : sujet et contraintes}

\begin{definition}[Sujet de production écrite]
Rédigez en chinois un texte de \cle{6 à 10 phrases} intitulé \zh{我的村子和我哥哥的城市} (\emph{Wǒ de cūnzi hé wǒ gēge de chéngshì.}) « Mon village et la ville de mon frère ». Vous y racontez la vie dans votre village, le départ de votre frère (ou d'un proche) pour la ville, les raisons de ce départ et votre point de vue.
\end{definition}

\begin{propriete}[Contraintes obligatoires]
Votre texte doit respecter \textbf{toutes} les contraintes suivantes :
\begin{enumerate}
\item Utiliser au moins \cle{3 connecteurs} parmi : \zh{因为…所以…} (parce que… donc…), \zh{但是} (mais), \zh{为了} (afin de, pour).
\item Employer au moins \cle{5 mots du glossaire} de la leçon (voir tableau ci-dessous).
\item Écrire les caractères \cle{sans erreur de traits} : chaque trait compte, surtout dans les caractères complexes du thème (\zh{村}, \zh{城}, \zh{离}, \zh{希}, \zh{望}).
\item Respecter la longueur : entre 6 et 10 phrases complètes (sujet + verbe + ponctuation chinoise 。).
\end{enumerate}
\end{propriete}

\begin{center}
\begin{tabularx}{\linewidth}{|X|X|X|X|}\hline
\rowcolor{popblueL} Caractères & Pinyin & Nature & Traduction \\ \hline
农村 & nóngcūn & nom & la campagne, le village \\ \hline
城市 & chéngshì & nom & la ville \\ \hline
村子 & cūnzi & nom & le village \\ \hline
离开 & líkāi & verbe & quitter, partir de \\ \hline
打工 & dǎgōng & verbe & travailler (comme migrant) \\ \hline
工作 & gōngzuò & verbe / nom & travailler ; le travail \\ \hline
希望 & xīwàng & verbe / nom & espérer ; l'espoir \\ \hline
生活 & shēnghuó & nom / verbe & la vie ; vivre \\ \hline
越来越 & yuè lái yuè & structure & de plus en plus \\ \hline
机会 & jīhuì & nom & l'occasion, la chance \\ \hline
\end{tabularx}
\end{center}

\courssub{Observer avant d'écrire : deux phrases modèles}

Avant de vous lancer, observez ces deux phrases qui peuvent encadrer votre texte.

\begin{exemplebox}[Phrase d'ouverture et phrase de conclusion]
\textbf{Ouverture :} \zh{我来自一个小农村。}\\
\emph{Wǒ láizì yí ge xiǎo nóngcūn.}\\
« Je viens d'un petit village. »\\[4pt]
\textbf{Conclusion :} \zh{我希望农村和城市都越来越好。}\\
\emph{Wǒ xīwàng nóngcūn hé chéngshì dōu yuè lái yuè hǎo.}\\
« J'espère que la campagne et la ville iront de mieux en mieux. »
\end{exemplebox}

\begin{propriete}[Structure \zh{越来越 + adjectif}]
La structure \zh{越来越 + adjectif} signifie « de plus en plus + adjectif » (ex. \zh{越来越好}, \zh{越来越忙}). Elle \textbf{ne prend jamais \zh{很} devant l'adjectif} : on dit \zh{越来越好} et non \zh{越来越很好}. Le sujet précède \zh{越来越} : \zh{城市越来越大} (La ville devient de plus en plus grande).
\end{propriete}

\courssub{Un texte modèle complet}

Lisez d'abord ce texte modèle (environ 230 caractères). Repérez les connecteurs, les mots du glossaire et la progression des idées.

\begin{textebox}[Modèle : 我的村子和我哥哥的城市]
我来自一个小农村。我的村子在雅温得附近，那里很安静，空气也很好。\\
\emph{Wǒ láizì yí ge xiǎo nóngcūn. Wǒ de cūnzi zài Yǎwēndé fùjìn, nàlǐ hěn ānjìng, kōngqì yě hěn hǎo.}\\
« Je viens d'un petit village. Mon village est près de Yaoundé ; c'est un endroit très calme et l'air y est bon. »\\[4pt]
但是，村子里工作机会不多，所以很多年轻人离开了农村。\\
\emph{Dànshì, cūnzi lǐ gōngzuò jīhuì bù duō, suǒyǐ hěn duō niánqīngrén líkāi le nóngcūn.}\\
« Mais il n'y a pas beaucoup d'occasions de travail au village, donc beaucoup de jeunes ont quitté la campagne. »\\[4pt]
三年前，我哥哥也离开了村子。为了找到好工作，他去了杜阿拉打工。\\
\emph{Sān nián qián, wǒ gēge yě líkāi le cūnzi. Wèile zhǎodào hǎo gōngzuò, tā qù le Dùālā dǎgōng.}\\
« Il y a trois ans, mon frère a lui aussi quitté le village. Afin de trouver un bon travail, il est parti travailler à Douala. »\\[4pt]
因为城市里有更多的工作，所以他现在生活得不错。他每个月都给家里打电话。\\
\emph{Yīnwèi chéngshì lǐ yǒu gèng duō de gōngzuò, suǒyǐ tā xiànzài shēnghuó de bùcuò. Tā měi ge yuè dōu gěi jiālǐ dǎ diànhuà.}\\
« Parce qu'il y a plus de travail en ville, il vit plutôt bien maintenant. Il téléphone à la famille chaque mois. »\\[4pt]
我很想念他，但是我也很喜欢我的村子。我希望农村和城市都越来越好。\\
\emph{Wǒ hěn xiǎngniàn tā, dànshì wǒ yě hěn xǐhuan wǒ de cūnzi. Wǒ xīwàng nóngcūn hé chéngshì dōu yuè lái yuè hǎo.}\\
« Il me manque beaucoup, mais j'aime aussi beaucoup mon village. J'espère que la campagne et la ville iront de mieux en mieux. »
\end{textebox}

\textbf{Questions de compréhension :}
\begin{enumerate}
\item Où se trouve le village du narrateur ? (\emph{Près de Yaoundé.})
\item Pourquoi les jeunes quittent-ils le village ? (\emph{Parce qu'il y a peu d'occasions de travail.})
\item Pourquoi le frère est-il parti à Douala ? (\emph{Pour trouver un bon travail.})
\item Combien de connecteurs le texte contient-il ? Lesquels ? (\emph{Trois : 但是, 为了, 因为…所以….})
\end{enumerate}

\courssub{La démarche en quatre étapes}

\begin{methode}[Du plan français à la copie au propre]
\begin{enumerate}
\item \textbf{Plan en français (5 minutes).} Notez vos idées en français : situation du village → problème (peu de travail) → départ du frère → sa vie en ville → vos sentiments → espoir final. Prévoyez où placer vos 3 connecteurs.
\item \textbf{Brouillon en chinois (20 minutes).} Rédigez phrase par phrase, avec des structures simples et sûres : Sujet + Verbe + Complément. N'écrivez que des caractères que vous savez tracer.
\item \textbf{Vérification des caractères complexes (5 minutes).} Relisez chaque caractère du thème trait par trait :
  \begin{itemize}
  \item \zh{村} (cūn) : clé du bois \zh{木} + \zh{寸} (pouce/mesure). Ne pas confondre avec \zh{材} (cái, clé bois + \zh{才}).
  \item \zh{城} (chéng) : clé de la terre \zh{土} + \zh{chéng} (accomplir).
  \item \zh{离} (lí) : caractère simplifié (forme traditionnelle \zh{離}), radical \zh{禸} (liù). Attention à ne pas ajouter de traits superflus.
  \item \zh{希} (xī) : clé du tissu \zh{巾} + \zh{禾} (grain).
  \item \zh{望} (wàng) : clé de la chair/lune \zh{月} + phonétique \zh{望} (composé de \zh{亡} au-dessus de \zh{王}).
  \end{itemize}
  Comptez les traits si vous avez un doute.
\item \textbf{Copie au propre (10 minutes).} Recopiez lentement, en caractères bien formés, réguliers et espacés. La ponctuation chinoise (，。？！) occupe une case entière.
\end{enumerate}
\end{methode}

\begin{methode}[Les 5 dernières minutes : la relecture]
Réservez toujours les \cle{cinq dernières minutes} à la relecture, jamais à de nouvelles phrases. Procédez ainsi :
\begin{enumerate}
\item Vérifiez chaque caractère complexe \textbf{trait par trait} (un trait manquant = une faute).
\item Vérifiez l'ordre des mots : le complément de lieu avec \zh{在} se place \textbf{avant} le verbe (\zh{我在村子生活} et non \zh{我生活在村子} à ce niveau).
\item Vérifiez que chaque phrase se termine par 。 et que vos 3 connecteurs sont bien présents.
\end{enumerate}
\end{methode}

\begin{attention}[Un trait peut changer le sens / Pièges grammaticaux]
\begin{itemize}
\item \textbf{Caractères :} \zh{村} (cūn, village) et \zh{材} (cái, matériau) ne diffèrent que par le composant de droite (\zh{寸} vs \zh{才}) ; \zh{土} (tǔ, terre, trait horizontal du bas plus long) et \zh{士} (shì, lettré, trait horizontal du haut plus long). À l'examen, \cle{la lisibilité compte autant que la grammaire} : un caractère illisible ou ambigu est compté faux.
\item \textbf{Grammaire :} Pour dire « la vie se passe bien », on utilise couramment \zh{过得不错} (guò de bùcuò). \zh{生活得不错} (shēnghuó de bùcuò) est possible mais insiste sur les \emph{conditions de vie}. Préférez \zh{过得好} à l'oral et \zh{生活得好} à l'écrit formel.
\item \textbf{Connecteur \zh{为了} :} Il introduit le \textbf{but} et se place \textbf{avant} le verbe d'action principal (souvent en début de phrase ou après le sujet) : \zh{为了找工作，他去了城市} et non \zh{他去了城市，为了找工作}.
\item \textbf{Paire \zh{因为…所以…} :} Les deux mots sont \textbf{indissociables} dans une phrase complexe. N'oubliez jamais \zh{所以} dans la seconde proposition.
\end{itemize}
\end{attention}

\courssub{Exercice résolu : corriger une copie d'élève}

\begin{exoresolu}[Corriger les erreurs d'un brouillon]
\textbf{Énoncé.} Voici trois phrases du brouillon d'un élève. Repérez et corrigez les erreurs (caractères, grammaire, connecteur).\\
1. \zh{我哥哥离开农村，为了找好工作。} (ordre des idées maladroit : le but est exprimé après l'action)\\
2. \zh{因为城市里有机会，他去了城市。但是} (connecteur incomplet, phrase inachevée)\\
3. L'élève a écrit \zh{材子} au lieu de \zh{村子}.
\tcblower
\textbf{Solution détaillée.}\\
1. Avec \zh{为了}, le but se place normalement \textbf{avant} l'action principale : \zh{为了找到好工作，我哥哥离开了农村。} \emph{Wèile zhǎodào hǎo gōngzuò, wǒ gēge líkāi le nóngcūn.} « Afin de trouver un bon travail, mon frère a quitté la campagne. » (Ajout de \zh{到} pour marquer le résultat de l'action \zh{找}).\\
2. La structure \zh{因为…所以…} fonctionne \textbf{en paire} : \zh{因为城市里有机会，所以他去了城市。} \emph{Yīnwèi chéngshì lǐ yǒu jīhuì, suǒyǐ tā qù le chéngshì.} « Parce qu'il y a des occasions en ville, il y est allé. » Le \zh{但是} isolé en fin de phrase est une faute : \zh{但是} introduit toujours une proposition complète (ex. \zh{但是他不想走}).\\
3. \zh{材} (cái, matériau) et \zh{村} (cūn, village) partagent la clé du bois \zh{木}, mais le second composant diffère : \zh{寸} (cùn) pour \zh{村}, \zh{才} (cái) pour \zh{材}. La forme correcte est \zh{村子} (cūnzi, le village). Un seul trait de différence change complètement le sens.
\end{exoresolu}

\courssub{La grille d'évaluation}

Votre texte sera noté sur 20 points selon quatre critères. Étudiez bien cette grille : elle vous indique exactement ce que le correcteur attend.

\begin{popfigure}
\begin{tikzpicture}[x=1cm,y=1cm]
\draw[fill=popblue,draw=popblue] (0,4) rectangle (10,4.8);
\node[text=white,font=\bfseries] at (5,4.4) {Grille d'évaluation — Production écrite (20 points)};
\draw[fill=popblueL,draw=popline] (0,3.2) rectangle (6,4);
\draw[fill=popblueL,draw=popline] (6,3.2) rectangle (8,4);
\draw[fill=popblueL,draw=popline] (8,3.2) rectangle (10,4);
\node[font=\bfseries] at (3,3.6) {Critère};
\node[font=\bfseries] at (7,3.6) {Barème};
\node[font=\bfseries] at (9,3.6) {Points};
\draw[fill=white,draw=popline] (0,2.4) rectangle (6,3.2);
\draw[fill=white,draw=popline] (6,2.4) rectangle (8,3.2);
\draw[fill=white,draw=popline] (8,2.4) rectangle (10,3.2);
\node[align=left,text width=5.4cm] at (3,2.8) {Exactitude des caractères (traits, radicaux)};
\node[align=center,text width=1.6cm] at (7,2.8) {0 à 5};
\node at (9,2.8) {/ 5};
\draw[fill=popcream,draw=popline] (0,1.6) rectangle (6,2.4);
\draw[fill=popcream,draw=popline] (6,1.6) rectangle (8,2.4);
\draw[fill=popcream,draw=popline] (8,1.6) rectangle (10,2.4);
\node[align=left,text width=5.4cm] at (3,2) {Grammaire et ordre des mots};
\node[align=center,text width=1.6cm] at (7,2) {0 à 5};
\node at (9,2) {/ 5};
\draw[fill=white,draw=popline] (0,0.8) rectangle (6,1.6);
\draw[fill=white,draw=popline] (6,0.8) rectangle (8,1.6);
\draw[fill=white,draw=popline] (8,0.8) rectangle (10,1.6);
\node[align=left,text width=5.4cm] at (3,1.2) {Connecteurs et cohérence (au moins 3)};
\node[align=center,text width=1.6cm] at (7,1.2) {0 à 5};
\node at (9,1.2) {/ 5};
\draw[fill=popcream,draw=popline] (0,0) rectangle (6,0.8);
\draw[fill=popcream,draw=popline] (6,0) rectangle (8,0.8);
\draw[fill=popcream,draw=popline] (8,0) rectangle (10,0.8);
\node[align=left,text width=5.4cm] at (3,0.4) {Respect du sujet et longueur (6 à 10 phrases)};
\node[align=center,text width=1.6cm] at (7,0.4) {0 à 5};
\node at (9,0.4) {/ 5};
\draw[fill=poporangeL,draw=popline] (0,-0.8) rectangle (6,0);
\draw[fill=poporangeL,draw=popline] (6,-0.8) rectangle (8,0);
\draw[fill=poporangeL,draw=popline] (8,-0.8) rectangle (10,0);
\node[font=\bfseries] at (3,-0.4) {Total};
\node at (7,-0.4) {};
\node[font=\bfseries] at (9,-0.4) {/ 20};
\end{tikzpicture}
\legende{Grille d'évaluation de la production écrite : quatre critères de 5 points chacun, pour un total de 20.}
\end{popfigure}

\begin{center}
\begin{tabularx}{\linewidth}{|X|X|X|}\hline
\rowcolor{popblueL} Critère & Ce qui est récompensé & Erreur qui fait perdre des points \\ \hline
Caractères (5 pts) & traits corrects, caractères lisibles & trait manquant, confusion 村/材 \\ \hline
Grammaire (5 pts) & ordre S + V + C, place de 在 & verbe avant le complément de lieu \\ \hline
Connecteurs (5 pts) & 因为…所以…, 但是, 为了 bien placés & 因为 sans 所以, 但是 en fin de phrase \\ \hline
Sujet et longueur (5 pts) & 6 à 10 phrases sur le thème & hors sujet, moins de 6 phrases \\ \hline
\end{tabularx}
\end{center}

\begin{aretenir}[Rappel express : vos 3 connecteurs obligatoires]
\begin{center}
\begin{tabularx}{\linewidth}{|X|X|X|}\hline
\rowcolor{popblueL} Connecteur & Structure & Exemple \\ \hline
\zh{因为…所以…} & \zh{因为} + Cause , \zh{所以} + Conséquence & \zh{因为村子里工作少，所以哥哥走了。} \\ \hline
\zh{但是} & Proposition 1 , \zh{但是} + Proposition 2 (opposition) & \zh{村子很美，但是工作少。} \\ \hline
\zh{为了} & \zh{为了} + But , Sujet + Verbe (action) & \zh{为了赚钱，他在城里打工。} \\ \hline
\end{tabularx}
\end{center}
\end{aretenir}

\courssub{Auto-évaluation et correction croisée}

\begin{experience}[Évaluer comme un examinateur]
\begin{enumerate}
\item \textbf{Auto-évaluation (5 minutes).} Relisez votre propre copie avec la grille : attribuez-vous une note sur 5 pour chaque critère et justifiez-la en une phrase en français.
\item \textbf{Correction croisée (10 minutes).} Échangez votre copie avec un voisin. Notez sa production avec la grille, entourez les caractères douteux et soulignez les connecteurs trouvés.
\item \textbf{Retour et amélioration (5 minutes).} Récupérez votre copie, lisez les remarques de votre pair, puis réécrivez les deux phrases les plus faibles.
\end{enumerate}
Cette activité vous entraîne à lire le chinois avec l'œil du correcteur : c'est le meilleur moyen de progresser avant l'examen.
\end{experience}

\begin{exercice}[À vous d'écrire]
Rédigez maintenant votre propre texte \zh{我的村子和我哥哥的城市} en suivant la démarche en quatre étapes. Contraintes : 6 à 10 phrases, au moins 3 connecteurs, au moins 5 mots du glossaire, caractères soignés. Temps conseillé : 40 minutes.
\end{exercice}

\begin{corrige}[Exercice : production écrite]
Il n'existe pas une seule bonne réponse, mais une copie réussie doit contenir : une ouverture de type \zh{我来自一个小农村。} ; au moins une phrase avec \zh{因为…所以…} (ex. \zh{因为村子里工作不多，所以哥哥去了城市。}) ; une opposition avec \zh{但是} ; un but exprimé avec \zh{为了} placé avant l'action ; une conclusion avec \zh{希望} et \zh{越来越}. Vérifiez enfin que les caractères \zh{村}, \zh{城}, \zh{离}, \zh{希}, \zh{望} sont tracés correctement et que la longueur est respectée (6 à 10 phrases).
\end{corrige}

\courssub{Complément utile pour l'examen}

Au baccalauréat, la production écrite de chinois est notée aussi sur la \cle{lisibilité des caractères}. Le correcteur doit pouvoir identifier chaque caractère sans hésitation. Entraînez-vous donc régulièrement à écrire \textbf{lentement et régulièrement} : des caractères de taille constante, des traits nets, un espacement égal. Dix minutes d'écriture soignée par jour valent mieux qu'une heure d'écriture rapide la veille de l'examen.

\begin{savaistu}[L'écriture soignée, une valeur partagée]
Au Cameroun comme en Chine, une belle écriture est un signe de sérieux et de respect du lecteur. En Chine, la calligraphie, \zh{书法} (\emph{shūfǎ}, littéralement « la méthode de l'écriture »), est un art millénaire pratiqué au pinceau et à l'encre ; les grands maîtres calligraphes y sont aussi célèbres que les peintres. Même à l'ère du téléphone portable, savoir tracer un caractère correctement à la main reste une compétence très valorisée à l'école chinoise — et dans votre examen camerounais !
\end{savaistu}

\newpage

\coursec{Activité d'intégration}

\coursec{Leçon 26 — Expression écrite : caractères complexes liés à l'exode rural}

\courssub{Objectifs de la leçon}

À l'issue de cette leçon (2 h), tu seras capable de :

\begin{itemize}
\item \textbf{Comprendre} un modèle de lettre personnelle sur l'exode rural (structure, formules de politesse, connecteurs).
\item \textbf{Mobiliser} le lexique du thème (minimum 6 mots : \cle{搬家}, \cle{离开}, \cle{农村}, \cle{城市}, \cle{挤}, \cle{租}, \cle{找工作}, \cle{希望}, \cle{困难}, \cle{机会}).
\item \textbf{Écrire correctement} les quatre caractères complexes ciblés : \cle{搬}, \cle{离}, \cle{挤}, \cle{租} (radicaux, ordre des traits, caractères proches).
\item \textbf{Structurer} un texte cohérent (8--10 phrases) en employant au moins 3 connecteurs logiques : \cle{因为…所以…}, \cle{虽然…但是…}, \cle{先…然后…}, \cle{以后}.
\item \textbf{Traduire} (thème/version) des phrases simples mobilisant ces structures et ce vocabulaire.
\item \textbf{Produire} une lettre personnelle complète (en-tête, corps, formule de fin, signature) et \textbf{évaluer} celle d'un pair à l'aide d'une grille critériée.
\end{itemize}

\courssub{Situation de communication}

\textbf{Contexte :} Tu t'appelles \textbf{André Mbarga}, 17 ans. Le mois dernier, tu as quitté ton village près d'\zh{埃博洛瓦} (Èbōluówǎ, Ebolowa) pour rejoindre \zh{雅温得} (Yǎwēndé, Yaoundé), chez ton oncle. Tu cherches du travail et souhaites poursuivre tes études. La vie urbaine est coûteuse, les transports sont saturés, tu partages une petite chambre. Pourtant, tu gardes espoir. Ce dimanche, tu écris à ta famille restée au village pour la rassurer, décrire ton quotidien et exposer tes projets.

\courssub{Modèle de texte commenté}

Lis attentivement ce début de lettre. Il sert de \cle{modèle de forme} (en-tête, appel, formules de politesse, signature).

\begin{textebox}[Modèle de début de lettre — André à sa mère]
亲爱的妈妈：\\
Qīn'ài de māma :\\
« Chère maman, »\\[4pt]
您好！我离开家已经一个月了，很想您。\\
Nín hǎo ! Wǒ líkāi jiā yǐjīng yí ge yuè le, hěn xiǎng nín.\\
« Bonjour ! Cela fait déjà un mois que j'ai quitté la maison, vous me manquez beaucoup. »\\[4pt]
我上个月搬到了雅温得，现在住在叔叔家。\\
Wǒ shàng ge yuè bāndào le Yǎwēndé, xiànzài zhù zài shūshu jiā.\\
« Le mois dernier, j'ai déménagé à Yaoundé, maintenant j'habite chez mon oncle. »\\[4pt]
城市很大，人很多，公共汽车很挤。\\
Chéngshì hěn dà, rén hěn duō, gōnggòng qìchē hěn jǐ.\\
« La ville est grande, il y a beaucoup de monde, les bus sont bondés. »
\end{textebox}

\begin{aretenir}[Analyse du modèle]
\begin{itemize}
\item \textbf{Formule d'appel} : \zh{亲爱的妈妈：} (Qīn'ài de māma :) — « Chère maman, ». \zh{亲爱的} (qīn'ài de) marque l'affection ; les deux points chinois \zh{：} sont obligatoires.
\item \textbf{Formule de politesse} : \zh{您好！} (Nín hǎo !) — \zh{您} (nín) est la forme polie de « vous » (à utiliser avec les aînés).
\item \textbf{Durée écoulée} : \zh{离开家已经一个月了} (líkāi jiā yǐjīng yí ge yuè le). Structure : \cle{Sujet + 离开 + Lieu + 已经 + Durée + 了}. \zh{已经…了} (yǐjīng… le) indique une action commencée dans le passé et qui se poursuit (ou dont l'état resultant dure).
\item \textbf{Verbe \zh{搬} (bān)} : \zh{搬到了} (bāndào le) = « déménager vers / arriver à destination en déménageant ». Complément de direction \zh{到} (dào) + lieu.
\item \textbf{Ordre des mots} : Complément de temps \zh{上个月} (shàng ge yuè) \textbf{avant} le verbe (\zh{我上个月搬到了…}).
\item \textbf{Description} : Phrases nominales parataxiques (sujet + adjectif) : \zh{城市很大，人很多} — style télégraphique courant à l'oral et dans l'écrit personnel.
\end{itemize}
\end{aretenir}

\courssub{Lexique écrit du thème}

\begin{center}
\begin{tabularx}{\linewidth}{|X|X|X|X|}
\hline
\rowcolor{popblueL} Caractères & Pinyin & Nature & Traduction \\
\hline
搬家 & bān jiā & verbe (VO) & déménager (litt. « déplacer la maison ») \\
\hline
离开 & lí kāi & verbe & quitter, partir de \\
\hline
农村 & nóng cūn & nom & campagne, village, zone rurale \\
\hline
城市 & chéng shì & nom & ville, zone urbaine \\
\hline
挤 & jǐ & adjectif & bondé, serré, plein à craquer \\
\hline
租 & zū & verbe & louer (payer pour utiliser) \\
\hline
找工作 & zhǎo gōng zuò & locution verbale & chercher du travail / un emploi \\
\hline
希望 & xī wàng & verbe / nom & espérer / espoir \\
\hline
困难 & kùn nan & nom / adjectif & difficulté / difficile \\
\hline
机会 & jī huì & nom & occasion, opportunité, chance \\
\hline
生活 & shēng huó & nom / verbe & vie, quotidien / vivre \\
\hline
继续 & jì xù & verbe & continuer \\
\hline
帮助 & bāng zhù & verbe & aider, aider à \\
\hline
身体健康 & shēn tǐ jiàn kāng & expression & bonne santé (formule de vœux) \\
\hline
\end{tabularx}
\end{center}

\begin{savaistu}[Exode rural : Cameroun \textbullet{} Chine]
Au \textbf{Cameroun}, l'exode rural (\zh{农村人口外流} nóngcūn rénkǒu wàiliú) alimente Yaoundé et Douala. Les jeunes quittent les zones de culture (cacao, café, banane plantain) pour l'informel urbain (\zh{小商贩} xiǎo shāngfàn, \zh{摩托车司机} mótuōchē sījī).\\
En \textbf{Chine}, le phénomène est massif (\zh{农民工} nóngmíngōng, « ouvriers paysans ») depuis les années 1980. Le système de \zh{户口} (hùkǒu, registre de résidence) crée une dualité : les migrants ruraux travaillent en ville (Bâtiment, usines, livraison \zh{外卖} wàimài) sans accès complet aux services urbains (école, santé). La Chine mène une politique d'\zh{城镇化} (chéngzhènhuà, urbanisation) et de \zh{乡村振兴} (xiāngcūn zhènxīng, revitalisation rurale) pour rééquilibrer.
\end{savaistu}

\courssub{Écriture des caractères complexes}

Les quatre caractères ci-dessous sont \cle{caractères phono-idéographiques} (形声字) : une clé (radical) donne le sens, une partie phonétique suggère le son. Maîtrise leur ordre des traits (règle générale : haut $\rightarrow$ bas, gauche $\rightarrow$ droite, horizontal avant vertical, trait central avant ailes).

\begin{propriete}[Fiches d'identité des 4 caractères ciblés]
\begin{center}
\begin{tabularx}{\linewidth}{|X|X|X|X|X|}
\hline
\rowcolor{popblueL} Caractère & Pinyin & Radical (clé) & Décomposition & Ordre des traits (résumé) \\
\hline
\zh{搬} & bān & \zh{扌} (shǒu, main) & \zh{扌} + \zh{般} (bān) & 1. \zh{扌} (3 traits :横, 竖钩, 提) \quad 2. \zh{般} (haut \zh{舟} puis \zh{八} + \zh{其}) \\
\hline
\zh{离} & lí & \zh{禸} (liú/zhí, rare) & \zh{离} (forme moderne) & Haut \zh{㐭} (tigre/veau) puis bas \zh{禸} ; 19 traits. Mémoriser : « s'éloigner du haut vers le bas » \\
\hline
\zh{挤} & jǐ & \zh{扌} (shǒu, main) & \zh{扌} + \zh{齐} (qí) & 1. \zh{扌} \quad 2. \zh{齐} (haut \zh{亠}, milieu \zh{丷}, bas \zh{田}) — « main qui aligne/serre » \\
\hline
\zh{租} & zū & \zh{禾} (hé, grain/riz) & \zh{禾} + \zh{且} (qiě) & 1. \zh{禾} (5 traits) \quad 2. \zh{且} (6 traits) — « grain + phonétique » = loyer payé en nature jadis \\
\hline
\end{tabularx}
\end{center}
\end{propriete}

\begin{attention}[Caractères proches — Ne pas confondre]
\begin{itemize}
\item \zh{租} (zū, louer, radical \zh{禾}) $\neq$ \zh{祖} (zǔ, ancêtre, radical \zh{礻} shì, autel).
\item \zh{挤} (jǐ, bondé, radical \zh{扌}) $\neq$ \zh{齐} (qí, aligné, complet, radical \zh{齐} lui-même).
\item \zh{离} (lí, quitter) $\neq$ \zh{丽} (lì, beau, radical \zh{一}).
\item \zh{搬} (bān, déplacer) $\neq$ \zh{般} (bān, sorte, radical \zh{舟}).
\end{itemize}
\end{attention}

\begin{exercice}[Entraînement à l'écriture]
\begin{enumerate}
\item Écris 5 fois chaque caractère en respectant l'ordre des traits. Dis le nom des traits à voix haute (ex. : \zh{搬} : \emph{héng, shùgōu, tí…}).
\item Complète avec le bon caractère : \zh{我\_\_\_家了。} (déménagé) \quad \zh{公交车很\_\_\_。} (bondé) \quad \zh{我\_\_\_开农村。} (quitté) \quad \zh{我们\_\_\_了一个房间。} (loué)
\item \textbf{Dictée rapide} : Le professeur dicte : \zh{搬家, 离开, 挤, 租房子, 城市, 农村, 找工作, 机会, 困难, 希望}.
\end{enumerate}
\end{exercice}

\begin{corrige}[Exercice : Entraînement à l'écriture]
2. \zh{搬} \quad \zh{挤} \quad \zh{离} \quad \zh{租} \\
3. Vérifie les tons : bān jiā (1-1), lí kāi (2-1), jǐ (3), zū fáng zi (1-2-3), chéng shì (2-4), nóng cūn (2-1), zhǎo gōng zuò (3-1-4), jī huì (1-4), kùn nan (4-2), xī wàng (1-4).
\end{corrige}

\courssub{Structures et connecteurs du texte}

Pour articuler ton récit (cause, opposition, chronologie, projection), utilise ces \cle{connecteurs} indispensables à l'écrit comme à l'oral (HSK 3--4).

\begin{propriete}[Connecteurs logiques pour la lettre]
\begin{center}
\begin{tabularx}{\linewidth}{|X|X|X|}
\hline
\rowcolor{popblueL} Structure & Exemple (Caractères + Pinyin) & Traduction \\
\hline
\cle{因为…，所以…} (yīnwèi…, suǒyǐ…) & 因为农村没有工作，所以我搬到了城市。 (Yīnwèi nóngcūn méiyǒu gōngzuò, suǒyǐ wǒ bāndào le chéngshì.) & Parce qu'il n'y a pas de travail à la campagne, j'ai déménagé en ville. \\
\hline
\cle{虽然…，但是/可是…} (suīrán…, dànshì/kěshì…) & 虽然城市很挤，但是机会很多。 (Suīrán chéngshì hěn jǐ, dànshì jīhuì hěn duō.) & Bien que la ville soit bondée, il y a beaucoup d'opportunités. \\
\hline
\cle{先…，然后…} (xiān…, ránhòu…) & 我先找工作，然后继续学习。 (Wǒ xiān zhǎo gōngzuò, ránhòu jìxù xuéxí.) & Je cherche d'abord du travail, ensuite je continue mes études. \\
\hline
\cle{以后} (yǐhòu) en début ou milieu de phrase & 以后我想回农村帮助家人。 (Yǐhòu wǒ xiǎng huí nóngcūn bāngzhù jiārén.) & Plus tard, je veux retourner à la campagne aider ma famille. \\
\hline
\end{tabularx}
\end{center}
\end{propriete}

\begin{attention}[Pièges grammaticaux fréquents (Francophones)]
\begin{itemize}
\item \textbf{Ordre Temps/Lieu} : \cle{Sujet + Temps + Lieu + Verbe}. \zh{我上个月在雅温得租了房子。} (Wǒ shàng ge yuè zài Yǎwēndé zū le fángzi.) $\neq$ *\zh{我租了房子上个月在雅温得。}
\item \textbf{\zh{离开} (lí kāi) transitif direct} : \zh{离开家}, \zh{离开农村}. Pas de préposition \zh{从} (cóng) après \zh{离开}. (Mais : \zh{从农村离开} = « partir \textbf{de} la campagne »).
\item \textbf{\zh{了} (le) d'accompli} : Place \zh{了} \textbf{après le verbe} (ou V+Objet). \zh{我租了房间。} \zh{我搬到了雅温得。} Ne mets pas \zh{了} en fin de phrase si l'objet est long non quantifié (ex. \zh{我租了房间了} possible mais \zh{我租了一个房间} standard).
\item \textbf{\zh{挤} (jǐ) est un adjectif verbe} : \zh{车很挤} (Le bus est bondé). Pas de verbe \zh{是} (shì) avant. $\neq$ *\zh{车是挤。}
\item \textbf{Tons} : \zh{搬} (bān, 1\textsuperscript{er} ton haut plat), \zh{挤} (jǐ, 3\textsuperscript{e} ton bas descendant-remontant). Ne dis pas *bǎn ni *jí.
\end{itemize}
\end{attention}

\courssub{Thème et version guidés}

Avant la production libre, entraîne-toi à basculer du français au chinois (thème) et du chinois au français (version) sur des phrases types de la lettre.

\begin{exemplebox}[Thème guidé — Français $\rightarrow$ Chinois]
\textbf{Consigne} : Traduis en caractères + pinyin. Respecte l'ordre Sujet--Temps--Lieu--Verbe et les connecteurs.
\begin{enumerate}
\item Parce que le village est petit, je suis parti à la ville. \\
$\rightarrow$ \zh{因为村子小，所以我去了城市。} (Yīnwèi cūnzi xiǎo, suǒyǐ wǒ qù le chéngshì.)
\item Bien que le loyer soit cher, j'ai loué une chambre. \\
$\rightarrow$ \zh{虽然房租贵，但是我租了一个房间。} (Suīrán fángzū guì, dànshì wǒ zū le yí ge fángjiān.)
\item D'abord je mange, ensuite je cherche du travail. \\
$\rightarrow$ \zh{我先吃饭，然后找工作。} (Wǒ xiān chī fàn, ránhòu zhǎo gōngzuò.)
\item Plus tard, j'espère retourner au village. \\
$\rightarrow$ \zh{以后我希望回农村。} (Yǐhòu wǒ xīwàng huí nóngcūn.)
\item J'ai quitté la maison depuis deux mois. \\
$\rightarrow$ \zh{我离开家已经两个月了。} (Wǒ líkāi jiā yǐjīng liǎng ge yuè le.)
\end{enumerate}
\end{exemplebox}

\begin{exemplebox}[Version guidée — Chinois $\rightarrow$ Français]
\textbf{Consigne} : Traduis en français naturel. Identifie le connecteur et le temps (accompli \zh{了}, durée \zh{已经…了}).
\begin{enumerate}
\item \zh{因为想学习中文，所以我来了中国。} \\
$\rightarrow$ Parce que je veux apprendre le chinois, je suis venu en Chine.
\item \zh{虽然工作很累，可是工资不低。} \\
$\rightarrow$ Bien que le travail soit fatigant, le salaire n'est pas bas.
\item \zh{我先租房子，然后找工作。} \\
$\rightarrow$ Je loue d'abord une chambre, ensuite je cherche du travail.
\item \zh{以后我们一起回喀麦隆。} \\
$\rightarrow$ Plus tard, nous rentrerons ensemble au Cameroun.
\item \zh{城市人多车挤，生活有点儿困难。} \\
$\rightarrow$ En ville il y a beaucoup de monde et les voitures sont bondées, la vie est un peu difficile.
\end{enumerate}
\end{exemplebox}

\courssub{Production écrite et évaluation}

\begin{methode}[Rédaction de ta lettre — 7 étapes]
\begin{enumerate}
\item \textbf{Planifie} en français les 4 paragraphes logiques : (a) Raisons du départ (cause \zh{因为…所以…}) ; (b) Situation actuelle : logement, vie quotidienne (description) ; (c) Difficultés vs espoirs (opposition \zh{虽然…但是…}) ; (d) Projets chronologiques (\zh{先…然后…}, \zh{以后}).
\item \textbf{Sélectionne} 6+ mots du glossaire (coche-les) et 3+ connecteurs.
\item \textbf{Rédige} une phrase brouillon par idée en chinois (caractères + pinyin).
\item \textbf{Assemble} en lettre complète : Appel \zh{亲爱的妈妈：} -- Corps (8--10 phrases) -- Formule finale \zh{祝你们身体健康！} -- Signature \zh{你们的儿子 André}.
\item \textbf{Relis} : Ordre des mots ? \zh{了} bien placé ? Caractères \zh{搬离挤租} corrects ? Tons notés ? Ponctuation chinoise ?
\item \textbf{Échange} avec un camarade pour correction croisée (grille ci-dessous).
\item \textbf{Recopie} au propre (version finale pour portfolio/affichage).
\end{enumerate}
\end{methode}

\begin{exoresolu}[Lettre modèle commentée — « Lettre d'André à sa famille »]
\textbf{Consigne} : Rédige la lettre (8--10 phrases, 3 connecteurs min., 6 mots glossaire, 4 caractères complexes corrects).
\tcblower
亲爱的妈妈：\\
Qīn'ài de māma :\\
« Chère maman, »\\[4pt]
您好！我离开农村已经一个月了，很想您和家人。\\
Nín hǎo ! Wǒ líkāi nóngcūn yǐjīng yí ge yuè le, hěn xiǎng nín hé jiārén.\\
« Bonjour ! Cela fait un mois que j'ai quitté le village, vous me manquez beaucoup, vous et la famille. »\\[4pt]
因为农村没有工作，所以我搬到了雅温得。\\
Yīnwèi nóngcūn méiyǒu gōngzuò, suǒyǐ wǒ bāndào le Yǎwēndé.\\
« Parce qu'il n'y avait pas de travail au village, j'ai déménagé à Yaoundé. »\\[4pt]
现在我租了一个小房间，每天很早起床。\\
Xiànzài wǒ zū le yí ge xiǎo fángjiān, měi tiān hěn zǎo qǐchuáng.\\
« Maintenant j'ai loué une petite chambre, je me lève très tôt chaque jour. »\\[4pt]
虽然城市很大，公共汽车很挤，但是这里机会很多。\\
Suīrán chéngshì hěn dà, gōnggòng qìchē hěn jǐ, dànshì zhèlǐ jīhuì hěn duō.\\
« Bien que la ville soit grande et les bus bondés, il y a beaucoup d'opportunités ici. »\\[4pt]
我先找工作，然后继续学习中文。\\
Wǒ xiān zhǎo gōngzuò, ránhòu jìxù xuéxí Zhōngwén.\\
« Je cherche d'abord du travail, ensuite je continue à apprendre le chinois. »\\[4pt]
生活有一点儿困难，可是我不怕。\\
Shēnghuó yǒu yìdiǎnr kùnnan, kěshì wǒ bú pà.\\
« La vie est un peu difficile, mais je n'ai pas peur. »\\[4pt]
以后我希望回农村帮助你们。\\
Yǐhòu wǒ xīwàng huí nóngcūn bāngzhù nǐmen.\\
« Plus tard, j'espère retourner au village pour vous aider. »\\[4pt]
祝你们身体健康！\\
Zhù nǐmen shēntǐ jiànkāng !\\
« Je vous souhaite une bonne santé ! »\\[4pt]
你们的儿子 André\\
Nǐmen de érzi André\\[6pt]
\textbf{Commentaire détaillé des choix de langue :}
\begin{itemize}
\item \textbf{Connecteurs (4/3 requis)} : \zh{因为…所以…} (phrase 3), \zh{虽然…但是…} (phrase 5), \zh{先…然后…} (phrase 6), \zh{以后} (phrase 8).
\item \textbf{Lexique glossaire (9/6 requis)} : \zh{离开}, \zh{农村}, \zh{搬}, \zh{租}, \zh{挤}, \zh{找工作}, \zh{机会}, \zh{困难}, \zh{希望}. (+ \zh{生活}, \zh{继续}, \zh{帮助}).
\item \textbf{Grammaire} : Temps avant verbe (\zh{上个月}, \zh{每天}, \zh{现在}) ; \zh{了} accompli (\zh{搬到了}, \zh{租了}) ; \zh{已经…了} durée ; \zh{一点儿} atténuation.
\item \textbf{Caractères complexes} : \zh{搬} (扌+般), \zh{离} (结构上下), \zh{挤} (扌+齐), \zh{租} (禾+且) — tracés conformes.
\item \textbf{Registre} : \zh{您} (politesse), \zh{亲爱的} (affection), \zh{身体健康} (formule épistolaire standard).
\end{itemize}
\end{exoresolu}

\begin{exercice}[Production écrite individuelle — Évaluation formative]
\textbf{Sujet} : Écris \textbf{ta propre lettre} (tu es André ou un autre personnage : Marie, 18 ans, partie de Buea pour Douala ; ou Paul, 17 ans, parti de Maroua pour Yaoundé). Adapte les détails (ville, parent d'accueil, type de travail visé).
\begin{itemize}
\item Format : Lettre manuscrite sur feuille A5 (ou traitement de texte police \emph{Kaiti} / \emph{SimSun} 14 pt).
\item Contraintes obligatoires : 8--10 phrases ; 3 connecteurs différents ; 6 mots glossaire ; 4 caractères ciblés sans faute ; formules d'appel/de fin.
\item Durée : 30 min en classe.
\end{itemize}
\end{exercice}

\begin{corrige}[Grille d'évaluation par les pairs (20 points) — À donner à chaque correcteur]
\begin{center}
\begin{tabularx}{\linewidth}{|X|c|X|}
\hline
\rowcolor{popblueL} \textbf{Critère} & \textbf{Points} & \textbf{Indicateurs de réussite} \\
\hline
Connecteurs (3 min.) & 3 & 1 pt par connecteur \textbf{bien construit} et \textbf{bien placé} (因为所以, 虽然但是, 先然后, 以后). \\
\hline
Lexique glossaire (6 min.) & 6 & 1 pt par mot \textbf{correctement employé} (caractère + sens + place dans phrase). \\
\hline
Caractères complexes & 2 & 0,5 pt par caractère (\zh{搬}, \zh{离}, \zh{挤}, \zh{租}) \textbf{sans erreur de trait/radical}. \\
\hline
Grammaire / Syntaxe & 4 & Ordre S-T-L-V respecté ; \zh{了} accompli correct ; \zh{离开} + lieu direct ; négation \zh{没(有)} ; \zh{一点儿} ; pas de \zh{是} + adj. \\
\hline
Cohérence / Forme lettre & 5 & Appel + formule politesse + 4 paragraphes logiques + formule finale + signature ; écriture lisible, ponctuation chinoise. \\
\hline
\textbf{TOTAL} & \textbf{20} & \\
\hline
\end{tabularx}
\end{center}
\textbf{Note} : Si note $<$ 12/20 $\rightarrow$ réécriture obligatoire après remédiation ciblée (caractères ou connecteur faible).
\end{corrige}

\courssub{Bilan de la leçon et prolongements}

\begin{aretenir}[Ce que tu as maîtrisé]
\begin{itemize}
\item \textbf{Compétence écrite} : Produire un texte personnel structuré (lettre) sur un thème de société (exode rural).
\item \textbf{Lexique} : Maîtriser 10 mots-clés du domaine (migration, logement, travail, sentiments).
\item \textbf{Caractères} : Écrire \zh{搬}, \zh{离}, \zh{挤}, \zh{租} en analysant radical + phonétique + ordre des traits.
\item \textbf{Grammaire} : Articuler cause (\zh{因为…所以…}), opposition (\zh{虽然…但是…}), chronologie (\zh{先…然后…}, \zh{以后}) ; placer compléments temps/lieu ; utiliser \zh{了} (accompli) et \zh{已经…了} (durée).
\item \textbf{Socioculturel} : Comparer réalités migratoires Cameroun/Chine (vocabulaire \zh{农民工}, \zh{户口}, \zh{城镇化}).
\item \textbf{Méthodologie} : Planifier $\rightarrow$ Rédiger $\rightarrow$ Relire $\rightarrow$ Évaluer pair $\rightarrow$ Recopier (processus écriture).
\end{itemize}
\end{aretenir}

\begin{experience}[Prolongement oral — Débat / Exposé (15 min)]
\textbf{Thème} : \zh{留在城市还是回农村？} (Rester en ville ou retourner à la campagne ?)
\begin{enumerate}
\item \textbf{Préparation (5 min)} : En binôme, prépare 3 arguments \textbf{pour} la ville (emploi, études, santé, modernité) et 3 arguments \textbf{pour} le village (famille, coût vie, terre, traditions). Note mots-clés chinois.
\item \textbf{Débat (5 min)} : Binôme A = « Rester en ville », Binôme B = « Retourner au village ». Utilise : \zh{我认为…因为…} (Je pense que… car…), \zh{但是…} (Mais…), \zh{而且…} (De plus…).
\item \textbf{Synthèse (5 min)} : Chaque binôme présente une conclusion commune : \zh{我们建议…因为…} (Nous suggérons… car…). Le professeur note au tableau les arguments en caractères.
\end{enumerate}
\end{experience}

\begin{methode}[Pour l'évaluation sommative du Module 3 (Type Bac)]
\begin{enumerate}
\item \textbf{Compréhension écrite} : Lettre/email court (60--80 caractères) sur migration, études, travail $\rightarrow$ QCM / Vrai-Faux / Questions ouvertes.
\item \textbf{Expression écrite} : Sujet type « Tu es [Prénom], tu as quitté [Village] pour [Ville]. Écris à ta famille (80--100 caractères). Contraintes : 3 connecteurs, 8 mots unité 3, 5 caractères complexes vus en classe. »
\item \textbf{Traduction} : 5 phrases Thème (Fr $\rightarrow$ Ch) + 5 phrases Version (Ch $\rightarrow$ Fr) sur structures leçon (比, 把, 被, 因为所以, 虽然但是, 了/过/着).
\item \textbf{Caractères} : Dictée 15 mots (dont 5 complexes) + Écriture guidée 5 caractères (ordre traits + radical).
\end{enumerate}
Garde ta lettre corrigée dans ton \textbf{portfolio} : elle constitue une trace de travail pour l'oral de fin d'année et l'examen écrit.
\end{methode}

\newpage

\coursec{Méthodes et exercices résolus}

\coursec{Leçon 26 — Expression écrite : caractères complexes liés à l'exode rural}

\courssub{Modèle de texte commenté}

\begin{textebox}[Modèle : Les jeunes quittent le village pour la ville]
\zh{现在很多年轻人离开农村去大城市。} \\
\emph{Xiànzài hěn duō niánqīngrén líkāi nóngcūn qù dà chéngshì.} \\
« Aujourd'hui, beaucoup de jeunes quittent la campagne pour les grandes villes. » \\[4pt]
\zh{因为在农村工作很少，工资也不高。} \\
\emph{Yīnwèi zài nóngcūn gōngzuò hěn shǎo, gōngzī yě bù gāo.} \\
« Parce qu'à la campagne le travail est rare et les salaires ne sont pas élevés. » \\[4pt]
\zh{他们离开巴富萨姆，来到雅温得找工作。} \\
\emph{Tāmen líkāi Bāfùsàmǔ, láidào Yǎwēndé zhǎo gōngzuò.} \\
« Ils quittent Bafoussam et arrivent à Yaoundé pour chercher du travail. » \\[4pt]
\zh{所以农村的老人越来越多，年轻人越来越少。} \\
\emph{Suǒyǐ nóngcūn de lǎorén yuè lái yuè duō, niánqīngrén yuè lái yuè shǎo.} \\
« Donc à la campagne il y a de plus en plus de personnes âgées et de moins en moins de jeunes. » \\[4pt]
\zh{我希望政府帮助农村发展。} \\
\emph{Wǒ xīwàng zhèngfǔ bāngzhù nóngcūn fāzhǎn.} \\
« J'espère que le gouvernement aidera la campagne à se développer. »
\end{textebox}

\begin{aretenir}[Structure-type du texte explicatif (5 phrases)]
\begin{enumerate}
\item \textbf{Sujet + situation actuelle} : \zh{现在 + Sujet + Verbe + Lieu}. (Présenter le phénomène).
\item \textbf{Cause} : \zh{因为 + Raison (Lieu + Verbe/Adjectif)}. (Expliquer le départ).
\item \textbf{Mouvement précis} : \zh{Sujet + 离开 + Lieu\_origine + 来到/去 + Lieu\_destination + But}. (Décrire la migration).
\item \textbf{Conséquence} : \zh{所以 + Sujet + 越来越 + Adjectif}. (Montrer l'impact structurel).
\item \textbf{Opinion / Souhait} : \zh{Sujet + 希望/建议 + Sujet\_action + Verbe}. (Ouvrir une perspective).
\end{enumerate}
\cle{Règle d'or} : En chinois, le complément de lieu et de temps se place \textbf{avant} le verbe. Pas de phrase sans verbe.
\end{aretenir}

\courssub{Lexique écrit du thème : l'exode rural}

\begin{center}
\begin{tabularx}{\linewidth}{|X|X|X|X|}
\hline
\rowcolor{popblueL} \textbf{Caractères} & \textbf{Pinyin} & \textbf{Nature} & \textbf{Traduction} \\
\hline
农村 & nóngcūn & nom & campagne, zone rurale \\
\hline
城市 & chéngshì & nom & ville, zone urbaine \\
\hline
离开 & líkāi & verbe & quitter, partir de \\
\hline
来到 & láidào & verbe & arriver à, venir à \\
\hline
找工作 & zhǎo gōngzuò & v. + obj. & chercher du travail \\
\hline
发展 & fāzhǎn & verbe & se développer, développer \\
\hline
工作 & gōngzuò & nom / verbe & travail, emploi / travailler \\
\hline
工资 & gōngzī & nom & salaire \\
\hline
房租 & fángzū & nom & loyer \\
\hline
老人 & lǎorén & nom & personne âgée, vieux \\
\hline
年轻人 & niánqīngrén & nom & jeune, jeune personne \\
\hline
越来越 & yuè lái yuè & adv. & de plus en plus \\
\hline
因为 & yīnwèi & conj. & parce que \\
\hline
所以 & suǒyǐ & adv. & donc, par conséquent \\
\hline
但是 & dànshì & conj. & mais, cependant \\
\hline
而且 & érqiě & conj. & de plus, en plus \\
\hline
政府 & zhèngfǔ & nom & gouvernement \\
\hline
家乡 & jiāxiāng & nom & village natal, pays natal \\
\hline
想家 & xiǎng jiā & v. composé & avoir le mal du pays, manquer la maison \\
\hline
种地 & zhòng dì & v. + obj. & cultiver la terre \\
\hline
工厂 & gōngchǎng & nom & usine \\
\hline
杜阿拉 & Dù'ālā & n. propre & Douala \\
\hline
雅温得 & Yǎwēndé & n. propre & Yaoundé \\
\hline
巴富萨姆 & Bāfùsàmǔ & n. propre & Bafoussam \\
\hline
\end{tabularx}
\end{center}

\courssub{Savoir-faire 1 : Produire un texte simple sur l'exode rural}

\begin{methode}[Rédiger un texte simple sur l'exode rural en 5 étapes]
\begin{enumerate}
\item \textbf{Poser le sujet dès la première phrase} : utilisez \zh{现在很多年轻人离开农村} (xiànzài hěn duō niánqīngrén líkāi nóngcūn) ou \zh{目前} (mùqián, « actuellement ») pour situer le propos.
\item \textbf{Donner la cause} avec \zh{因为} (yīnwèi) en tête de proposition : \zh{因为农村工作少，收入低} (yīnwèi nóngcūn gōngzuò shǎo, shōurù dī).
\item \textbf{Décrire le mouvement} avec \zh{离开} (líkāi) + lieu d'origine et \zh{去} (qù) / \zh{来到} (láidào) + ville : \zh{他们离开家乡，去杜阿拉、雅温得找工作}.
\item \textbf{Montrer la conséquence} avec \zh{所以} (suǒyǐ) + \zh{越来越} (yuè lái yuè) : \zh{所以农村剩下老人越来越多}.
\item \textbf{Conclure} par une opinion (\zh{我认为} wǒ rènwéi) ou un souhait (\zh{我希望} wǒ xīwàng, \zh{建议} jiànyì).
\end{enumerate}
\cle{Réflexe} : Un texte de 5 à 8 phrases suffit. Chaque phrase suit l'ordre \cle{Sujet + Temps + Lieu + Verbe + Complément}. Jamais de phrase sans verbe. Ponctuation chinoise pleine chasse (。 ，).
\end{methode}

\begin{exoresolu}[Rédaction guidée : le départ des jeunes de Bafoussam]
Énoncé : Rédigez un texte de 5 à 6 phrases en chinois sur : « Les jeunes quittent le village de Bafoussam pour Yaoundé : causes et conséquences. » Utilisez au moins \zh{因为}, \zh{所以} et \zh{越来越}.
\tcblower
\textbf{Étape 1 — Plan.} Sujet (départ) → Cause (peu d'emplois, bas salaires) → Mouvement (Bafoussam → Yaoundé) → Conséquence (vieillissement village) → Conclusion (souhait développement).\\
\textbf{Étape 2 — Rédaction phrase par phrase.}
\begin{itemize}
\item Phrase 1 (Sujet) : \zh{现在很多年轻人离开农村去大城市。} \emph{Xiànzài hěn duō niánqīngrén líkāi nóngcūn qù dà chéngshì.} « Aujourd'hui beaucoup de jeunes quittent la campagne pour les grandes villes. »
\item Phrase 2 (Cause) : \zh{因为在农村工作很少，工资也不高。} \emph{Yīnwèi zài nóngcūn gōngzuò hěn shǎo, gōngzī yě bù gāo.} « Parce qu'à la campagne le travail est rare et les salaires ne sont pas élevés. »
\item Phrase 3 (Mouvement) : \zh{他们离开巴富萨姆，来到雅温得找工作。} \emph{Tāmen líkāi Bāfùsàmǔ, láidào Yǎwēndé zhǎo gōngzuò.} « Ils quittent Bafoussam et viennent à Yaoundé chercher du travail. »
\item Phrase 4 (Conséquence) : \zh{所以农村的老人越来越多，年轻人越来越少。} \emph{Suǒyǐ nóngcūn de lǎorén yuè lái yuè duō, niánqīngrén yuè lái yuè shǎo.} « Donc à la campagne il y a de plus en plus de vieux et de moins en moins de jeunes. »
\item Phrase 5 (Conclusion) : \zh{我希望政府帮助农村发展。} \emph{Wǒ xīwàng zhèngfǔ bāngzhù nóngcūn fāzhǎn.} « J'espère que le gouvernement aidera la campagne à se développer. »
\end{itemize}
\textbf{Étape 3 — Vérification.} Ordre S + T + L + V respecté ; \zh{因为} en tête de cause, \zh{所以} en tête de conséquence ; \zh{越来越} + adjectif sans \zh{很}. Connecteurs exigés présents. Texte final correct.
\end{exoresolu}

\courssub{Savoir-faire 2 : Écrire correctement les caractères complexes du thème}

\begin{methode}[Mémoriser et tracer les caractères complexes : 5 étapes]
\begin{enumerate}
\item \textbf{Décomposer} en clé (radical) + partie phonétique. Ex. : \zh{离} (lí) = \zh{禸} + \zh{隹} (simplifié : clé \zh{禸}) ; \zh{城} (chéng) = clé \zh{土} (terre) + \zh{成} (phonétique) ; \zh{村} (cūn) = clé \zh{木} (bois) + \zh{寸}.
\item \textbf{Respecter l'ordre des traits} : ① Haut → Bas ② Gauche → Droite ③ Horizontal avant Vertical ④ Extérieur avant Intérieur ⑤ Fermer la boîte en dernier.
\item \textbf{Comparer les caractères proches} (paires à ne pas confondre) :
\begin{itemize}
\item \zh{城} (chéng, ville, clé \zh{土}) vs \zh{诚} (chéng, sincère, clé \zh{讠}).
\item \zh{村} (cūn, village, clé \zh{木}) vs \zh{树} (shù, arbre, clé \zh{木}) vs \zh{林} (lín, forêt, clé \zh{木}).
\item \zh{找} (zhǎo, chercher, clé \zh{扌} main) vs \zh{我} (wǒ, je, clé \zh{戈}).
\item \zh{离} (lí, quitter) vs \zh{禽} (qín, volaille — clé \zh{隹} oiseau).
\end{itemize}
\item \textbf{Écrire chaque caractère 5 fois} en prononçant le pinyin à voix haute (mémoire gestuelle + auditive).
\item \textbf{Se relire} en vérifiant la présence de la clé attendue : ville/lieu → souvent \zh{土} ou \zh{阝} ; action manuelle → \zh{扌} ; parole → \zh{讠} ; nature/bois → \zh{木}.
\end{enumerate}
\end{methode}

\begin{popfigure}
\begin{tikzpicture}[node distance=1.2cm, every node/.style={font=\small, align=center}]
\node (cheng) [draw=popblue, thick, rounded corners, fill=popblueL, minimum width=2.5cm, align=center]{\zh{城}\\\emph{chéng}\\\textbf{ville}\\\textit{Clé : 土 (terre)}};
\node (cheng2) [right=of cheng, draw=poporange, thick, rounded corners, fill=poporangeL, minimum width=2.5cm, align=center]{\zh{诚}\\\emph{chéng}\\\textbf{sincère}\\\textit{Clé : 讠 (parole)}};
\node (cun) [below=of cheng, draw=popgreen, thick, rounded corners, fill=popgreenL, minimum width=2.5cm, align=center]{\zh{村}\\\emph{cūn}\\\textbf{village}\\\textit{Clé : 木 (bois)}};
\node (shu) [right=of cun, draw=popgreen, thick, rounded corners, fill=popgreenL, minimum width=2.5cm, align=center]{\zh{树}\\\emph{shù}\\\textbf{arbre}\\\textit{Clé : 木 (bois)}};
\node (zhao) [below=of cun, draw=poppink, thick, rounded corners, fill=poppinkL, minimum width=2.5cm, align=center]{\zh{找}\\\emph{zhǎo}\\\textbf{chercher}\\\textit{Clé : 扌 (main)}};
\node (wo) [right=of zhao, draw=poppurple, thick, rounded corners, fill=poppurpleL, minimum width=2.5cm, align=center]{\zh{我}\\\emph{wǒ}\\\textbf{je}\\\textit{Clé : 戈 (hallebard)}};
\draw[popline, thick, ->] (cheng.east) -- (cheng2.west) node[midway, above, font=\tiny] {Même son, clés diff.};
\draw[popline, thick, ->] (cun.east) -- (shu.west) node[midway, above, font=\tiny] {Même clé, sens diff.};
\draw[popline, thick, ->] (zhao.east) -- (wo.west) node[midway, above, font=\tiny] {Formes très proches};
\end{tikzpicture}
\legende{Caractères proches à distinguer : clés et sens}
\end{popfigure}

\begin{exoresolu}[Dictée et analyse de caractères]
Énoncé : a) Écrivez en caractères : nóngcūn, chéngshì, líkāi, fāzhǎn, gōngzuò. b) Pour \zh{城} et \zh{村}, indiquez la clé et citez 2 autres mots du thème contenant cette clé. c) Distinguez \zh{找} et \zh{我} par la clé et le sens.
\tcblower
\textbf{a) Réponses :} \zh{农村} (nóngcūn), \zh{城市} (chéngshì), \zh{离开} (líkāi), \zh{发展} (fāzhǎn), \zh{工作} (gōngzuò).\\
\textbf{b) Analyse des clés :}
\begin{itemize}
\item \zh{城} (chéng) : clé \zh{土} (tǔ, terre, 3 traits, à gauche). Mots du thème : \zh{地} (dì, terre/sol), \zh{场} (chǎng, place/terrain), \zh{在} (zài, être à — contient \zh{土} en bas). Logique : la ville est construite sur la terre, murs de terre.
\item \zh{村} (cūn) : clé \zh{木} (mù, bois/arbre, 4 traits, à gauche). Mots du thème : \zh{林} (lín, forêt), \zh{树} (shù, arbre). Logique : le village traditionnel est entouré de bois, maisons en bois.
\end{itemize}
\textbf{c) \zh{找} vs \zh{我} :}
\begin{itemize}
\item \zh{找} (zhǎo) : clé \zh{扌} (shǒu, main, 3 traits). Sens : action de la main pour chercher. Structure : gauche-droite.
\item \zh{我} (wǒ) : clé \zh{戈} (gē, hallebarde/arme, 4 traits). Sens originel : arme, puis « moi » (emprunt phonétique). Structure : haut-bas (anciennement).
\end{itemize}
\cle{Piège} : Écrire \zh{我工作} pour « je cherche du travail » → Faux. Correct : \zh{我找工作} (wǒ zhǎo gōngzuò).
\end{exoresolu}

\courssub{Savoir-faire 3 : Structurer un texte avec les connecteurs logiques}

\begin{propriete}[Tableau récapitulatif des connecteurs du thème]
\begin{center}
\begin{tabularx}{\linewidth}{|X|X|X|}
\hline
\rowcolor{popblueL} \textbf{Fonction} & \textbf{Connecteur(s)} & \textbf{Position / Structure} \\
\hline
Cause & \zh{因为} (yīnwèi) & En tête de la prop. de cause. Paire possible \zh{因为……所以……} \\
\hline
Conséquence & \zh{所以} (suǒyǐ) & En tête de la prop. de conséquence. Jamais seul pour une cause. \\
\hline
Opposition & \zh{但是} (dànshì) / \zh{可是} (kěshì) & En tête de 2\textsuperscript{e} prop. : \zh{Phrase 1，但是 Phrase 2.} \\
\hline
Addition & \zh{而且} (érqiě) / \zh{也} (yě) & \zh{而且} en tête de prop. ; \zh{也} avant le verbe. \\
\hline
Évolution & \zh{越来越} (yuè lái yuè) + Adj. & \zh{Sujet + 越来越 + Adj.} \textbf{Jamais} \zh{很} devant. \\
\hline
Enchaînement & \zh{先……然后……} (xiān... ránhòu...) & \zh{Sujet + 先 + V1，然后 + V2.} \\
\hline
\end{tabularx}
\end{center}
\end{propriete}

\begin{methode}[Choisir et placer le bon connecteur]
\begin{enumerate}
\item \textbf{Identifier la relation logique} entre deux idées (cause/effet, opposition, ajout, temps).
\item \textbf{Placer le connecteur en tête de la deuxième proposition} (sauf \zh{也}, \zh{越来越}).
\item \textbf{Ne jamais relier deux phrases complètes par une simple virgule} en chinois écrit formel ; utiliser un connecteur ou un point。
\item \textbf{Vérifier la paire \zh{因为……所以……}} : en chinois elle est correcte (contrairement au français « parce que... donc »).
\end{enumerate}
\end{methode}

\begin{exoresolu}[Compléter un texte à trous avec les connecteurs]
Énoncé : Complétez avec \zh{因为}, \zh{所以}, \zh{但是}, \zh{而且}, \zh{越来越} :\\
\zh{农村的生活很安静，\_\_\_\_ 工作很少。\_\_\_\_ 很多年轻人去城市。城市很大，\_\_\_\_ 很热闹，\_\_\_\_ 房租很贵。在农村，老人\_\_\_\_ 多。}
\tcblower
\textbf{Raisonnement phrase par phrase :}
\begin{enumerate}
\item « Vie calme, \_\_\_\_ travail rare » : Opposition → \zh{但是}. \zh{农村的生活很安静，但是工作很少。}
\item « \_\_\_\_ beaucoup de jeunes vont en ville » : Conséquence de la phrase 1 → \zh{所以}. \zh{所以很多年轻人去城市。}
\item « Ville grande, \_\_\_\_ très animée » : Addition → \zh{而且}. \zh{城市很大，而且很热闹。}
\item « \_\_\_\_ loyer très cher » : Opposition (animé MAIS cher) → \zh{但是}. \zh{但是房租很贵。}
\item « Vieux \_\_\_\_ nombreux » : Évolution → \zh{越来越}. \zh{在农村，老人越来越多。}
\end{enumerate}
\textbf{Texte complet :} \zh{农村的生活很安静，但是工作很少。所以很多年轻人去城市。城市很大，而且很热闹，但是房租很贵。在农村，老人越来越多。}\\
\cle{Vérification} : \zh{因为} non nécessaire (pas de prop. de cause explicite à introduire) ; chaque connecteur en tête de prop.
\end{exoresolu}

\courssub{Savoir-faire 4 : Traduire du français vers le chinois (Thème)}

\begin{methode}[Réussir le thème en 4 étapes]
\begin{enumerate}
\item \textbf{Identifier le verbe principal} et son équivalent chinois exact.
\item \textbf{Reconstruire dans l'ordre chinois} : \cle{Sujet + Temps + Lieu (\zh{在}) + Verbe + Objet}. Le lieu/temps précèdent le verbe (contrairement au français).
\item \textbf{Choisir les mots-outils} : \zh{了} (action accomplie), \zh{的} (adjectif + nom), \textbf{pas de \zh{是} devant adjectif} (\zh{很贵} ≠ \zh{是很贵}), classificateur si Num + Nom.
\item \textbf{Relire} : Tons du pinyin, ponctuation chinoise 。，, ordre des mots, classificateurs.
\end{enumerate}
\end{methode}

\begin{exoresolu}[Thème : trois phrases sur l'exode rural]
Énoncé : Traduisez en chinois : a) « Beaucoup de jeunes quittent le village pour aller travailler à Douala. » b) « Parce que la vie en ville est très chère, certains veulent retourner à la campagne. » c) « Le gouvernement veut développer les campagnes. »
\tcblower
\textbf{a)} Verbes : « quitter » = \zh{离开} ; « aller travailler » = \zh{去……工作}. Lieu avant verbe.\\
\zh{很多年轻人离开农村去杜阿拉工作。} \emph{Hěn duō niánqīngrén líkāi nóngcūn qù Dù'ālā gōngzuò.}\\
\cle{Justification} : \zh{很多} + nom sans classificateur ; enchaînement verbes \zh{离开} + \zh{去} sans mot-outil.\\
\textbf{b)} Cause \zh{因为} en tête ; « vie en ville » = \zh{城市的生活} (\zh{的} relation) ; « chère » = \zh{很贵} (pas \zh{是}) ; « certains » = \zh{有的人} / \zh{一些人} ; « retourner » = \zh{回} / \zh{回去}.\\
\zh{因为城市的生活很贵，所以有的人想回农村。} \emph{Yīnwèi chéngshì de shēnghuó hěn guì, suǒyǐ yǒu de rén xiǎng huí nóngcūn.}\\
\cle{Justification} : Paire \zh{因为……所以……} correcte en chinois ; \zh{想} + V = « vouloir faire ».\\
\textbf{c)} « Vouloir développer » = \zh{想发展} / \zh{打算发展} ; « campagnes » = \zh{农村}.\\
\zh{政府想发展农村。} \emph{Zhèngfǔ xiǎng fāzhǎn nóngcūn.}\\
\cle{Vérification finale} : Ordre S + V + O respecté ; lieu/temps jamais rejetés en fin de phrase.
\end{exoresolu}

\courssub{Savoir-faire 5 : Traduire du chinois vers le français (Version)}

\begin{methode}[Réussir la version en 4 étapes]
\begin{enumerate}
\item \textbf{Découper} la phrase chinoise en blocs : Sujet / Temps / Lieu (\zh{在}...) / Verbe / Objet.
\item \textbf{Traduire bloc par bloc} puis réordonner selon la syntaxe française (Lieu/Temps souvent après le verbe).
\item \textbf{Rendre les mots-outils} : \zh{了} = accompli (passé composé) ; \zh{过} = expérience (« avoir déjà ») ; \zh{越来越} = « de plus en plus » ; \zh{比} = comparatif « plus... que ».
\item \textbf{Produire un français naturel} : accords, articles, pas de mot-à-mot (ex. \zh{想家} = « avoir le mal du pays », pas « penser maison »).
\end{enumerate}
\end{methode}

\begin{exoresolu}[Version : texte sur l'exode rural]
Énoncé : Traduisez en français :\\
\zh{我的表哥以前在村里种地。去年他离开了家乡，来到大城市找工作。现在他在一个工厂工作，工资比在农村高很多。但是他说，城市的生活很累，他想家。}
\emph{Wǒ de biǎogē yǐqián zài cūn lǐ zhòng dì. Qùnián tā líkāi le jiāxiāng, láidào dà chéngshì zhǎo gōngzuò. Xiànzài tā zài yí ge gōngchǎng gōngzuò, gōngzī bǐ zài nóngcūn gāo hěn duō. Dànshì tā shuō, chéngshì de shēnghuó hěn lèi, tā xiǎng jiā.}
\tcblower
\textbf{Découpage et traduction :}
\begin{enumerate}
\item \zh{我的表哥} (mon cousin aîné) + \zh{以前} (autrefois) + \zh{在村里} (au village) + \zh{种地} (cultivait la terre) → « Mon cousin cultivait autrefois la terre au village. » (\zh{以前} = habitude passée → imparfait).
\item \zh{去年} (l'an dernier) + \zh{他离开了家乡} (\zh{了} = accompli) + \zh{来到大城市找工作} → « L'an dernier, il a quitté son village natal et est parti chercher du travail dans une grande ville. »
\item \zh{现在他在一个工厂工作} (classif. \zh{个}) + \zh{工资比在农村高很多} (\zh{比} comparatif) → « Maintenant, il travaille dans une usine et son salaire est bien plus élevé qu'à la campagne. »
\item \zh{但是他说} + \zh{城市的生活很累} + \zh{他想家} → « Mais il dit que la vie en ville est épuisante et que son pays lui manque. »
\end{enumerate}
\textbf{Texte final :} « Mon cousin cultivait autrefois la terre au village. L'an dernier, il a quitté son village natal pour chercher du travail dans une grande ville. Maintenant, il travaille dans une usine et son salaire est bien plus élevé qu'à la campagne. Mais il dit que la vie en ville est épuisante et que son pays lui manque. »\\
\cle{Vigilance} : \zh{了} → passé composé ; \zh{比} → « plus... que » ; \zh{想家} = verbe composé (mal du pays).
\end{exoresolu}

\courssub{Production écrite et évaluation}

\begin{exercice}[Entraînement 1 : Rédaction courte]
Rédigez un texte de 6 phrases en chinois sur le thème : « Un jeune de ton village part à Douala. » Contraintes : Utiliser \zh{因为}, \zh{所以}, \zh{但是}, \zh{越来越}, \zh{离开}, \zh{来到}, \zh{找工作}, \zh{工资}, \zh{房租}, \zh{想家}. Écrire caractères + pinyin + traduction.
\end{exercice}

\begin{exercice}[Entraînement 2 : Thème / Version]
\begin{enumerate}
\item \textbf{Thème :} Traduisez en chinois : « Autrefois, mon oncle cultivait le riz à la campagne. L'année dernière, il est parti à Yaoundé. Maintenant, il conduit un taxi. Il gagne plus d'argent, mais il trouve la vie fatigante. Il veut rentrer au village l'année prochaine. » (Voc. aide : 种水稻 zhòng shuǐdào, 开出租车 kāi chūzūchē, 赚钱 zhuànqián, 明年 míngnián, 回去 huíqù).
\item \textbf{Version :} Traduisez en français : \zh{我的妹妹不想在农村种地。她觉得农村太安静了。所以她去了杜阿拉。在那里她找到了一份工作。虽然工资不高，但她很开心。} (Voc. : 妹妹 mèimei, 觉得 juéde, 太...了 tài...le, 虽然...但... suīrán... dàn..., 开心 kāixīn).
\end{enumerate}
\end{exercice}

\begin{corrige}[Exercices 1 et 2]
\textbf{Ex. 1 (Proposition) :}
\zh{现在很多年轻人离开家乡。} (Xiànzài hěn duō niánqīngrén líkāi jiāxiāng.)
\zh{因为农村工作很少，工资低。} (Yīnwèi nóngcūn gōngzuò hěn shǎo, gōngzī dī.)
\zh{他们离开村子，来到杜阿拉找工作。} (Tāmen líkāi cūnzi, láidào Dù'ālā zhǎo gōngzuò.)
\zh{所以村里老人越来越多。} (Suǒyǐ cūn lǐ lǎorén yuè lái yuè duō.)
\zh{但是城市房租很贵，生活很累。} (Dànshì chéngshì fángzū hěn guì, shēnghuó hěn lèi.)
\zh{他们常常想家，希望将来回乡发展。} (Tāmen chángcháng xiǎng jiā, xīwàng jiānglái huí xiāng fāzhǎn.)

\textbf{Ex. 2.1 Thème :}
\zh{我叔叔以前在农村种水稻。} (Wǒ shūshu yǐqián zài nóngcūn zhòng shuǐdào.)
\zh{去年他离开家乡，来到雅温得。} (Qùnián tā líkāi jiāxiāng, láidào Yǎwēndé.)
\zh{现在他开出租车。} (Xiànzài tā kāi chūzūchē.)
\zh{他赚的钱比在农村多。} (Tā zhuàn de qián bǐ zài nóngcūn duō.)
\zh{但是他觉得城市生活很累。} (Dànshì tā juéde chéngshì shēnghuó hěn lèi.)
\zh{他明年想回去农村。} (Tā míngnián xiǎng huíqù nóngcūn.)

\textbf{Ex. 2.2 Version :}
« Ma jeune sœur ne veut pas cultiver la terre à la campagne. Elle trouve la campagne trop calme. Donc elle est allée à Douala. Là-bas, elle a trouvé un travail. Bien que le salaire ne soit pas élevé, elle est très contente. »
\end{corrige}

\courssub{Élément socioculturel : Exode rural Cameroun / Chine}

\begin{savaistu}[L'exode rural : regards croisés Cameroun — Chine]
\textbf{Cameroun :} L'exode rural massif vers Douala et Yaoundé (mais aussi Bafoussam, Ngaoundéré) concerne surtout les jeunes (15–35 ans) fuyant le sous-emploi agricole, l'enclavement et la pauvreté. Conséquences : vieillissement des villages, pression sur l'urbain (logement, chômage informel), perte de savoir-faire agricoles. Politiques : \emph{Programme d'Appui au Développement Rural}, décentralisation, promotion de l'agro-industrie pour fixer les jeunes.

\textbf{Chine :} Plus grande migration interne de l'histoire (années 1980–2020) : \zh{农民工} (nóngmíngōng, « ouvriers paysans ») — 290 millions (2023) quittent les provinces intérieures (Henan, Sichuan, Anhui) pour les mégalopoles côtières (Shenzhen, Dongguan, Shanghai, Pékin). Système \zh{户口} (hùkǒu, registre de résidence) : limite l'accès aux services urbains (école, santé, logement) pour les migrants. Depuis 2014 : réforme \zh{新型城镇化} (nouvelle urbanisation) + développement \zh{乡村振兴} (revitalisation rurale) pour rééquilibrer.

\textbf{Comparaison :} Au Cameroun, migration souvent définitive, informelle, peu encadrée. En Chine, migration historiquement circulaire (retour au village pour moissons/fêtes), encadrée par le \zh{户口}, aujourd'hui en transition vers urbanisation permanente. Point commun : le \cle{défi du développement rural} pour retenir la jeunesse.
\end{savaistu}

\courssub{Expression orale : Débat guidé}

\begin{experience}[Débat : « Faut-il rester au village ou partir en ville ? »]
\textbf{Consigne :} En groupes de 4. Chaque groupe tire un rôle : ① Jeune parti en ville (réussi) ② Jeune parti en ville (échec) ③ Parent resté au village ④ Maire / Chef de village.
\textbf{Préparation (10 min) :} Chaque élève prépare 3 arguments en chinois (mots-clés + phrases types du cours).
\textbf{Débat (15 min) :} Tour de table. Le modérateur (professeur) note au tableau les structures entendues : \zh{因为……所以……}, \zh{虽然……但是……}, \zh{越来越}, \zh{比}.
\textbf{Synthèse (5 min) :} Vote à main levée : « La ville est-elle la seule solution ? ». Écriture collective d'une phrase de conclusion au tableau.
\textbf{Évaluation :} Grille : Prononciation (tons), Vocabulaire thème, Structures (connecteurs), Interaction.
\end{experience}

\begin{attention}[Les 5 erreurs qui coûtent des points au Bac]
\begin{enumerate}
\item \textbf{Tons oubliés ou en chiffres.} Écrire « wo3 xiang3 jia1 » est sanctionné ; il faut \zh{wǒ xiǎng jiā}. Un ton faux change le sens : \zh{mǎ} (马, cheval) ≠ \zh{mā} (妈, maman) ≠ \zh{mà} (骂, gronder).
\item \textbf{Confondre caractères proches.} \zh{找} (zhǎo, chercher, clé \zh{扌}) / \zh{我} (wǒ, je, clé \zh{戈}) ; \zh{城} (chéng, ville, clé \zh{土}) / \zh{诚} (chéng, sincère, clé \zh{讠}) ; \zh{村} (cūn, village, clé \zh{木}) / \zh{树} (shù, arbre, clé \zh{木}). \cle{Vérifiez la clé !}
\item \textbf{Calquer l'ordre français (Lieu/Temps après verbe).} Faux : \zh{他工作在城市} / \zh{我去年离开了农村} (si « l'an dernier » mis à la fin). Correct : \zh{他在城市工作} / \zh{我去年离开了农村} (Temps en tête).
\item \textbf{Mots-outils mal employés.} Pas de \zh{是} devant adjectif : \zh{房租很贵} (pas \zh{房租是很贵}). Pas de \zh{了} après verbe d'état/habitude. Pas de \zh{很} avec \zh{越来越} : \zh{越来越多} (jamais \zh{很越来越多}).
\item \textbf{Classificateurs oubliés ou en trop.} \zh{一个工厂} (yí ge gōngchǎng), \zh{一份工作} (yí fèn gōngzuò), \zh{一辆出租车} (yí liàng chūzūchē). Mais \zh{很多年轻人} (SANS classificateur). \cle{Version} : Ne pas traduire \zh{因为……所以……} par « Parce que... donc » (faute de style français) ; choisir l'un des deux.
\end{enumerate}
\end{attention}

\newpage

\coursec{Exercices}

\courssub{Je vérifie mes connaissances}

\begin{exercice}[QCM : vocabulaire de l'exode rural]
Choisissez la bonne réponse parmi les trois propositions. Une seule réponse est correcte.
\begin{enumerate}
\item Comment dit-on « déménager / s'installer ailleurs » en chinois ?
      \begin{itemize}
      \item A. 离开
      \item B. 搬家
      \item C. 返回
      \end{itemize}
\item Quel mot signifie « loyer » ?
      \begin{itemize}
      \item A. 房租
      \item B. 工资
      \item C. 物价
      \end{itemize}
\item Le caractère \zh{挤} (jǐ) dans \zh{拥挤} signifie :
      \begin{itemize}
      \item A. Calme
      \item B. Bondé, serré
      \item C. Propre
      \end{itemize}
\item Pour exprimer « bien que / même si », on utilise la conjonction :
      \begin{itemize}
      \item A. 因为…所以…
      \item B. 虽然…但是…
      \item C. 不但…而且…
      \end{itemize}
\end{enumerate}
\end{exercice}

\begin{exercice}[Vrai ou Faux : caractères et structures]
Indiquez si les affirmations sont vraies (V) ou fausses (F). Justifiez brièvement en français pour les réponses fausses.
\begin{enumerate}
\item Le radical de \zh{搬} est \zh{扌} (main) et il indique un rapport avec l'action de la main.
\item Les caractères \zh{村} et \zh{材} se prononcent tous les deux \emph{cūn} et s'écrivent de la même façon.
\item Dans la phrase \zh{他离开家乡去城市打工}, le verbe \zh{离开} est suivi directement du lieu sans préposition.
\item La structure \zh{越来越} + adjectif s'écrit \zh{越来越贵} et non \zh{越贵越来}.
\end{enumerate}
\end{exercice}

\begin{exercice}[Vocabulaire : association caractères / pinyin / sens]
Reliez chaque caractère (colonne 1) à son pinyin (colonne 2) et à sa traduction (colonne 3) en traçant des flèches.
\begin{center}
\begin{tabularx}{\linewidth}{|X|X|X|}\hline
\rowcolor{popblueL} \textbf{Caractère} & \textbf{Pinyin} & \textbf{Traduction} \\ \hline
搬 & cūn & loyer \\ \hline
村 & jǐ & village \\ \hline
离 & bān & quitter, s'éloigner \\ \hline
挤 & lí & bondé, serré \\ \hline
租 & zū & déménager, déplacer \\ \hline
\end{tabularx}
\end{center}
\end{exercice}

\begin{exercice}[Pinyin et tons : discrimination auditive]
Écoutez les syllabes (ou lisez-les mentalement) et écrivez le pinyin avec les tons corrects pour les mots suivants. Attention aux tons changeants (3e ton + 3e ton).
\begin{enumerate}
\item \zh{搬家} \hspace{2cm}
\item \zh{农村} \hspace{2cm}
\item \zh{离开} \hspace{2cm}
\item \zh{拥挤} \hspace{2cm}
\item \zh{房租} \hspace{2cm}
\end{enumerate}
\end{exercice}

\courssub{Je m'entraîne}

\begin{exercice}[Thème : phrases à traduire en chinois]
Traduisez les phrases suivantes en caractères chinois (simplifiés). Ajoutez le pinyin avec tons.
\begin{enumerate}
\item Parce qu'il n'y a pas de travail au village, il part en ville.
\item Bien que la ville soit pratique, le loyer est trop cher.
\item De plus en plus de jeunes quittent la campagne.
\item Mes parents ne veulent pas déménager, ils aiment le village.
\item La route est très bondée, il y a trop de voitures.
\end{enumerate}
\end{exercice}

\begin{exercice}[Version : phrases à traduire en français]
Traduisez les phrases suivantes en français correct.
\begin{textebox}[Phrases pour version]
1. 村里没有工厂，所以年轻人都搬去城市了。\\
2. 虽然城市工作机会多，但是生活压力很大。\\
3. 我哥哥去年离开家乡，在深圳打工。\\
4. 城市里房租越来越贵，住房很拥挤。\\
5. 因为交通不方便，所以很少人回村里。
\end{textebox}
\end{exercice}

\begin{exercice}[Caractères complexes : écriture et radicaux]
Écrivez les caractères suivants en respectant l'ordre des traits (indiquez le nombre de traits). Entourez ou soulignez le radical (clé) de chaque caractère. Complétez le tableau.
\begin{center}
\begin{tabularx}{\linewidth}{|X|X|X|X|X|}\hline
\rowcolor{popblueL} Caractère & Nombre de traits & Radical (clé) & Sens du radical & Ordre des traits (description brève) \\ \hline
搬 & & & & \\ \hline
离 & & & & \\ \hline
挤 & & & & \\ \hline
租 & & & & \\ \hline
\end{tabularx}
\end{center}
\end{exercice}

\begin{exercice}[Discrimination de caractères proches : complétez les phrases]
Choisissez le bon caractère entre parenthèses pour compléter chaque phrase. Écrivez la phrase complète en caractères.
\begin{enumerate}
\item 我家在山里的一个小 \zh{(村 / 材)} 里。
\item 请帮我把书 \zh{(搬 / 般)} 到楼上。
\item 地铁里人太多了，非常 \zh{(挤 / 济)}。
\item 这 \zh{(村 / 材)} 料很好，很结实。
\item 他 \zh{(搬 / 般)} 家了，新地址我不知道。
\end{enumerate}
\end{exercice}

\begin{exercice}[Connecteurs et structure : remise en ordre]
Remettez les phrases suivantes dans l'ordre logique pour former un paragraphe cohérent sur l'exode rural. Utilisez les connecteurs comme indices. Écrivez le texte complet en caractères.
\begin{enumerate}
\item \zh{因为 村里 没有 工作}
\item \zh{所以 他们 搬去 城市}
\item \zh{虽然 城市 生活 方便}
\item \zh{但是 房租 很贵}
\item \zh{现在 越来越多 年轻人 离开 家乡}
\item \zh{这是 一个 严重 问题}
\end{enumerate}
\end{exercice}

\begin{exercice}[Texte à trous : mots-outils et vocabulaire clé]
Complétez le texte suivant avec les mots manquants choisis dans la liste ci-dessous. Écrivez les caractères dans les cases.
\zh{因为，所以，虽然，但是，越来越多，搬家，离开，拥挤，房租，工作}

\zh{现在，\_\_\_\_\_\_ 青年人 \_\_\_\_\_\_ 家乡 \_\_\_\_\_\_ 城市打工。\_\_\_\_\_\_ 村里 没有 好 \_\_\_\_\_\_，\_\_\_\_\_\_ 他们 必须 \_\_\_\_\_\_。\_\_\_\_\_\_ 城市 机会 多，\_\_\_\_\_\_ 城市 很 \_\_\_\_\_\_，\_\_\_\_\_\_ 也 很贵。}
\end{exercice}

\courssub{Je me prépare au Bac}

\begin{exercice}[Compréhension de l'écrit et langue : texte original]
Lisez le texte suivant et répondez aux questions en français.

\begin{textebox}[L'exode rural : défis et espoirs]
现在，喀麦隆 和 中国 的 农村 都 面临 同一个 问题：年轻人 离开 村庄 去 城市。\\
因为 农村 的 土地 少，收成 不 稳定，没有 工厂，所以 找 工作 很 难。\\
虽然 城市 生活 方便，有 医院、商场 和 好 学校，但是 房租 很贵，交通 也 很 拥挤。\\
很多 年轻人 在 城市 当 服务员、司机 或 在 工地 干活。他们 工资 不 高，压力 很大。\\
专家 说：如果 农村 有 好 路、好 网络 和 工厂，年轻人 就 不想 走 了。\\
我们 应该 发展 农村，让 家乡 变 好。
\end{textebox}

\textbf{Questions de compréhension :}
\begin{enumerate}
\item Selon le texte, quels sont les deux pays concernés par l'exode rural ?
\item Citez trois raisons pour lesquelles les jeunes quittent le village (d'après le premier paragraphe).
\item Citez deux avantages et deux inconvénients de la vie en ville mentionnés dans le texte.
\item Quel genre de travail font souvent les jeunes en ville ? Leur situation est-elle facile ?
\item Quelle solution les experts proposent-ils pour freiner l'exode ?
\end{enumerate}

\textbf{Questions de langue :}
\begin{enumerate}
\item Identifiez dans le texte la structure exprimant la concession (bien que... mais...). Recopiez la phrase complète en caractères.
\item Trouvez dans le texte le mot qui signifie « récolte / rendement agricole ».
\item Transformez la phrase \zh{年轻人 离开 村庄} en utilisant \zh{越来越多} au début de la phrase.
\item Écrivez les caractères correspondant au pinyin : \emph{zhǔjiā} (expert), \emph{fāzhǎn} (développer).
\end{enumerate}
\end{exercice}

\begin{exercice}[Thème et Version type Baccalauréat]
\textbf{Thème (Français $\rightarrow$ Chinois) :} Traduisez le texte suivant en chinois (caractères simplifiés + pinyin). Environ 60--80 caractères.
\begin{quote}
« Parce qu'il n'y a pas assez de travail au village, de plus en plus de jeunes déménagent en ville. Bien que la vie en ville soit pratique, le loyer est cher et c'est très bondé. C'est un problème sérieux. »
\end{quote}

\textbf{Version (Chinois $\rightarrow$ Français) :} Traduisez le texte suivant en français courant.
\begin{textebox}[Texte pour version]
虽然 农村 空气 好，环境 美，但是 年轻人 不 想 留下。因为 没有 工作，工资 低。他们 搬去 城市，想 赚 多 钱。但是 城市 房租 贵，生活 累。希望 政府 能 发展 农村 经济。
\end{textebox}
\end{exercice}

\begin{exercice}[Expression écrite : sujet type Bac]
\textbf{Sujet :} \emph{« Pourquoi les jeunes quittent-ils le village pour la ville ? »}

\textbf{Consignes :}
\begin{itemize}
\item Rédigez un texte court en chinois (caractères simplifiés).
\item Longueur : \textbf{8 phrases minimum} (environ 100--120 caractères).
\item Utilisez \textbf{au moins 3 connecteurs logiques} différents parmi : \zh{因为…所以…}，\zh{虽然…但是…}，\zh{不但…而且…}，\zh{首先…其次…}，\zh{因此}。
\item Utilisez au moins \textbf{5 mots du vocabulaire de la leçon} : \zh{搬家，离开，村庄，城市，工作，房租，拥挤，方便，压力，发展}。
\item Soignez l'ordre des mots (Sujet + Temps + Lieu + Verbe + Complément) et l'écriture des caractères.
\item Ajoutez le pinyin avec tons sous votre texte ou à la suite.
\end{itemize}
\end{exercice}

\newpage

\coursec{Corrigés des exercices}

\begin{corrige}[Exercice 1 — QCM : vocabulaire de l'exode rural]
\begin{enumerate}
\item \textbf{B.} \zh{搬家} (bān jiā) : « déménager ». \zh{离开} (líkāi) signifie « quitter » et \zh{返回} (fǎnhuí) « retourner ».
\item \textbf{A.} \zh{房租} (fángzū) : « loyer » (房 = maison + 租 = louer). \zh{工资} (gōngzī) = salaire ; \zh{物价} (wùjià) = niveau des prix.
\item \textbf{B.} \zh{挤} (jǐ) signifie « bondé, serré » ; \zh{拥挤} (yōngjǐ) = « bondé, encombré » (foule, circulation).
\item \textbf{B.} \zh{虽然…但是…} (suīrán… dànshì…) exprime la concession « bien que… mais… ». \zh{因为…所以…} = cause/conséquence ; \zh{不但…而且…} = « non seulement… mais aussi… ».
\end{enumerate}
\end{corrige}

\begin{corrige}[Exercice 2 — Vrai ou Faux : caractères et structures]
\begin{enumerate}
\item \textbf{V.} \zh{搬} contient bien la clé de la main \zh{扌} (variante de 手) à gauche : déménager est une action faite avec les mains.
\item \textbf{F.} \zh{村} (cūn, village) et \zh{材} (cái, matériau) ne se prononcent pas pareil : les tons diffèrent (1\textsuperscript{er} ton / 2\textsuperscript{e} ton). Ils partagent seulement la clé du bois \zh{木} et l'élément phonétique \zh{才}.
\item \textbf{V.} \zh{离开} est un verbe qui prend directement son complément de lieu : \zh{离开家乡} (líkāi jiāxiāng), sans préposition. En français on dit « quitter \emph{son} village », en chinois aucun mot-outil n'est nécessaire.
\item \textbf{V.} La structure correcte est \zh{越来越} + adjectif : \zh{越来越贵} (yuè lái yuè guì, « de plus en plus cher »). \zh{越贵越来} est un ordre de mots fautif.
\end{enumerate}
\end{corrige}

\begin{corrige}[Exercice 3 — Association caractères / pinyin / sens]
\begin{center}
\begin{tabularx}{\linewidth}{|X|X|X|}\hline
\rowcolor{popblueL} \textbf{Caractère} & \textbf{Pinyin} & \textbf{Traduction} \\ \hline
搬 & bān & déménager, déplacer \\ \hline
村 & cūn & village \\ \hline
离 & lí & quitter, s'éloigner \\ \hline
挤 & jǐ & bondé, serré \\ \hline
租 & zū & loyer, louer \\ \hline
\end{tabularx}
\end{center}
\textbf{Astuce mémorisation :} \zh{搬} et \zh{挤} portent la clé de la main \zh{扌} (actions) ; \zh{村} porte la clé du bois \zh{木} ; \zh{租} porte la clé du grain \zh{禾} (autrefois le loyer se payait en récolte).
\end{corrige}

\begin{corrige}[Exercice 4 — Pinyin et tons]
\begin{enumerate}
\item \zh{搬家} : \emph{bān jiā} (1\textsuperscript{er} ton + 1\textsuperscript{er} ton)
\item \zh{农村} : \emph{nóng cūn} (2\textsuperscript{e} ton + 1\textsuperscript{er} ton)
\item \zh{离开} : \emph{lí kāi} (2\textsuperscript{e} ton + 1\textsuperscript{er} ton)
\item \zh{拥挤} : \emph{yōng jǐ} (1\textsuperscript{er} ton + 3\textsuperscript{e} ton)
\item \zh{房租} : \emph{fáng zū} (2\textsuperscript{e} ton + 1\textsuperscript{er} ton)
\end{enumerate}
\textbf{Point de vigilance :} aucun de ces mots ne contient deux 3\textsuperscript{es} tons consécutifs, donc aucune mutation de ton n'est nécessaire ici. Ne confondez pas \emph{lí} (离, 2\textsuperscript{e} ton, « quitter ») avec \emph{lǐ} (里, 3\textsuperscript{e} ton, « intérieur ») ni avec \emph{lì} (力, 4\textsuperscript{e} ton, « force »).
\end{corrige}

\begin{corrige}[Exercice 5 — Thème : phrases à traduire en chinois]
\begin{enumerate}
\item \zh{因为村里没有工作，所以他去城市。}\\
\emph{Yīnwèi cūn lǐ méiyǒu gōngzuò, suǒyǐ tā qù chéngshì.}\\
Structure de cause : \zh{因为} + cause, \zh{所以} + conséquence.
\item \zh{虽然城市很方便，但是房租太贵了。}\\
\emph{Suīrán chéngshì hěn fāngbiàn, dànshì fángzū tài guì le.}\\
Concession : \zh{虽然…但是…} ; \zh{太…了} renforce l'adjectif.
\item \zh{越来越多的年轻人离开农村。}\\
\emph{Yuè lái yuè duō de niánqīngrén líkāi nóngcūn.}\\
\zh{越来越多} + \zh{的} + nom ; \zh{离开} prend directement son complément.
\item \zh{我父母不想搬家，他们喜欢村庄。}\\
\emph{Wǒ fùmǔ bù xiǎng bān jiā, tāmen xǐhuan cūnzhuāng.}\\
Négation de la volonté : \zh{不想} + verbe.
\item \zh{路上很拥挤，车太多了。}\\
\emph{Lù shang hěn yōngjǐ, chē tài duō le.}
\end{enumerate}
\textbf{Points de vigilance :} ne pas mettre de préposition après \zh{离开} ; penser au \zh{的} après \zh{越来越多} ; le complément de lieu se place avant le verbe (\zh{去城市}).
\end{corrige}

\begin{corrige}[Exercice 6 — Version : phrases à traduire en français]
\begin{enumerate}
\item \zh{村里没有工厂，所以年轻人都搬去城市了。} : « Il n'y a pas d'usine au village, donc les jeunes sont tous partis s'installer en ville. » (\zh{都} = tous ; \zh{了} marque l'accompli.)
\item \zh{虽然城市工作机会多，但是生活压力很大。} : « Bien qu'il y ait beaucoup d'opportunités de travail en ville, la pression de la vie est très forte. »
\item \zh{我哥哥去年离开家乡，在深圳打工。} : « Mon grand frère a quitté son village natal l'année dernière et travaille comme ouvrier migrant à Shenzhen. » (\zh{打工} = travailler comme journalier/migrant.)
\item \zh{城市里房租越来越贵，住房很拥挤。} : « En ville, le loyer devient de plus en plus cher et les logements sont très exigus. »
\item \zh{因为交通不方便，所以很少有人回村里。} : « Comme les transports ne sont pas pratiques, peu de gens retournent au village. » (\zh{很少} + nom = « très peu de… ».)
\end{enumerate}
\end{corrige}

\begin{corrige}[Exercice 7 — Caractères complexes : écriture et radicaux]
\begin{center}
\begin{tabularx}{\linewidth}{|X|X|X|X|X|}\hline
\rowcolor{popblueL} Caractère & Traits & Radical & Sens du radical & Ordre des traits \\ \hline
搬 & 13 & 扌 & main (action) & d'abord 扌 à gauche (3 traits), puis 般 à droite, de haut en bas \\ \hline
离 & 10 & 亠 & couvercle / tête & de haut en bas : 亠, puis la partie centrale, puis le bas \\ \hline
挤 & 9 & 扌 & main (presser) & 扌 à gauche (3 traits), puis 齐 à droite de haut en bas \\ \hline
租 & 10 & 禾 & grain, céréale & 禾 à gauche (5 traits), puis 且 à droite \\ \hline
\end{tabularx}
\end{center}
\textbf{Règles générales rappelées :} gauche avant droite, haut avant bas, horizontal avant vertical, l'extérieur avant l'intérieur. Le radical est toujours la clé sémantique : ici \zh{扌} (main) pour \zh{搬} et \zh{挤}, \zh{禾} (grain) pour \zh{租}.
\end{corrige}

\begin{corrige}[Exercice 8 — Discrimination de caractères proches]
\begin{enumerate}
\item \zh{我家在山里的一个小村里。} (\zh{村} cūn = village ; \zh{材} cái = matériau.)
\item \zh{请帮我把书搬到楼上。} (\zh{搬} bān = déplacer ; \zh{般} bān n'est qu'un élément phonétique, jamais employé seul ici.)
\item \zh{地铁里人太多了，非常挤。} (\zh{挤} jǐ = bondé ; \zh{济} jì apparaît dans \zh{经济}, économie.)
\item \zh{这材料很好，很结实。} (\zh{材料} cáiliào = matériau ; \zh{结实} jiēshi = solide, robuste.)
\item \zh{他搬家了，新地址我不知道。} (\zh{搬家} = déménager.)
\end{enumerate}
\textbf{Point de vigilance :} \zh{村}/\zh{材} et \zh{搬}/\zh{般} partagent le même élément phonétique mais pas le même radical : c'est la clé (木, 扌) qui porte le sens.
\end{corrige}

\begin{corrige}[Exercice 9 — Connecteurs et structure : remise en ordre]
Ordre logique : 5 $\rightarrow$ 1 $\rightarrow$ 2 $\rightarrow$ 3 $\rightarrow$ 4 $\rightarrow$ 6.

\zh{现在，越来越多的年轻人离开家乡。因为村里没有工作，所以他们搬去城市。虽然城市生活方便，但是房租很贵。这是一个严重的问题。}

\emph{Xiànzài, yuè lái yuè duō de niánqīngrén líkāi jiāxiāng. Yīnwèi cūn lǐ méiyǒu gōngzuò, suǒyǐ tāmen bān qù chéngshì. Suīrán chéngshì shēnghuó fāngbiàn, dànshì fángzū hěn guì. Zhè shì yí ge yánzhòng de wèntí.}

Traduction : « Aujourd'hui, de plus en plus de jeunes quittent leur village natal. Parce qu'il n'y a pas de travail au village, ils s'installent en ville. Bien que la vie en ville soit pratique, le loyer est très cher. C'est un problème sérieux. »

\textbf{Justification :} on ouvre par le constat général (5), puis la chaîne cause-conséquence \zh{因为…所以…} (1-2), puis la concession \zh{虽然…但是…} (3-4), et on conclut par le jugement (6).
\end{corrige}

\begin{corrige}[Exercice 10 — Texte à trous]
\zh{现在，\textbf{越来越多}青年人\textbf{离开}家乡\textbf{去}城市打工。\textbf{因为}村里没有好\textbf{工作}，\textbf{所以}他们必须\textbf{搬家}。\textbf{虽然}城市机会多，\textbf{但是}城市很\textbf{拥挤}，\textbf{房租}也很贵。}

\emph{Xiànzài, yuè lái yuè duō qīngniánrén líkāi jiāxiāng qù chéngshì dǎgōng. Yīnwèi cūn lǐ méiyǒu hǎo gōngzuò, suǒyǐ tāmen bìxū bānjiā. Suīrán chéngshì jīhuì duō, dànshì chéngshì hěn yōngjǐ, fángzū yě hěn guì.}

Traduction : « Aujourd'hui, de plus en plus de jeunes quittent leur village pour aller travailler en ville. Comme il n'y a pas de bon travail au village, ils doivent déménager. Bien que la ville offre beaucoup d'opportunités, elle est très bondée et le loyer est aussi très cher. »

\textbf{Note :} le troisième trou attend \zh{去} (aller) — mot de direction déjà connu ; les autres mots viennent de la liste. Remarquez la structure \zh{离开} + lieu + \zh{去} + lieu + verbe : « quitter X pour aller faire Y ».
\end{corrige}

\begin{corrige}[Exercice 11 — Compréhension de l'écrit et langue]
\textbf{Questions de compréhension :}
\begin{enumerate}
\item Les deux pays concernés sont le \textbf{Cameroun} (\zh{喀麦隆}) et la \textbf{Chine} (\zh{中国}).
\item Trois raisons : la terre est rare au village (\zh{土地少}), les récoltes sont instables (\zh{收成不稳定}) et il n'y a pas d'usines (\zh{没有工厂}), donc il est difficile de trouver du travail.
\item Avantages : la vie est pratique, il y a des hôpitaux, des magasins et de bonnes écoles. Inconvénients : le loyer est très cher et la circulation est très encombrée.
\item Ils travaillent souvent comme serveurs, chauffeurs ou ouvriers sur les chantiers. Non, leur situation n'est pas facile : leurs salaires sont bas et la pression est forte.
\item Les experts proposent de développer la campagne : construire de bonnes routes, un bon réseau Internet et des usines, pour que les jeunes ne veuillent plus partir.
\end{enumerate}
\textbf{Questions de langue :}
\begin{enumerate}
\item Structure de concession \zh{虽然…但是…} : \zh{虽然城市生活方便，有医院、商场和好学校，但是房租很贵，交通也很拥挤。}
\item « Récolte / rendement agricole » = \zh{收成} (shōucheng).
\item \zh{越来越多的年轻人离开村庄。} (\emph{Yuè lái yuè duō de niánqīngrén líkāi cūnzhuāng.}) — ne pas oublier le \zh{的}.
\item \emph{zhuānjiā} = \zh{专家} (expert) ; \emph{fāzhǎn} = \zh{发展} (développer).
\end{enumerate}
\end{corrige}

\begin{corrige}[Exercice 12 — Thème et Version type Baccalauréat]
\textbf{Thème — proposition de traduction :}

\zh{因为村里没有足够的工作，所以越来越多的年轻人搬去城市。虽然城市生活很方便，但是房租很贵，也很拥挤。这是一个严重的问题。}

\emph{Yīnwèi cūn lǐ méiyǒu zúgòu de gōngzuò, suǒyǐ yuè lái yuè duō de niánqīngrén bān qù chéngshì. Suīrán chéngshì shēnghuó hěn fāngbiàn, dànshì fángzū hěn guì, yě hěn yōngjǐ. Zhè shì yí ge yánzhòng de wèntí.}

\textbf{Version — proposition de traduction :}

« Bien que l'air soit pur et le paysage beau à la campagne, les jeunes ne veulent pas y rester, car il n'y a pas de travail et les salaires sont bas. Ils s'installent en ville, espérant gagner plus d'argent. Mais en ville le loyer est cher et la vie est fatigante. Espérons que le gouvernement puisse développer l'économie rurale. »

\textbf{Barème indicatif (sur 10) :} exactitude du vocabulaire (3 pts) ; structures \zh{因为…所以…} et \zh{虽然…但是…} correctes (3 pts) ; ordre des mots et mots-outils \zh{的}/\zh{了} (2 pts) ; exactitude des caractères (2 pts).

\textbf{Points de vigilance :} \zh{搬去城市} et non \zh{搬在城市} ; \zh{越来越} + adjectif ; \zh{的} obligatoire après \zh{越来越多} ; en version, \zh{希望} introduit un souhait (« espérons que… »).
\end{corrige}

\begin{corrige}[Exercice 13 — Expression écrite : sujet type Bac]
\textbf{Proposition de rédaction modèle :}

\zh{现在，越来越多的年轻人离开村庄去城市。为什么呢？首先，因为村里没有工厂，找工作很难，所以他们必须搬家。其次，农村的工资很低，年轻人想赚更多的钱。虽然城市生活很方便，有医院和好学校，但是房租很贵，交通也很拥挤。很多年轻人不但工资不高，而且压力很大。因此，我认为政府应该发展农村，建工厂和好路。如果家乡变好了，年轻人就不想走了。}

\emph{Xiànzài, yuè lái yuè duō de niánqīngrén líkāi cūnzhuāng qù chéngshì. Wèishénme ne? Shǒuxiān, yīnwèi cūn lǐ méiyǒu gōngchǎng, zhǎo gōngzuò hěn nán, suǒyǐ tāmen bìxū bānjiā. Qícì, nóngcūn de gōngzī hěn dī, niánqīngrén xiǎng zhuàn gèng duō de qián. Suīrán chéngshì shēnghuó hěn fāngbiàn, yǒu yīyuàn hé hǎo xuéxiào, dànshì fángzū hěn guì, jiāotōng yě hěn yōngjǐ. Hěn duō niánqīngrén búdàn gōngzī bù gāo, érqiě yālì hěn dà. Yīncǐ, wǒ rènwéi zhèngfǔ yīnggāi fāzhǎn nóngcūn, jiàn gōngchǎng hé hǎo lù. Rúguǒ jiāxiāng biàn hǎo le, niánqīngrén jiù bù xiǎng zǒu le.}

Traduction : « Aujourd'hui, de plus en plus de jeunes quittent le village pour la ville. Pourquoi ? D'abord, comme il n'y a pas d'usines au village, trouver du travail est difficile, donc ils doivent déménager. Ensuite, les salaires à la campagne sont bas et les jeunes veulent gagner plus d'argent. Bien que la vie en ville soit pratique, avec des hôpitaux et de bonnes écoles, le loyer est cher et la circulation encombrée. Beaucoup de jeunes ont non seulement des salaires bas, mais aussi une forte pression. C'est pourquoi je pense que le gouvernement doit développer la campagne, construire des usines et de bonnes routes. Si le village natal s'améliore, les jeunes ne voudront plus partir. »

\textbf{Barème indicatif (sur 20) :} respect des consignes (longueur, 3 connecteurs, 5 mots du thème) : 5 pts ; correction grammaticale (ordre des mots, \zh{的}/\zh{了}, structures) : 6 pts ; richesse et exactitude du vocabulaire : 4 pts ; exactitude des caractères : 3 pts ; cohérence et organisation du texte : 2 pts.

\textbf{Points de vigilance :}
\begin{itemize}
\item Connecteurs employés ici : \zh{首先…其次…}, \zh{因为…所以…}, \zh{虽然…但是…}, \zh{不但…而且…}, \zh{因此} (5 connecteurs, minimum exigé : 3).
\item Vocabulaire du thème employé : \zh{离开}, \zh{村庄}, \zh{搬家}, \zh{房租}, \zh{拥挤}, \zh{方便}, \zh{压力}, \zh{发展} (8 mots, minimum exigé : 5).
\item Attention à \zh{不但} qui se prononce \emph{búdàn} devant un 4\textsuperscript{e} ton (mutation de 不).
\item Ne pas oublier le \zh{的} après \zh{越来越多} et dans \zh{更多的钱}.
\item La structure \zh{如果…就…} (rúguǒ… jiù…) exprime la condition : « si… alors… » ; le \zh{了} final marque le changement d'état.
\end{itemize}
\end{corrige}

\newpage

\coursec{Fiche bilan}

\begin{popfigure}
\begin{tikzpicture}[
  mindmap,
  grow cyclic,
  every node/.style={concept, execute at begin node=\hskip0pt},
  concept color=popblue,
  level 1/.style={level distance=4.5cm, sibling angle=60, font=\large\bfseries},
  level 2/.style={level distance=3cm, sibling angle=45, font=\normalsize},
  branch1/.style={concept color=poporange, font=\normalsize\bfseries},
  branch2/.style={concept color=popgreen, font=\normalsize\bfseries},
  branch3/.style={concept color=poppurple, font=\normalsize\bfseries},
  branch4/.style={concept color=poppink, font=\normalsize\bfseries},
  branch5/.style={concept color=popgold, font=\normalsize\bfseries},
  branch6/.style={concept color=popdark, font=\normalsize\bfseries},
  text=white
]
\node[concept, text width=3.2cm, align=center, font=\Large\bfseries] {Leçon 26\\Exode rural\\表达写作}
  child[branch1] {node[concept, text width=2.5cm, align=center] {Vocabulaire\\clé\\重点词汇}}
  child[branch2] {node[concept, text width=2.5cm, align=center] {Caractères\\complexes\\复杂汉字}}
  child[branch3] {node[concept, text width=2.5cm, align=center] {Connecteurs\\logiques\\连接词}}
  child[branch4] {node[concept, text width=2.5cm, align=center] {Structures\\grammaticales\\语法结构}}
  child[branch5] {node[concept, text width=2.5cm, align=center] {Thème /\ Version\\翻译练习}}
  child[branch6] {node[concept, text width=2.5cm, align=center] {Contexte socio-\\culturel\\社会文化}};
\end{tikzpicture}
\legende{Carte mentale de la Leçon 26 : Expression écrite sur l'exode rural}
\end{popfigure}

\begin{bilan}

\end{bilan}

\courssub{Lexique clé du thème (HSK 3--4)}

\begin{center}
\begin{tabularx}{\linewidth}{|X|X|X|X|}
\hline
\rowcolor{popblueL} \textbf{Caractères} & \textbf{Pinyin} & \textbf{Nature} & \textbf{Traduction} \\ \hline
农村 & nóngcūn & n. & campagne, zone rurale \\ \hline
城市 & chéngshì & n. & ville, zone urbaine \\ \hline
农民 & nóngmín & n. & paysan, agriculteur \\ \hline
打工 & dǎgōng & v. & faire des petits boulots, travailler comme migrant \\ \hline
外出 & wàichū & v. & sortir, partir (travailler ailleurs) \\ \hline
务工 & wùgōng & v. & aller travailler (terme formel) \\ \hline
留守 & liúshǒu & v./adj. & rester sur place (enfants, personnes âgées) \\ \hline
留守儿童 & liúshǒu értóng & n. & enfants laissés derrière \\ \hline
空巢 & kōngcháo & n./adj. & nid vide (personnes âgées seules) \\ \hline
收入 & shōurù & n. & revenu, gains \\ \hline
机会 & jīhuì & n. & opportunité, chance \\ \hline
发展 & fāzhǎn & v./n. & développer, développement \\ \hline
适应 & shìyìng & v. & s'adapter \\ \hline
困难 & kùnnán & n./adj. & difficulté, difficile \\ \hline
改变 & gǎibiàn & v. & changer, transformer \\ \hline
希望 & xīwàng & n./v. & espoir, espérer \\ \hline
\end{tabularx}
\end{center}

\vspace{0.5em}
\courssub{Règles de grammaire et structures à connaître par cœur}

\begin{center}
\begin{tabularx}{\linewidth}{|X|X|X|}
\hline
\rowcolor{popblueL} \textbf{Structure / Règle} & \textbf{Exemple (Caractères + Pinyin)} & \textbf{Traduction / Note} \\ \hline
\emph{Parce que... donc...} (Cause $\rightarrow$ Conséquence) & \zh{因为农村收入低，所以他外出打工。} (Yīnwèi nóngcūn shōurù dī, suǒyǐ tā wàichū dǎgōng.) & « Parce que les revenus ruraux sont bas, il part travailler dehors. » \textbf{Ne pas omettre 所以}. \\ \hline
\emph{De plus en plus} (Évolution) & \zh{留守儿童越来越多了。} (Liúshǒu értóng yuè lái yuè duō le.) & « Les enfants laissés derrière sont de plus en plus nombreux. » \textbf{+ 了} obligatoire pour le changement d'état. \\ \hline
\emph{Comparatif 比} (A plus que B) & \zh{城市的机会比农村多。} (Chéngshì de jīhuì bǐ nóngcūn duō.) & « Les opportunités en ville sont plus nombreuses qu'à la campagne. » Pas de \zh{更} après l'adj. \\ \hline
\emph{Changement d'état / Aspect accompli \zh{了}} & \zh{他已经适应城市生活了。} (Tā yǐjīng shìyìng chéngshì shēnghuó le.) & « Il s'est déjà adapté à la vie urbaine. » Marque la nouvelle situation. \\ \hline
\emph{Connecteurs : 但是 / 而且 / 因此 / 一方面...另一方面...} & \zh{他想回家，但是收入不够。} (Tā xiǎng huí jiā, dànshì shōurù bù gòu.) & « Il veut rentrer, mais les revenus ne suffisent pas. » \textbf{Ordre : Suj + Adv + V}. \\ \hline
\emph{Structure \zh{把} (Disposition objet)} & \zh{农民把土地承包给别人。} (Nóngmín bǎ tǔdì chéngbāo gěi biérén.) & « Le paysan confie la terre à d'autres. » Objet connu, action aboutie. \\ \hline
\end{tabularx}
\end{center}

\vspace{0.5em}
\courssub{Méthodes express}

\begin{methode}[Écrire un caractère complexe (ex: \zh{留}, \zh{务}, \zh{适})]
\begin{enumerate}
\item \textbf{Identifier le radical (clé) :} \zh{留} ($\rightarrow$ 刀/刂), \zh{务} ($\rightarrow$ 力), \zh{适} ($\rightarrow$ 辶).
\item \textbf{Décomposer :} Haut $\rightarrow$ Bas, Gauche $\rightarrow$ Droite, Extérieur $\rightarrow$ Intérieur $\rightarrow$ Fermeture.
\item \textbf{Respecter l'ordre des traits :} Horizontales avant verticales, traits centraux avant latéraux, point final en dernier.
\item \textbf{Vérifier la structure :} \zh{留} (haut-bas), \zh{务} (haut-bas), \zh{适} (gauche-droite / semi-entourant).
\item \textbf{S'entraîner :} 5 fois en respectant le carré imaginaire, puis de mémoire.
\end{enumerate}
\end{methode}

\begin{methode}[Traduction Thème (Fr $\rightarrow$ Zh) / Version (Zh $\rightarrow$ Fr)]
\begin{enumerate}
\item \textbf{Analyser} la phrase source : Sujet, Verbe, Compléments (Lieu, Temps, Manière), Particules (\zh{了, 的, 得, 地}).
\item \textbf{Ordonner} selon la syntaxe chinoise : Sujet + Temps + Lieu + Manière + Verbe + Complément / Particule.
\item \textbf{Choisir} le vocabulaire précis (ex: \zh{外出务工} vs \zh{打工}, \zh{留守儿童} vs \zh{孤儿}).
\item \textbf{Vérifier} : Tons (pinyin), Caractères (simplifiés), Ponctuation chinoise (，。), \zh{了} (aspect/changement), \zh{的/得/地}.
\end{enumerate}
\end{methode}

\begin{methode}[Rédaction d'un paragraphe cohérent (80--100 caractères)]
\begin{enumerate}
\item \textbf{Plan :} Phrase d'accroche (Situation) $\rightarrow$ Causes (\zh{因为...所以...}) $\rightarrow$ Conséquences (\zh{导致...}, \zh{越来越...}) $\rightarrow$ Solutions / Avenir (\zh{应该...}, \zh{希望...}).
\item \textbf{Relier} avec connecteurs : \zh{首先}, \zh{其次}, \zh{然而}, \zh{因此}, \zh{总之}.
\item \textbf{Écrire} des phrases courtes (SVO), éviter les phrases trop longues sans virgule.
\item \textbf{Relire} : Ordre mots (Temps/Lieu avant Verbe), \zh{了} final, accords \zh{的/得/地}.
\end{enumerate}
\end{methode}



\begin{aretenir}[Checklist avant le Bac -- Leçon 26]
$\square$ Je connais par cœur le vocabulaire de l'exode rural (Chinois $\leftrightarrow$ Français, Pinyin + Tons). \\
$\square$ J'écris correctement les 15--20 caractères complexes du thème (ordre traits, radicaux, structure). \\
$\square$ Je distingue \zh{留守} (rester) de \zh{外出} (partir) et \zh{务工} (travail formel) de \zh{打工} (job). \\
$\square$ Je construis une phrase cause/effet avec \zh{因为...所以...} sans inverser l'ordre. \\
$\square$ J'utilise \zh{越来越 + Adj + 了} pour exprimer une évolution (ex: \zh{越来越严重了}). \\
$\square$ Je forme le comparatif \zh{A 比 B Adj} sans ajouter \zh{更} ni \zh{很}. \\
$\square$ Je place correctement les compléments de Lieu/Temps \textbf{avant} le verbe. \\
$\square$ J'utilise \zh{了} en fin de phrase pour marquer le changement d'état (nouvelle situation). \\
$\square$ Je choisis le bon connecteur : \zh{但是} (contraste), \zh{而且} (addition), \zh{因此} (conséquence). \\
$\square$ Je respecte la structure « Situation -- Causes -- Conséquences -- Avenir » en rédaction. \\
$\square$ En thème/version, je vérifie \zh{的/得/地} et la ponctuation chinoise pleine chasse. \\
$\square$ Je cite un exemple concret Cameroun (ex: \zh{雅温得}, \zh{杜阿拉}, \zh{恩多莱}) ou Chine (\zh{春运}, \zh{农民工}).
\end{aretenir}

\findecours

\end{document}
